\chapter{Análisis y diseño}\label{cap:analisis}

\section{Metodología}\label{sec:metodologia}

Extreme Programming (XP) es una metodología ágil de desarrollo de software propuesta por Beck y Andres \cite{beck2004}, orientada a equipos pequeños que trabajan con requisitos técnicos complejos o inciertos. Sus prácticas centrales incluyen programación en pares, desarrollo guiado por pruebas (TDD), spikes arquitectónicos para resolver incertidumbre técnica, e integración continua con propiedad colectiva del código.

Se eligió XP por tres razones vinculadas a las características del proyecto: primero, el equipo está compuesto por dos integrantes, lo que favorece la programación en pares como práctica central en lugar de ceremonias de coordinación más pesadas; segundo, la integración de mecanismos criptográficos (firma digital, sellado de tiempo, cifrado autenticado) y la vinculación con servicios externos como la Autoridad de Sellado de Tiempo introducen alta incertidumbre técnica, que XP aborda mediante \emph{Spikes} (pruebas de concepto) y ciclos cortos de retroalimentación; tercero, el Desarrollo Guiado por Pruebas (TDD) resulta especialmente adecuado para validar la corrección de los servicios criptográficos, donde un error silencioso puede invalidar toda la evidencia generada.

Para abordar la complejidad del sistema, el proyecto estará segmentado en dos fases de iteraciones con objetivos específicos. El bloque inicial se enfocará en la planificación y la resolución de riesgos técnicos mediante \emph{Spikes}, dejando el desarrollo iterativo continuo para el segundo bloque.

\begin{figure}[H]
  \centering
  \includegraphics[width=0.7\linewidth]{image4.png}
  \caption{Ciclo de la metodología XP (Extreme Programming) aplicada al proyecto.}
  \label{fig:xp-metodologia}
\end{figure}

\textbf{Plan de Iteraciones (10 semanas)}. El proyecto se estructura en dos bloques principales de 5 semanas cada uno. El primer bloque corresponde a la fase de planeación, investigación y diseño, enfocado en mitigar riesgos y modelar la solución. El segundo bloque abarca el ciclo continuo de codificación, integración y pruebas.

\textbf{Fase I: Exploración, Diseño y Spikes Arquitectónicos}. Durante estas 5 semanas, las iteraciones se centran en el levantamiento ágil, el análisis normativo y el diseño técnico del sistema, preparando el terreno para la implementación.

\begin{itemize}
  \item \textbf{Iteración 1 (Semana 1) - Requerimientos y Fundamentos Criptográficos:} Investigación de la normativa mexicana referente al manejo de firmas digitales y comprobación de tiempo, y recolección de normativas sobre la Ley de Protección de Datos Personales. Investigación teórica de criptografía (Hash, RSA, Firma Digital) y levantamiento ágil de requerimientos funcionales y no funcionales.

  \item \textbf{Iteración 2 (Semana 2) - Análisis de Entorno y Sistemas Similares:} Investigación del Código de Procedimientos Penales Mexicanos aplicable al contexto digital. Análisis de sistemas similares en el mercado para la detección de brechas de seguridad y operativas.

  \item \textbf{Iteración 3 (Semana 3) - Análisis de Riesgos y Modelado Ágil:} Elaboración del análisis y matriz de riesgos. Traducción de los requerimientos mediante la redacción de Historias de Usuario priorizadas y diseño de los Casos de Uso del sistema.

  \item \textbf{Iteración 4 (Semana 4) - Diseño Arquitectónico y Modelo de Datos:} Diseño detallado de la arquitectura del sistema (Backend, Frontend y flujos de seguridad). Modelado de la base de datos híbrida, integrando bases de datos SQL para relaciones procesales y bases de datos Clave-Valor.

  \item \textbf{Iteración 5 (Semana 5) - Estrategia de Calidad y Cierre Documental:} Creación de la matriz de pruebas (funcionales, de integración y criptográficas). Cierre y validación documental del diseño técnico para la conclusión de la primera fase.
\end{itemize}

\textbf{Fase II: Construcción, Pruebas y Refactorización}. Una vez validado el diseño, el desarrollo entra en un ciclo estricto de codificación iterativa mediante XP, utilizando programación en pares, Desarrollo Guiado por Pruebas (TDD) y entregas continuas. El orden de las iteraciones antepone el núcleo criptográfico a la construcción de vistas: el análisis de riesgos (sección~\ref{sec:riesgos}) prescribe resolver primero los riesgos criptográficos graves y recomienda no avanzar hacia el desarrollo de vistas web hasta haber comprobado técnicamente la validez del sellado de tiempo.

\begin{itemize}
  \item \textbf{Iteración 6 (Semana 6) - Andamiaje del Workspace y Núcleo Criptográfico Base:} Configuración del workspace de Cargo con sus cinco crates y del pipeline de integración continua (formato, análisis estático, pruebas y política de dependencias). Implementación del servicio de hash SHA-256 y del cifrado autenticado AES-256-GCM con envelope encryption, proporcionando integridad y confidencialidad en reposo. Primer contacto con el sellado de tiempo mediante el Spike de consumo y parseo del token RFC 3161.

  \item \textbf{Iteración 7 (Semana 7) - PKI Interna y Firma Digital:} Implementación de la infraestructura de llave pública interna: CA raíz, emisión de certificados X.509 y lista de certificados revocados (CRL). Integración de la firma digital RSA-3072 con SHA-256 sobre los certificados emitidos. Inicio del flujo de sellado de tiempo conforme a RFC 3161.

  \item \textbf{Iteración 8 (Semana 8) - Sellado de Tiempo, Autenticación y Auditoría:} Conclusión del flujo de sellado de tiempo RFC 3161. Implementación de la autenticación (hashing de contraseñas con Argon2id y segundo factor TOTP) y de la cadena de auditoría con encadenamiento de hashes. Verificación integral de la evidencia y demostración por línea de comandos, hito de criptografía demostrada de extremo a extremo.

  \item \textbf{Iteración 9 (Semana 9) - Persistencia, Control de Acceso y Gestión Procesal:} Implementación de la persistencia híbrida PostgreSQL/Redis y del control de acceso basado en roles (RBAC) con JWT. Desarrollo de la gestión procesal y documental —módulos CRUD para expedientes, casos, contratos y documentos— integrando los servicios criptográficos ya probados. Desarrollo de una interfaz web mínima.

  \item \textbf{Iteración 10 (Semana 10) - Pruebas Finales y Cierre:} Ejecución de la matriz de pruebas diseñada en la Fase I. Pruebas de usabilidad con usuarios representativos. Corrección de errores, refactorización final de código y elaboración de los manuales técnicos y de usuario.
\end{itemize}

\section{Historias de Usuario}\label{sec:historias-usuario}

Las historias de usuario constituyen el instrumento de captura de requerimientos empleado por la metodología XP. Cada historia se redacta desde la perspectiva del actor que interactúa con el sistema, en lenguaje del dominio jurídico, estableciendo el rol, la necesidad y el beneficio esperado. Las historias se organizan en siete secciones funcionales que cubren la totalidad de los flujos del prototipo, desde la gestión de identidad hasta la operación y monitoreo del despacho. La formalización técnica de cada historia se realiza posteriormente en los casos de uso (sección~\ref{sec:casos-uso}), donde se especifican los flujos, precondiciones, postcondiciones y los mecanismos concretos de implementación.

\subsection{Autenticación, Acceso y Gestión de Identidad}\label{subsec:hu-autenticacion}

\begin{itemize}
  \item \textbf{HU-01 — Registro de despacho.} Como socio fundador, quiero registrar mi despacho en el sistema proporcionando los datos de la firma, para disponer de un espacio propio y aislado donde gestionar los expedientes y al equipo de trabajo.

  \item \textbf{HU-02 — Selección de plan de servicio.} Como socio fundador, quiero elegir un plan de servicio (Starter o Professional) al momento del registro, para que el sistema habilite las funcionalidades y los límites de usuarios correspondientes al nivel contratado.

  \item \textbf{HU-03 — Invitación de miembros al equipo.} Como socio del despacho, quiero invitar nuevos colaboradores mediante correo electrónico asignándoles un perfil (Litigante o Paralegal), para que puedan acceder al sistema con los permisos apropiados a su función.

  \item \textbf{HU-04 — Aceptación de invitación.} Como colaborador invitado, quiero aceptar la invitación recibida por correo y completar mi registro, para incorporarme al despacho con el perfil que me fue asignado.

  \item \textbf{HU-05 — Inicio de sesión con credenciales.} Como usuario del sistema, quiero autenticarme con mi correo y contraseña, para acceder a las funciones habilitadas según mi perfil.

  \item \textbf{HU-06 — Inicio de sesión con certificado digital.} Como socio del despacho, quiero iniciar sesión presentando mi certificado digital, para acceder al sistema con un nivel de seguridad superior que vincule mi identidad de forma verificable con la sesión.

  \item \textbf{HU-07 — Verificación en dos pasos.} Como cliente del despacho, quiero confirmar mi identidad con un código temporal generado en mi dispositivo después de ingresar mi contraseña, para que mi expediente esté protegido ante el robo o filtración de la contraseña.

  \item \textbf{HU-08 — Recuperación de contraseña.} Como usuario del sistema, quiero solicitar la recuperación de mi contraseña mediante correo electrónico, para restablecer el acceso a mi cuenta sin necesidad de contactar al administrador.

  \item \textbf{HU-09 — Cierre de sesión por inactividad.} Como usuario del sistema, quiero que mi sesión se cierre automáticamente tras un periodo de inactividad, para prevenir accesos no autorizados si abandono mi equipo de trabajo sin cerrar sesión.

  \item \textbf{HU-10 — Restricción de acceso por perfil.} Como socio del despacho, quiero que cada miembro del equipo solo pueda acceder a los módulos y operaciones que corresponden a su perfil, para garantizar que únicamente las personas autorizadas realicen operaciones sensibles.

  \item \textbf{HU-11 — Baja de miembros del equipo.} Como socio del despacho, quiero dar de baja a un colaborador, para revocar de inmediato su acceso al sistema y a los expedientes del despacho.
\end{itemize}

\subsection{Gestión Procesal Penal (Ciclo de Vida del Expediente)}\label{subsec:hu-procesal}

\begin{itemize}
  \item \textbf{HU-12 — Creación de expediente penal.} Como Litigante, quiero registrar un nuevo expediente ingresando el Número Único de Caso, la Carpeta Judicial, el tipo de delito y los datos generales del asunto, para que quede disponible en el sistema con estado inicial “Investigación”.

  \item \textbf{HU-13 — Registro de participantes del caso.} Como Litigante, quiero vincular a un expediente a las personas y autoridades que participan en el proceso (imputado, víctima, defensor, juez, Ministerio Público, entre otros) con sus datos identificadores y su rol dentro del proceso, para conformar el directorio de participantes del caso.

  \item \textbf{HU-14 — Avance de etapa procesal.} Como Litigante, quiero avanzar un expediente a la siguiente etapa procesal cargando el documento que lo habilita (escrito de acusación para pasar a Intermedia, auto de apertura para pasar a Juicio Oral), para que el sistema registre el cambio y habilite las funciones correspondientes a la nueva etapa.

  \item \textbf{HU-15 — Bloqueo de avance sin documento habilitante.} Como Litigante, quiero que el sistema no permita avanzar de etapa si no se ha cargado el documento procesal requerido, para evitar inconsistencias en la secuencia del expediente que comprometan su validez.

  \item \textbf{HU-16 — Historial de transiciones del expediente.} Como Litigante, quiero consultar la línea de tiempo completa de las transiciones de etapa de un expediente, con la fecha, el responsable y el documento asociado a cada cambio, para verificar que la secuencia procesal esté íntegra.

  \item \textbf{HU-17 — Actualización de datos del expediente.} Como Litigante, quiero actualizar los datos generales de un expediente (delitos, identificadores complementarios, información general), para mantener la información vigente sin alterar el historial de transiciones.

  \item \textbf{HU-18 — Vista de expedientes asignados.} Como Paralegal, quiero visualizar únicamente los expedientes que me han sido asignados, para concentrarme en los casos relevantes a mi trabajo sin exposición a información de otros asuntos.
\end{itemize}

\subsection{Calendario Jurídico, Plazos y Alertas}\label{subsec:hu-calendario}

\begin{itemize}
  \item \textbf{HU-19 — Programación de audiencia.} Como Litigante, quiero agendar una audiencia seleccionando su tipo (inicial, intermedia, juicio oral, individualización de sanciones), fecha, hora, participantes y modalidad, para que quede vinculada al expediente correspondiente.

  \item \textbf{HU-20 — Activación automática de plazos tras audiencia.} Como Litigante, quiero que al registrar el resultado de una audiencia (por ejemplo, vinculación a proceso), el sistema calcule y active automáticamente los plazos derivados (como el de investigación complementaria), para no depender del cálculo manual de los términos.

  \item \textbf{HU-21 — Cómputo de plazos con días inhábiles.} Como Litigante, quiero que el sistema calcule las fechas de vencimiento descontando correctamente sábados, domingos y días inhábiles del calendario judicial de la entidad federativa, para obtener fechas precisas conforme a las reglas de cómputo del CNPP.

  \item \textbf{HU-22 — Configuración del calendario judicial.} Como socio del despacho, quiero configurar los días inhábiles por entidad federativa y fuero, para que los cómputos de plazos reflejen los periodos de receso de cada jurisdicción donde litiga el despacho.

  \item \textbf{HU-23 — Alertas de vencimiento próximo.} Como Litigante, quiero recibir notificaciones automáticas por correo y en el sistema cuando un plazo esté próximo a vencer, para tomar acciones oportunas y evitar consecuencias procesales irreversibles.

  \item \textbf{HU-24 — Calendario unificado del despacho.} Como Litigante, quiero visualizar en un calendario interactivo todas las audiencias programadas y los plazos límite de mis casos, con colores diferenciados por tipo de evento y acceso directo al expediente, para gestionar mi agenda procesal desde una sola vista.
\end{itemize}

\subsection{Gestión Documental y Cadena de Custodia}\label{subsec:hu-documental}

\begin{itemize}
  \item \textbf{HU-25 — Carga de documento al expediente.} Como Litigante, quiero cargar un documento procesal asociándolo a un expediente, para que quede almacenado de forma protegida y disponible para consulta por las personas autorizadas.

  \item \textbf{HU-26 — Clasificación de documentos.} Como Litigante, quiero asignar tipo documental, clasificación y etiquetas a cada documento cargado, para facilitar su organización, búsqueda y recuperación dentro del expediente.

  \item \textbf{HU-27 — Consulta de documentos según asignación.} Como Paralegal, quiero consultar y descargar los documentos de los expedientes que tengo asignados, para apoyar la gestión del caso sin acceder a información de asuntos no relacionados.

  \item \textbf{HU-28 — Registro automático de movimientos sobre documentos.} Como Litigante, quiero que cada operación sobre un documento o evidencia (carga, consulta, descarga, movimiento) quede registrada automáticamente con la fecha, la persona responsable y el tipo de operación, para contar con un historial completo que respalde la cadena de custodia del expediente.

  \item \textbf{HU-29 — Protección del historial de movimientos.} Como Litigante, quiero que ningún usuario, incluido el socio del despacho, pueda modificar ni eliminar los registros del historial de movimientos de documentos, para garantizar que la trazabilidad sea confiable ante cualquier requerimiento judicial.
\end{itemize}

\subsection{Servicios de Protección de la Información}\label{subsec:hu-proteccion}

\begin{itemize}
  \item \textbf{HU-30 — Huella digital del documento al momento de carga.} Como Litigante, quiero que al cargar un documento al sistema, se genere automáticamente una huella digital única del archivo, para contar con una referencia que permita comprobar en el futuro si el documento fue alterado.

  \item \textbf{HU-31 — Alerta de alteración en consulta posterior.} Como Litigante, quiero que cada vez que consulte un documento, el sistema verifique automáticamente que no haya sido modificado desde su carga y me alerte de inmediato si detecta alguna discrepancia, para tener certeza de la inalterabilidad de mis expedientes.

  \item \textbf{HU-32 — Protección del contenido almacenado.} Como Litigante, quiero que los documentos de mis expedientes se almacenen protegidos de forma que solo las partes autorizadas puedan acceder a su contenido, para que una intrusión al sistema no exponga la información de mis clientes.

  \item \textbf{HU-33 — Detección de manipulación en el almacenamiento.} Como Litigante, quiero que si alguien manipula directamente los archivos almacenados en el servidor, el sistema lo detecte y rechace la entrega del documento afectado notificándome la anomalía, para saber que cualquier intento de alteración será evidenciado.
\end{itemize}

\subsection{Firma de Contratos y Constancia de Fecha}\label{subsec:hu-firma}

\begin{itemize}
  \item \textbf{HU-34 — Firma de contrato con identidad verificable.} Como Litigante, quiero firmar un contrato de manera que mi identidad quede vinculada de forma verificable con el contenido del documento, para que ninguna de las partes pueda negar posteriormente haber suscrito el acuerdo.

  \item \textbf{HU-35 — Constancia de fecha y hora de la firma.} Como Litigante, quiero que al firmar un contrato, el sistema obtenga automáticamente una constancia emitida por una autoridad certificadora que acredite la fecha y hora exacta de la firma, para que el documento tenga valor probatorio conforme a la normatividad mexicana.

  \item \textbf{HU-36 — Verificación de firma y constancia de un contrato.} Como usuario con acceso al contrato, quiero poder ejecutar una verificación que me indique si la firma es válida, si el certificado del firmante está vigente y si la constancia de fecha es auténtica, para evaluar la confiabilidad del documento sin depender de un tercero.

  \item \textbf{HU-37 — Detección de alteración en contrato firmado.} Como usuario con acceso al contrato, quiero que si el documento fue modificado después de firmado, la verificación me indique claramente que la firma ya no es válida y me explique la causa, para disponer de evidencia de que el contrato fue alterado.
\end{itemize}

\subsection{Operación y Monitoreo}\label{subsec:hu-operacion}

\begin{itemize}
  \item \textbf{HU-38 — Tablero de control del despacho.} Como socio del despacho, quiero consultar un tablero con indicadores clave (casos activos, plazos por vencer próximamente, contratos pendientes de firma, carga de trabajo por abogado), para tener visibilidad del estado general de la operación.

  \item \textbf{HU-39 — Generación de informes.} Como socio del despacho, quiero generar informes filtrados por periodo, abogado responsable o estado de los casos, exportables en formato descargable, para analizar el desempeño y las métricas operativas de la firma.

  \item \textbf{HU-40 — Registro general de actividad.} Como socio del despacho, quiero consultar un registro completo e inalterable de todas las acciones realizadas en el sistema (quién realizó la acción, qué hizo, sobre qué recurso y cuándo), para disponer de trazabilidad ante auditorías internas o requerimientos judiciales.
\end{itemize}


\section{Casos de Uso}\label{sec:casos-uso}

Los casos de uso formalizan las historias de usuario definidas en la \cref{sec:historias-usuario}, especificando los flujos de interacción entre los actores y el sistema, las condiciones que deben cumplirse antes y después de cada operación, y los escenarios alternos que el prototipo debe manejar. Cada caso de uso agrupa una o más historias de usuario en un flujo funcional cohesivo e incluye la trazabilidad hacia los requerimientos funcionales y no funcionales que lo respaldan. La especificación técnica de los mecanismos criptográficos (algoritmos, estándares y protocolos) se detalla en los requerimientos de la \cref{sec:requerimientos}; los casos de uso se limitan a describir el comportamiento observable del sistema desde la perspectiva del actor.

\subsection{Autenticación, Acceso y Gestión de Identidad}\label{sec:cu-autenticacion}

\subsubsection{CU-01 --- Registro de despacho y selección de plan}

\begin{table}[H]
  \centering
  \caption{Caso de uso CU-01: Registro de despacho y selección de plan.}
  \label{tab:cu-01}
  \begin{tabularx}{\linewidth}{@{}>{\bfseries}l X@{}}
    \toprule
    ID & \textbf{CU-01} \\
    \midrule
    Nombre & Registro de despacho y selección de plan \\
    Actor principal & Socio fundador \\
    Actores secundarios & Sistema de correo transaccional \\
    Precondiciones & El usuario no posee una cuenta activa en el sistema. \\
    Flujo principal & 1. El usuario accede al formulario de registro. 2. Ingresa los datos del despacho (nombre de la firma, domicilio fiscal, datos de contacto). 3. Selecciona un plan de servicio (Starter o Professional). 4. Establece sus credenciales de acceso (correo y contraseña). 5. El sistema valida los datos ingresados y aprovisiona el espacio de trabajo del despacho con el plan seleccionado. 6. El sistema envía un correo de confirmación al socio fundador. 7. El socio fundador confirma su correo electrónico mediante el enlace recibido. 8. El sistema activa la cuenta y redirige al tablero de control del despacho. \\
    Flujos alternos & 4a. La contraseña no cumple las políticas de complejidad → el sistema indica los requisitos faltantes. 5a. El correo electrónico ya está registrado → el sistema informa que la cuenta existe y ofrece la opción de recuperación. 7a. El enlace de confirmación ha expirado → el sistema permite reenviar el correo de confirmación. \\
    Postcondiciones & El despacho queda registrado con su espacio de datos inicializado, el plan activo y el socio fundador autenticado con rol Owner. Las credenciales se almacenan mediante funciones de hashing resistentes. \\
    Derivado de & HU-01, HU-02 \\
    RF relacionados & RF-01, RF-02 \\
    RNF relacionados & RNF-01, RNF-02 \\
    \bottomrule
  \end{tabularx}
\end{table}

\begin{figure}[H]
  \centering
  \includegraphics[width=0.9\linewidth]{image5.png}
  \caption{Diagrama CU-01. Registro de despacho y selección de plan.}
  \label{fig:cu-01}
\end{figure}

\subsubsection{CU-02 --- Gestión del ciclo de vida de miembros del equipo}

\begin{table}[H]
  \centering
  \caption{Caso de uso CU-02: Gestión del ciclo de vida de miembros del equipo.}
  \label{tab:cu-02}
  \begin{tabularx}{\linewidth}{@{}>{\bfseries}l X@{}}
    \toprule
    ID & \textbf{CU-02} \\
    \midrule
    Nombre & Gestión del ciclo de vida de miembros del equipo \\
    Actor principal & Owner \\
    Actores secundarios & Colaborador invitado, sistema de correo transaccional \\
    Precondiciones & El Owner tiene sesión activa. \\
    Flujo principal (alta) & 1. El Owner accede al módulo de gestión de equipo. 2. Ingresa el correo electrónico del nuevo miembro y selecciona el perfil a asignar (Litigante o Paralegal). 3. El sistema verifica que el despacho no ha alcanzado el límite de usuarios del plan contratado. 4. El sistema genera un enlace de invitación con vigencia limitada y lo envía por correo. 5. El colaborador invitado accede al enlace, completa su registro con sus datos personales y establece su contraseña. 6. El sistema valida los datos, crea la cuenta asociada al despacho con el perfil indicado y registra la operación en la bitácora. 7. El nuevo miembro accede al sistema con las funciones habilitadas según su perfil. \\
    Flujo principal (baja) & 1. El Owner accede al módulo de gestión de equipo. 2. Selecciona al miembro que desea dar de baja. 3. El sistema solicita confirmación de la acción. 4. El Owner confirma la baja. 5. El sistema revoca de inmediato el acceso del colaborador, invalida sus sesiones activas y registra la operación en la bitácora. 6. Los expedientes y documentos asociados al colaborador permanecen en el sistema sin modificación; únicamente se revoca el acceso. \\
    Flujos alternos & 2a. (Alta) El correo ya pertenece a un miembro activo del despacho → el sistema rechaza la invitación e informa el motivo. 3a. (Alta) El despacho ha alcanzado el límite de usuarios → el sistema bloquea la invitación e indica la necesidad de actualizar el plan. 5a. (Alta) El enlace de invitación ha expirado → el sistema muestra un mensaje informativo; el Owner debe reenviar la invitación. 2b. (Baja) El Owner intenta darse de baja a sí mismo → el sistema rechaza la operación, ya que el despacho requiere al menos un Owner activo. \\
    Postcondiciones & (Alta) El colaborador queda registrado en el despacho con el perfil asignado. (Baja) El acceso del colaborador queda revocado de inmediato. En ambos casos, la operación se registra en la bitácora de actividad. \\
    Derivado de & HU-03, HU-04, HU-11 \\
    RF relacionados & RF-03 \\
    RNF relacionados & RNF-01, RNF-02 \\
    \bottomrule
  \end{tabularx}
\end{table}

\begin{figure}[H]
  \centering
  \includegraphics[width=0.9\linewidth]{image6.png}
  \caption{Diagrama CU-02. Gestión del ciclo de vida de miembros del equipo.}
  \label{fig:cu-02}
\end{figure}

\subsubsection{CU-03 --- Inicio de sesión y gestión de sesiones}

\begin{table}[H]
  \centering
  \caption{Caso de uso CU-03: Inicio de sesión y gestión de sesiones.}
  \label{tab:cu-03}
  \begin{tabularx}{\linewidth}{@{}>{\bfseries}l X@{}}
    \toprule
    ID & \textbf{CU-03} \\
    \midrule
    Nombre & Inicio de sesión y gestión de sesiones \\
    Actor principal & Usuario del sistema (Owner, Litigante, Paralegal, Cliente) \\
    Actores secundarios & --- \\
    Precondiciones & El usuario posee una cuenta activa en el sistema. \\
    Flujo principal & 1. El usuario accede a la pantalla de inicio de sesión. 2. Ingresa su correo electrónico y contraseña. 3. El sistema verifica las credenciales contra el hash almacenado. 4. Si el usuario tiene habilitada la verificación en dos pasos, el sistema solicita el código temporal generado en su dispositivo. 5. El usuario ingresa el código y el sistema lo valida. 6. El sistema crea una sesión autenticada, registra el inicio de sesión en la bitácora y redirige al usuario al módulo correspondiente a su perfil. \\
    Flujos alternos & 3a. Las credenciales son incorrectas → el sistema informa el error sin revelar cuál campo es inválido. 3b. El usuario es Owner y elige autenticarse con certificado digital → el sistema genera un desafío criptográfico; el usuario firma con su clave privada; el sistema verifica la firma contra el certificado registrado. 1a. El usuario selecciona la opción de recuperación de contraseña → el sistema solicita el correo electrónico asociado a la cuenta → envía un enlace de restablecimiento con vigencia limitada → el usuario accede al enlace y establece una nueva contraseña → el sistema la almacena mediante funciones de hashing resistentes, invalida las sesiones activas previas y confirma el restablecimiento. 5a. El código temporal es inválido o ha expirado → el sistema rechaza el acceso y permite reintentar. 6a. La sesión permanece inactiva por el tiempo configurado → el sistema la invalida y requiere nueva autenticación, preservando los datos no guardados. \\
    Postcondiciones & El usuario queda autenticado con una sesión activa. El inicio de sesión, incluyendo método de autenticación utilizado y dirección IP, se registra en la bitácora. \\
    Derivado de & HU-05, HU-06, HU-07, HU-08, HU-09 \\
    RF relacionados & RF-01, RF-02 \\
    RNF relacionados & RNF-02, RNF-05 \\
    \bottomrule
  \end{tabularx}
\end{table}

\begin{figure}[H]
  \centering
  \includegraphics[width=0.9\linewidth]{image7.png}
  \caption{Diagrama CU-03. Inicio de sesión y gestión de sesiones.}
  \label{fig:cu-03}
\end{figure}

\subsubsection{CU-04 --- Control de acceso por perfil}

\begin{table}[H]
  \centering
  \caption{Caso de uso CU-04: Control de acceso por perfil.}
  \label{tab:cu-04}
  \begin{tabularx}{\linewidth}{@{}>{\bfseries}l X@{}}
    \toprule
    ID & \textbf{CU-04} \\
    \midrule
    Nombre & Control de acceso por perfil \\
    Actor principal & Usuario autenticado (cualquier rol) \\
    Actores secundarios & --- \\
    Precondiciones & El usuario tiene sesión activa. \\
    Flujo principal & 1. El usuario solicita acceso a un módulo o recurso del sistema. 2. El sistema evalúa el perfil del usuario contra la matriz de permisos del módulo solicitado. 3. Si el perfil tiene permiso, el sistema concede el acceso y presenta la interfaz correspondiente. \\
    Flujos alternos & 3a. El perfil no tiene permiso sobre el módulo solicitado → el sistema deniega el acceso con un mensaje explicativo y registra el intento en la bitácora. 3b. El usuario intenta acceder a un recurso fuera de su ámbito de autorización (por ejemplo, manipulando identificadores de expedientes no asignados) → el sistema devuelve un error de autorización sin revelar la existencia del recurso. \\
    Postcondiciones & El acceso queda concedido o denegado. Cada intento de acceso, exitoso o fallido, se registra en la bitácora. \\
    Derivado de & HU-10 \\
    RF relacionados & RF-02 \\
    RNF relacionados & RNF-01, RNF-05 \\
    \bottomrule
  \end{tabularx}
\end{table}

\begin{figure}[H]
  \centering
  \includegraphics[width=0.9\linewidth]{image8.png}
  \caption{Diagrama CU-04. Control de acceso por perfil.}
  \label{fig:cu-04}
\end{figure}

\subsection{Gestión Procesal Penal (Ciclo de Vida del Expediente)}\label{sec:cu-gestion-procesal}

\subsubsection{CU-05 --- Registro y administración de expediente penal}

\begin{table}[H]
  \centering
  \caption{Caso de uso CU-05: Registro y administración de expediente penal.}
  \label{tab:cu-05}
  \begin{tabularx}{\linewidth}{@{}>{\bfseries}l X@{}}
    \toprule
    ID & \textbf{CU-05} \\
    \midrule
    Nombre & Registro y administración de expediente penal \\
    Actor principal & Litigante \\
    Actores secundarios & Paralegal (consulta) \\
    Precondiciones & El Litigante tiene sesión activa y pertenece al despacho. \\
    Flujo principal & 1. El Litigante accede al módulo de gestión procesal y selecciona la opción de nuevo expediente. 2. Ingresa los identificadores oficiales: Número Único de Caso (NUC) y Carpeta Judicial. 3. Registra el tipo de delito, los datos generales del asunto y los metadatos asociados. 4. El sistema valida que los campos obligatorios estén completos y que los identificadores no estén duplicados dentro del despacho. 5. El sistema persiste el expediente con estado inicial “Investigación”, genera la marca de tiempo de creación y registra la operación en la bitácora. \\
    Flujos alternos & 4a. Faltan campos obligatorios (NUC o Carpeta Judicial) → el sistema rechaza el registro e indica los campos faltantes. 4b. El NUC ya existe en el despacho → el sistema informa la duplicidad. \\
    Postcondiciones & El expediente queda registrado con estado “Investigación”, disponible para los usuarios autorizados del despacho. \\
    Derivado de & HU-12, HU-17, HU-18 \\
    RF relacionados & RF-04 \\
    RNF relacionados & RNF-04, RNF-10 \\
    \bottomrule
  \end{tabularx}
\end{table}

\begin{figure}[H]
  \centering
  \includegraphics[width=0.9\linewidth]{image9.png}
  \caption{Diagrama CU-05. Registro y administración de expediente penal.}
  \label{fig:cu-05}
\end{figure}

\subsubsection{CU-06 --- Gestión del directorio de participantes}

\begin{table}[H]
  \centering
  \caption{Caso de uso CU-06: Gestión del directorio de participantes.}
  \label{tab:cu-06}
  \begin{tabularx}{\linewidth}{@{}>{\bfseries}l X@{}}
    \toprule
    ID & \textbf{CU-06} \\
    \midrule
    Nombre & Gestión del directorio de participantes \\
    Actor principal & Litigante \\
    Actores secundarios & --- \\
    Precondiciones & Existe un expediente activo al cual vincular participantes. \\
    Flujo principal & 1. El Litigante accede al expediente y selecciona el módulo de directorio de participantes. 2. Selecciona el tipo de sujeto procesal (imputado, víctima, defensor, juez de control, Ministerio Público, entre otros). 3. Ingresa los datos identificadores del participante según su tipo (nombre, CURP, cédula profesional, órgano jurisdiccional, según aplique). 4. El sistema valida los datos requeridos según el tipo de sujeto procesal y registra al participante vinculado al expediente. \\
    Flujos alternos & 3a. Para un juez de control, no se proporciona el órgano jurisdiccional ni la FIREL → el sistema rechaza el registro e indica los campos faltantes. 4a. El participante ya existe en el directorio del caso con el mismo rol → el sistema informa la duplicidad. \\
    Postcondiciones & El participante queda vinculado al expediente con su rol procesal. Un mismo individuo puede tener roles distintos en expedientes diferentes. \\
    Derivado de & HU-13 \\
    RF relacionados & RF-06 \\
    RNF relacionados & RNF-10 \\
    \bottomrule
  \end{tabularx}
\end{table}

\begin{figure}[H]
  \centering
  \includegraphics[width=0.9\linewidth]{image10.png}
  \caption{Diagrama CU-06. Gestión del directorio de participantes.}
  \label{fig:cu-06}
\end{figure}

\subsubsection{CU-07 --- Transición de etapa procesal}

\begin{table}[H]
  \centering
  \caption{Caso de uso CU-07: Transición de etapa procesal.}
  \label{tab:cu-07}
  \begin{tabularx}{\linewidth}{@{}>{\bfseries}l X@{}}
    \toprule
    ID & \textbf{CU-07} \\
    \midrule
    Nombre & Transición de etapa procesal \\
    Actor principal & Litigante \\
    Actores secundarios & --- \\
    Precondiciones & El expediente se encuentra en una etapa procesal que admite transición (Investigación, Intermedia o Juicio Oral). \\
    Flujo principal & 1. El Litigante accede al expediente y solicita avanzar a la siguiente etapa procesal. 2. El sistema solicita la carga del documento habilitante correspondiente (escrito de acusación para Investigación → Intermedia; auto de apertura para Intermedia → Juicio Oral). 3. El Litigante carga el documento procesal. 4. El sistema almacena el documento con las protecciones de integridad y confidencialidad correspondientes. 5. El sistema registra la transición de forma inmutable: marca de tiempo, usuario responsable y documento asociado. 6. El sistema habilita los módulos correspondientes a la nueva etapa. \\
    Flujos alternos & 2a. El Litigante intenta avanzar de etapa sin cargar el documento habilitante → el sistema rechaza la transición e indica el documento requerido. 3a. El documento cargado no corresponde al formato admitido → el sistema rechaza la carga. \\
    Postcondiciones & El expediente se encuentra en la nueva etapa procesal. La transición queda registrada de forma inmutable en el historial del expediente y no puede revertirse. \\
    Derivado de & HU-14, HU-15, HU-16 \\
    RF relacionados & RF-05 \\
    RNF relacionados & RNF-04, RNF-10 \\
    \bottomrule
  \end{tabularx}
\end{table}

\begin{figure}[H]
  \centering
  \includegraphics[width=0.9\linewidth]{image11.png}
  \caption{Diagrama CU-07. Transición de etapa procesal.}
  \label{fig:cu-07}
\end{figure}

\subsection{Calendario Jurídico, Plazos y Alertas}\label{sec:cu-calendario}

\subsubsection{CU-08 --- Programación de audiencia y activación de plazos}

\begin{table}[H]
  \centering
  \caption{Caso de uso CU-08: Programación de audiencia y activación de plazos.}
  \label{tab:cu-08}
  \begin{tabularx}{\linewidth}{@{}>{\bfseries}l X@{}}
    \toprule
    ID & \textbf{CU-08} \\
    \midrule
    Nombre & Programación de audiencia y activación de plazos \\
    Actor principal & Litigante \\
    Actores secundarios & Sistema de cómputo de plazos \\
    Precondiciones & Existe un expediente activo con una etapa procesal que admite la audiencia a programar. \\
    Flujo principal & 1. El Litigante accede al expediente y selecciona la opción de programar audiencia. 2. Selecciona el tipo de audiencia del catálogo tipificado conforme al CNPP (inicial, intermedia, juicio oral, individualización de sanciones). 3. Ingresa fecha, hora, participantes y modalidad (presencial o videoconferencia). 4. El sistema valida la coherencia procesal (el tipo de audiencia corresponde a la etapa vigente del expediente). 5. El sistema registra la audiencia vinculada al expediente. 6. Al registrar el resultado de la audiencia (por ejemplo, vinculación a proceso), el sistema activa automáticamente los plazos derivados, calculando las fechas de vencimiento conforme a las reglas del artículo 94 del CNPP (véase §2.2.5) y al calendario de días inhábiles configurado. \\
    Flujos alternos & 4a. El tipo de audiencia no corresponde a la etapa procesal vigente → el sistema rechaza la programación. 6a. El resultado de la audiencia no genera plazos derivados → el sistema registra el resultado sin activar cómputos adicionales. \\
    Postcondiciones & La audiencia queda registrada y visible en el calendario del despacho. Los plazos derivados, si los hay, quedan activos con sus fechas de vencimiento calculadas y las alertas preventivas programadas. \\
    Derivado de & HU-19, HU-20, HU-21 \\
    RF relacionados & RF-07, RF-08 \\
    RNF relacionados & RNF-10 \\
    \bottomrule
  \end{tabularx}
\end{table}

\begin{figure}[H]
  \centering
  \includegraphics[width=0.9\linewidth]{image12.png}
  \caption{Diagrama CU-08. Programación de audiencia y activación de plazos.}
  \label{fig:cu-08}
\end{figure}

\subsubsection{CU-09 --- Configuración del calendario judicial}

\begin{table}[H]
  \centering
  \caption{Caso de uso CU-09: Configuración del calendario judicial.}
  \label{tab:cu-09}
  \begin{tabularx}{\linewidth}{@{}>{\bfseries}l X@{}}
    \toprule
    ID & \textbf{CU-09} \\
    \midrule
    Nombre & Configuración del calendario judicial \\
    Actor principal & Owner \\
    Actores secundarios & --- \\
    Precondiciones & El Owner tiene sesión activa. \\
    Flujo principal & 1. El Owner accede al módulo de configuración del despacho. 2. Selecciona la entidad federativa y el fuero para los cuales configurará los días inhábiles. 3. Registra los periodos de receso judicial, días feriados y días inhábiles específicos de la jurisdicción. 4. El sistema almacena la configuración y recalcula los plazos activos que dependan de ese calendario. \\
    Flujos alternos & 3a. Se registra un día inhábil que afecta un plazo con vencimiento inminente (menos de 48 horas) → el sistema recalcula la fecha y emite una alerta al Litigante responsable. \\
    Postcondiciones & El calendario judicial del despacho queda actualizado. Todos los cómputos de plazos futuros utilizan la configuración vigente. \\
    Derivado de & HU-22 \\
    RF relacionados & RF-08 \\
    RNF relacionados & RNF-10 \\
    \bottomrule
  \end{tabularx}
\end{table}

\begin{figure}[H]
  \centering
  \includegraphics[width=0.9\linewidth]{image13.png}
  \caption{Diagrama CU-09. Configuración del calendario judicial.}
  \label{fig:cu-09}
\end{figure}

\subsubsection{CU-10 --- Monitoreo de plazos y emisión de alertas}

\begin{table}[H]
  \centering
  \caption{Caso de uso CU-10: Monitoreo de plazos y emisión de alertas.}
  \label{tab:cu-10}
  \begin{tabularx}{\linewidth}{@{}>{\bfseries}l X@{}}
    \toprule
    ID & \textbf{CU-10} \\
    \midrule
    Nombre & Monitoreo de plazos y emisión de alertas \\
    Actor principal & Sistema (proceso automático) \\
    Actores secundarios & Litigante (receptor de alertas), sistema de correo transaccional \\
    Precondiciones & Existen plazos activos en el sistema. \\
    Flujo principal & 1. El sistema ejecuta periódicamente la verificación de plazos activos. 2. Para cada plazo próximo a vencer (dentro del umbral configurado: 48 h y 24 h por defecto), el sistema genera una notificación. 3. La notificación se envía al Litigante responsable por correo electrónico y se muestra como alerta en el tablero de control. 4. El envío de la notificación se registra en la bitácora de actividad. \\
    Flujos alternos & 2a. El plazo ya venció sin que se registrara la actuación correspondiente → el sistema emite una alerta de nivel crítico. \\
    Postcondiciones & El Litigante ha sido notificado de la proximidad del vencimiento. El envío queda registrado en la bitácora. \\
    Derivado de & HU-23, HU-24 \\
    RF relacionados & RF-09, RF-10 \\
    RNF relacionados & RNF-06 \\
    \bottomrule
  \end{tabularx}
\end{table}

\begin{figure}[H]
  \centering
  \includegraphics[width=0.9\linewidth]{image14.png}
  \caption{Diagrama CU-10. Monitoreo de plazos y emisión de alertas.}
  \label{fig:cu-10}
\end{figure}

\subsection{Gestión Documental y Cadena de Custodia}\label{sec:cu-documental}

\subsubsection{CU-11 --- Carga y clasificación de documentos}

\begin{table}[H]
  \centering
  \caption{Caso de uso CU-11: Carga y clasificación de documentos.}
  \label{tab:cu-11}
  \begin{tabularx}{\linewidth}{@{}>{\bfseries}l X@{}}
    \toprule
    ID & \textbf{CU-11} \\
    \midrule
    Nombre & Carga y clasificación de documentos \\
    Actor principal & Litigante \\
    Actores secundarios & Paralegal (consulta posterior) \\
    Precondiciones & Existe un expediente al cual asociar el documento. El usuario tiene perfil con permisos de carga. \\
    Flujo principal & 1. El Litigante accede al repositorio documental del expediente. 2. Selecciona el archivo a cargar (PDF, DOCX u otro formato admitido). 3. Asigna tipo documental, clasificación y etiquetas al documento. 4. El sistema calcula la huella digital del archivo al momento de la carga. 5. El sistema cifra el documento y lo almacena en el repositorio. 6. El sistema registra en la cadena de custodia: tipo de operación (carga), usuario responsable, marca de tiempo, huella digital del documento y estado del documento. \\
    Flujos alternos & 2a. El formato del archivo no está admitido → el sistema rechaza la carga e informa los formatos permitidos. 2b. El archivo excede el tamaño máximo admitido → el sistema rechaza la carga. \\
    Postcondiciones & El documento queda almacenado cifrado, indexado con su huella digital y disponible para consulta por los perfiles autorizados. La operación queda registrada en la cadena de custodia. \\
    Derivado de & HU-25, HU-26, HU-28, HU-30, HU-32 \\
    RF relacionados & RF-12, RF-13, RF-14 \\
    RNF relacionados & RNF-03, RNF-04 \\
    \bottomrule
  \end{tabularx}
\end{table}

\begin{figure}[H]
  \centering
  \includegraphics[width=0.9\linewidth]{image15.png}
  \caption{Diagrama CU-11. Carga y clasificación de documentos.}
  \label{fig:cu-11}
\end{figure}

\subsubsection{CU-12 --- Consulta de documentos con verificación de integridad}

\begin{table}[H]
  \centering
  \caption{Caso de uso CU-12: Consulta de documentos con verificación de integridad.}
  \label{tab:cu-12}
  \begin{tabularx}{\linewidth}{@{}>{\bfseries}l X@{}}
    \toprule
    ID & \textbf{CU-12} \\
    \midrule
    Nombre & Consulta de documentos con verificación de integridad \\
    Actor principal & Litigante o Paralegal \\
    Actores secundarios & --- \\
    Precondiciones & El usuario tiene perfil con acceso al expediente que contiene el documento. \\
    Flujo principal & 1. El usuario accede al repositorio documental del expediente. 2. Selecciona el documento a consultar. 3. El sistema verifica que el perfil del usuario tiene permiso de acceso al expediente. 4. El sistema descifra el documento. 5. El sistema recalcula la huella digital del documento descifrado y la compara con el valor de referencia almacenado al momento de la carga. 6. Si las huellas coinciden, el sistema entrega el documento al usuario. 7. La operación de consulta se registra en la cadena de custodia. \\
    Flujos alternos & 3a. El usuario no tiene permiso sobre el expediente → el sistema deniega el acceso. 4a. La verificación de la etiqueta de autenticación del cifrado falla (el contenido cifrado fue manipulado en el almacenamiento) → el sistema rechaza el descifrado, notifica al usuario y al Owner, y registra el incidente como alerta de seguridad. 5a. La huella digital recalculada no coincide con la referencia → el sistema alerta al usuario de que el documento fue modificado, bloquea la entrega y registra la discrepancia. \\
    Postcondiciones & El usuario recibe el documento íntegro, o recibe una alerta explícita de alteración. La consulta y su resultado quedan registrados en la cadena de custodia. \\
    Derivado de & HU-27, HU-29, HU-31, HU-33 \\
    RF relacionados & RF-11, RF-12, RF-13, RF-14 \\
    RNF relacionados & RNF-03, RNF-04 \\
    \bottomrule
  \end{tabularx}
\end{table}

\begin{figure}[H]
  \centering
  \includegraphics[width=0.9\linewidth]{image16.png}
  \caption{Diagrama CU-12. Consulta de documentos con verificación de integridad.}
  \label{fig:cu-12}
\end{figure}

\subsection{Firma de Contratos y Constancia de Fecha}\label{sec:cu-firma}

\subsubsection{CU-13 --- Firma digital de contrato con sellado de tiempo}

\begin{table}[H]
  \centering
  \caption{Caso de uso CU-13: Firma digital de contrato con sellado de tiempo.}
  \label{tab:cu-13}
  \begin{tabularx}{\linewidth}{@{}>{\bfseries}l X@{}}
    \toprule
    ID & \textbf{CU-13} \\
    \midrule
    Nombre & Firma digital de contrato con sellado de tiempo \\
    Actor principal & Litigante (firmante) \\
    Actores secundarios & TSA configurada (autoridad local en la demostración; PSC autorizado solo en una integración futura) \\
    Precondiciones & El contrato existe en el repositorio en formato PDF y está listo para firma. El firmante tiene un certificado digital X.509 vigente registrado en el sistema. \\
    Flujo principal & 1. El Litigante accede al contrato y selecciona la opción de firma digital. 2. El sistema presenta el contenido del contrato y solicita confirmación explícita de la acción de firma. 3. El Litigante confirma la firma. 4. El sistema calcula la huella digital del contrato. 5. El sistema firma la huella digital con la clave privada del certificado X.509 del firmante, vinculando criptográficamente su identidad con el contenido del documento. 6. El sistema envía la huella digital a la TSA configurada mediante el protocolo de sellado de tiempo. 7. La TSA retorna el token RFC 3161 que acredita criptográficamente el digest y su tiempo de generación. 8. El sistema almacena el contrato firmado junto con la firma digital, los metadatos del certificado y el token de sello de tiempo. 9. La operación se registra en la bitácora de actividad y en la cadena de custodia del expediente. \\
    Flujos alternos & 2a. El certificado del firmante ha expirado o está revocado → el sistema informa la situación e impide la firma. 6a. La TSA no responde dentro del tiempo límite → el sistema informa o encola la solicitud según la política explícita del emisor; no conmuta silenciosamente a otra autoridad. 7a. El token recibido no supera la validación de formato → el sistema rechaza el sello y registra el error. \\
    Postcondiciones & El contrato queda firmado con evidencia criptográfica verificable de la identidad del firmante y del momento de la firma. Si la autoridad es la TSA local, el resultado acredita el flujo técnico RFC 3161, pero no constituye por sí mismo una constancia NOM-151 de un PSC autorizado. \\
    Derivado de & HU-34, HU-35 \\
    RF relacionados & RF-15, RF-16 \\
    RNF relacionados & RNF-03, RNF-11 \\
    \bottomrule
  \end{tabularx}
\end{table}

\begin{figure}[H]
  \centering
  \includegraphics[width=0.9\linewidth]{image17.png}
  \caption{Diagrama CU-13. Firma digital de contrato con sellado de tiempo.}
  \label{fig:cu-13}
\end{figure}

\subsubsection{CU-14 --- Verificación de firma digital y sello de tiempo}

\begin{table}[H]
  \centering
  \caption{Caso de uso CU-14: Verificación de firma digital y sello de tiempo.}
  \label{tab:cu-14}
  \begin{tabularx}{\linewidth}{@{}>{\bfseries}l X@{}}
    \toprule
    ID & \textbf{CU-14} \\
    \midrule
    Nombre & Verificación de firma digital y sello de tiempo \\
    Actor principal & Usuario autenticado con acceso al contrato \\
    Actores secundarios & --- \\
    Precondiciones & El contrato ha sido firmado digitalmente y posee un token de sello de tiempo asociado. \\
    Flujo principal & 1. El usuario accede al contrato firmado y selecciona la opción de verificación. 2. El sistema recalcula la huella digital del contrato actual. 3. El sistema verifica la firma digital: descifra la firma con la clave pública del certificado X.509 del firmante y compara el resultado con la huella digital recalculada. 4. El sistema valida el estado del certificado del firmante (vigente, revocado o expirado) consultando la cadena de certificación. 5. El sistema verifica la autenticidad del token de sello de tiempo validando la firma del PSC. 6. El sistema presenta al usuario un reporte de verificación con el resultado de cada componente. \\
    Flujos alternos & 3a. La huella digital recalculada no coincide con la contenida en la firma → el sistema reporta que el documento fue modificado después de firmado e indica la causa de la invalidación. 4a. El certificado del firmante ha expirado o fue revocado → el sistema lo indica en el reporte sin invalidar necesariamente la firma (la firma era válida al momento del sellado). 5a. El token de sello de tiempo no supera la validación → el sistema reporta que la constancia de fecha no es verificable. \\
    Postcondiciones & El usuario dispone de un reporte de verificación con el estado de cada componente: validez de la firma, estado del certificado, validez del sello de tiempo e integridad del documento. \\
    Derivado de & HU-36, HU-37 \\
    RF relacionados & RF-17 \\
    RNF relacionados & RNF-11 \\
    \bottomrule
  \end{tabularx}
\end{table}

\begin{figure}[H]
  \centering
  \includegraphics[width=0.9\linewidth]{image18.png}
  \caption{Diagrama CU-14. Verificación de firma digital y sello de tiempo.}
  \label{fig:cu-14}
\end{figure}

\subsection{Operación y Monitoreo}\label{sec:cu-operacion}

\subsubsection{CU-15 --- Consulta del tablero de control}

\begin{table}[H]
  \centering
  \caption{Caso de uso CU-15: Consulta del tablero de control.}
  \label{tab:cu-15}
  \begin{tabularx}{\linewidth}{@{}>{\bfseries}l X@{}}
    \toprule
    ID & \textbf{CU-15} \\
    \midrule
    Nombre & Consulta del tablero de control \\
    Actor principal & Owner o Litigante \\
    Actores secundarios & --- \\
    Precondiciones & El usuario tiene sesión activa. \\
    Flujo principal & 1. El usuario accede al tablero de control del despacho. 2. El sistema consulta los indicadores operativos del despacho: casos activos, plazos por vencer en 48 h y 7 días, contratos pendientes de firma y carga de trabajo por abogado. 3. El sistema presenta los indicadores con diferenciación visual por nivel de urgencia. \\
    Flujos alternos & 2a. Existen plazos fatales por vencer en menos de 48 horas → el sistema los destaca con indicador de alerta de nivel crítico, diferenciado visualmente del resto de elementos. \\
    Postcondiciones & El usuario visualiza el estado actual del despacho. \\
    Derivado de & HU-38 \\
    RF relacionados & RF-18 \\
    RNF relacionados & RNF-07 \\
    \bottomrule
  \end{tabularx}
\end{table}

\begin{figure}[H]
  \centering
  \includegraphics[width=0.9\linewidth]{image19.png}
  \caption{Diagrama CU-15. Consulta del tablero de control.}
  \label{fig:cu-15}
\end{figure}

\subsubsection{CU-16 --- Generación de informes}

\begin{table}[H]
  \centering
  \caption{Caso de uso CU-16: Generación de informes.}
  \label{tab:cu-16}
  \begin{tabularx}{\linewidth}{@{}>{\bfseries}l X@{}}
    \toprule
    ID & \textbf{CU-16} \\
    \midrule
    Nombre & Generación de informes \\
    Actor principal & Owner \\
    Actores secundarios & Sistema de colas de trabajo \\
    Precondiciones & El Owner tiene sesión activa. \\
    Flujo principal & 1. El Owner accede al módulo de reportes. 2. Selecciona el tipo de informe y los filtros aplicables (rango de fechas, abogado responsable, estado de los casos). 3. El sistema encola la solicitud de generación como tarea en segundo plano. 4. El sistema confirma la recepción de la solicitud y muestra un indicador de progreso. 5. Al completar la generación, el sistema notifica al Owner y pone el informe disponible para descarga en formato PDF o CSV. \\
    Flujos alternos & 3a. El volumen de datos supera el umbral de procesamiento inmediato → el sistema informa el tiempo estimado y envía notificación por correo cuando el reporte esté listo. \\
    Postcondiciones & El informe queda disponible para descarga. La interfaz permanece responsiva durante el procesamiento. \\
    Derivado de & HU-39 \\
    RF relacionados & RF-19 \\
    RNF relacionados & RNF-06, RNF-07 \\
    \bottomrule
  \end{tabularx}
\end{table}

\begin{figure}[H]
  \centering
  \includegraphics[width=0.9\linewidth]{image20.png}
  \caption{Diagrama CU-16. Generación de informes.}
  \label{fig:cu-16}
\end{figure}

\subsubsection{CU-17 --- Consulta de bitácora de actividad}

\begin{table}[H]
  \centering
  \caption{Caso de uso CU-17: Consulta de bitácora de actividad.}
  \label{tab:cu-17}
  \begin{tabularx}{\linewidth}{@{}>{\bfseries}l X@{}}
    \toprule
    ID & \textbf{CU-17} \\
    \midrule
    Nombre & Consulta de bitácora de actividad \\
    Actor principal & Owner \\
    Actores secundarios & --- \\
    Precondiciones & El Owner tiene sesión activa. \\
    Flujo principal & 1. El Owner accede al módulo de bitácora de actividad. 2. Aplica filtros de búsqueda (rango de fechas, usuario, tipo de operación, recurso afectado). 3. El sistema presenta los registros cronológicamente, mostrando para cada entrada: identificador del usuario, tipo de operación, recurso afectado, marca de tiempo y dirección IP. \\
    Flujos alternos & --- (la bitácora es de solo lectura; no existen operaciones de modificación o eliminación). \\
    Postcondiciones & El Owner visualiza los registros solicitados. Los registros de la bitácora no pueden modificarse ni eliminarse por ningún usuario del sistema, incluido el Owner. \\
    Derivado de & HU-40 \\
    RF relacionados & RF-20 \\
    RNF relacionados & RNF-04 \\
    \bottomrule
  \end{tabularx}
\end{table}

\begin{figure}[H]
  \centering
  \includegraphics[width=0.9\linewidth]{image22.png}
  \caption{Diagrama CU-17. Consulta de bitácora de actividad.}
  \label{fig:cu-17}
\end{figure}


\section{Requerimientos del Sistema}\label{sec:requerimientos}

Los requerimientos se derivan de los objetivos definidos anteriormente y se organizan en funcionales y no funcionales. Los requerimientos funcionales definen los servicios que el sistema debe proporcionar y cómo debe reaccionar ante entradas particulares. Los requerimientos no funcionales describen restricciones y atributos de calidad que limitan las decisiones de diseño e implementación.

\subsection{Requerimientos Funcionales}\label{sec:requerimientos-funcionales}

Los requerimientos funcionales especifican los servicios que el sistema debe proporcionar, las entradas que acepta, las salidas que produce y el comportamiento esperado ante condiciones particulares. Cada requerimiento se deriva de las historias de usuario definidas en la \cref{sec:historias-usuario} y se formaliza mediante una estructura que incluye descripción, entradas, salidas, actores involucrados y criterio de aceptación verificable. Los requerimientos se agrupan en cinco categorías funcionales alineadas con los módulos del prototipo.

\subsubsection{Autenticación y Control de Acceso}

\begin{table}[H]
  \centering
  \caption{RF-01: Autenticación y gestión de sesiones}
  \label{tab:rf-01}
  \begin{tabularx}{\linewidth}{@{}L{3.2cm}X@{}}
    \toprule
    \textbf{ID} & RF-01 \\
    \midrule
    \textbf{Nombre} & Autenticación y gestión de sesiones \\
    \textbf{Descripción} & El sistema debe proveer mecanismos seguros de inicio de sesión, recuperación de contraseñas, expiración de sesiones por inactividad y protección contra fijación de sesión. \\
    \textbf{Entrada} & Credenciales del usuario (usuario/contraseña o certificado digital); solicitud de recuperación de contraseña. \\
    \textbf{Salida} & Sesión autenticada con token de acceso; notificación de expiración; redireccionamiento al inicio de sesión tras inactividad sin pérdida de datos no guardados. \\
    \textbf{Actores involucrados} & Todos los roles (Owner, Litigante, Paralegal, Cliente); sistema de gestión de sesiones. \\
    \textbf{Criterio de aceptación} & Dado un usuario inactivo por el tiempo configurado, cuando vence la sesión, entonces el sistema la invalida automáticamente y requiere nueva autenticación sin pérdida de datos no guardados. \\
    \bottomrule
  \end{tabularx}
\end{table}

\begin{table}[H]
  \centering
  \caption{RF-02: Control de acceso basado en roles (RBAC)}
  \label{tab:rf-02}
  \begin{tabularx}{\linewidth}{@{}L{3.2cm}X@{}}
    \toprule
    \textbf{ID} & RF-02 \\
    \midrule
    \textbf{Nombre} & Control de acceso basado en roles (RBAC) \\
    \textbf{Descripción} & El sistema debe asignar permisos de forma granular a través de roles predefinidos (Owner, Litigante, Paralegal, Cliente) para restringir el acceso a módulos y operaciones sensibles. \\
    \textbf{Entrada} & Rol asignado al usuario; recurso o módulo al que se intenta acceder. \\
    \textbf{Salida} & Permiso concedido o denegado; registro en bitácora del intento de acceso. \\
    \textbf{Actores involucrados} & Owner (configuración de roles); todos los usuarios como sujetos de control; sistema de autorización. \\
    \textbf{Criterio de aceptación} & Dado un usuario con rol Paralegal, cuando intenta acceder a un módulo restringido a Owner o Litigante, entonces el sistema deniega el acceso con mensaje explicativo y registra el intento en la bitácora. \\
    \bottomrule
  \end{tabularx}
\end{table}

\begin{table}[H]
  \centering
  \caption{RF-03: Gestión de equipos colaborativos}
  \label{tab:rf-03}
  \begin{tabularx}{\linewidth}{@{}L{3.2cm}X@{}}
    \toprule
    \textbf{ID} & RF-03 \\
    \midrule
    \textbf{Nombre} & Gestión de equipos colaborativos \\
    \textbf{Descripción} & El sistema debe permitir invitar nuevos miembros al despacho mediante correo electrónico, gestionar la aceptación de invitaciones y administrar el alta y baja de usuarios en el equipo. \\
    \textbf{Entrada} & Correo electrónico del nuevo miembro; rol a asignar; acción de aceptación del invitado. \\
    \textbf{Salida} & Invitación enviada; usuario creado y asociado al despacho; acceso habilitado según el rol asignado. \\
    \textbf{Actores involucrados} & Owner (invita y da de baja); usuario invitado; sistema de correo transaccional. \\
    \textbf{Criterio de aceptación} & Dado un Owner que envía una invitación, cuando el invitado la acepta dentro del periodo de validez del enlace, entonces queda registrado con el rol indicado; si el enlace vence, la invitación queda inactiva automáticamente. \\
    \bottomrule
  \end{tabularx}
\end{table}

\subsubsection{Gestión Procesal Penal}

\begin{table}[H]
  \centering
  \caption{RF-04: Expediente penal digital}
  \label{tab:rf-04}
  \begin{tabularx}{\linewidth}{@{}L{3.2cm}X@{}}
    \toprule
    \textbf{ID} & RF-04 \\
    \midrule
    \textbf{Nombre} & Expediente penal digital \\
    \textbf{Descripción} & El sistema debe permitir el registro y administración integral de casos penales, incluyendo identificadores oficiales (NUC, Carpeta Judicial), delitos y metadatos asociados. \\
    \textbf{Entrada} & NUC, Carpeta Judicial, tipo de delito, datos del caso y metadatos ingresados por el Litigante. \\
    \textbf{Salida} & Expediente registrado con identificadores únicos, estado inicial `Investigación' y timestamp de creación. \\
    \textbf{Actores involucrados} & Litigante, Paralegal; base de datos relacional. \\
    \textbf{Criterio de aceptación} & Dado un Litigante que registra un nuevo caso, cuando omite el NUC o la Carpeta Judicial, entonces el sistema rechaza el registro con indicación del campo faltante; cuando los datos son completos, el expediente queda persistido y disponible en el módulo de gestión procesal. \\
    \bottomrule
  \end{tabularx}
\end{table}

\begin{table}[H]
  \centering
  \caption{RF-05: Control de etapas procesales (CNPP)}
  \label{tab:rf-05}
  \begin{tabularx}{\linewidth}{@{}L{3.2cm}X@{}}
    \toprule
    \textbf{ID} & RF-05 \\
    \midrule
    \textbf{Nombre} & Control de etapas procesales (CNPP) \\
    \textbf{Descripción} & El sistema debe gestionar la transición entre etapas procesales (Investigación, Intermedia, Juicio Oral) conforme al CNPP, manteniendo un historial inmutable de cambios y validaciones procesales. \\
    \textbf{Entrada} & Documento procesal habilitante cargado por el usuario (ej. escrito de acusación para transitar de Investigación a Intermedia). \\
    \textbf{Salida} & Registro inmutable de la transición con marca de tiempo, identificador de usuario y documento asociado; módulo de la siguiente etapa habilitado. \\
    \textbf{Actores involucrados} & Litigante (Owner); sistema de gestión procesal; bitácora de auditoría. \\
    \textbf{Criterio de aceptación} & Dado un expediente en etapa de Investigación, cuando se intenta avanzar sin cargar el escrito de acusación, entonces el sistema rechaza la transición; cuando se carga el documento habilitante, la transición se registra de forma inmutable y no puede revertirse. \\
    \bottomrule
  \end{tabularx}
\end{table}

\begin{table}[H]
  \centering
  \caption{RF-06: Directorio de participantes}
  \label{tab:rf-06}
  \begin{tabularx}{\linewidth}{@{}L{3.2cm}X@{}}
    \toprule
    \textbf{ID} & RF-06 \\
    \midrule
    \textbf{Nombre} & Directorio de participantes \\
    \textbf{Descripción} & El sistema debe permitir la gestión de sujetos procesales vinculados a cada caso (imputados, víctimas, jueces, peritos), definiendo su rol procesal y estatus jurídico. \\
    \textbf{Entrada} & Datos del sujeto procesal (nombre, rol, cédula, CURP u otros identificadores según tipo); caso al que se vincula. \\
    \textbf{Salida} & Participante registrado y vinculado al expediente con su rol procesal; historial de cambios de rol disponible. \\
    \textbf{Actores involucrados} & Litigante; sistema de gestión de directorio. \\
    \textbf{Criterio de aceptación} & Dado un caso activo, cuando se registra un participante con rol `Juez de control', entonces el sistema valida que se proporcione el órgano jurisdiccional y la firma electrónica (FIREL); un mismo usuario puede tener roles distintos en casos diferentes sin conflicto. \\
    \bottomrule
  \end{tabularx}
\end{table}

\begin{table}[H]
  \centering
  \caption{RF-07: Gestión de audiencias especializadas}
  \label{tab:rf-07}
  \begin{tabularx}{\linewidth}{@{}L{3.2cm}X@{}}
    \toprule
    \textbf{ID} & RF-07 \\
    \midrule
    \textbf{Nombre} & Gestión de audiencias especializadas \\
    \textbf{Descripción} & El sistema debe permitir la programación y seguimiento de audiencias tipificadas según el CNPP, registrando resultados, asistentes y acuerdos derivados. \\
    \textbf{Entrada} & Tipo de audiencia (del catálogo CNPP), fecha, hora, participantes y modalidad (presencial o videoconferencia). \\
    \textbf{Salida} & Audiencia agendada vinculada al expediente; activación automática del cómputo de plazos derivados; registro de resultados y acuerdos. \\
    \textbf{Actores involucrados} & Litigante; sistema de calendario; módulo de plazos. \\
    \textbf{Criterio de aceptación} & Dado el registro de una `Audiencia inicial' con vinculación a proceso, cuando se confirma, entonces el sistema activa automáticamente el plazo de investigación complementaria (2 o 6 meses) y genera la alerta preventiva correspondiente. \\
    \bottomrule
  \end{tabularx}
\end{table}

\begin{table}[H]
  \centering
  \caption{RF-08: Blindaje de plazos procesales}
  \label{tab:rf-08}
  \begin{tabularx}{\linewidth}{@{}L{3.2cm}X@{}}
    \toprule
    \textbf{ID} & RF-08 \\
    \midrule
    \textbf{Nombre} & Blindaje de plazos procesales \\
    \textbf{Descripción} & El sistema debe permitir el registro y cálculo de términos constitucionales y plazos fatales, con derivación automática a partir de actuaciones o audiencias registradas. \\
    \textbf{Entrada} & Actuación o audiencia registrada que activa un plazo constitucional; calendario de días inhábiles configurado por entidad federativa. \\
    \textbf{Salida} & Plazo calculado con fecha y hora de vencimiento; alerta programada; registro en el módulo de plazos del expediente. \\
    \textbf{Actores involucrados} & Litigante; sistema de cómputo de plazos; calendario judicial configurable. \\
    \textbf{Criterio de aceptación} & Dado el registro de vinculación a proceso con plazo de 6 meses de investigación complementaria, entonces el sistema calcula la fecha de vencimiento excluyendo correctamente los días inhábiles del calendario configurado para la entidad federativa correspondiente. \\
    \bottomrule
  \end{tabularx}
\end{table}

\begin{table}[H]
  \centering
  \caption{RF-09: Sistema de alertas preventivas}
  \label{tab:rf-09}
  \begin{tabularx}{\linewidth}{@{}L{3.2cm}X@{}}
    \toprule
    \textbf{ID} & RF-09 \\
    \midrule
    \textbf{Nombre} & Sistema de alertas preventivas \\
    \textbf{Descripción} & El sistema debe generar notificaciones automáticas ante la proximidad de vencimientos de plazos o audiencias programadas. \\
    \textbf{Entrada} & Plazos activos registrados; reglas de anticipación configurables (ej. alerta a 48 h y a 24 h del vencimiento). \\
    \textbf{Salida} & Notificación enviada al Litigante responsable por correo y visible en el dashboard del sistema. \\
    \textbf{Actores involucrados} & Litigante; sistema de notificaciones; módulo de colas de trabajo. \\
    \textbf{Criterio de aceptación} & Dado un plazo de vinculación a proceso con vencimiento en 48 horas, cuando el sistema ejecuta la verificación periódica de plazos, entonces envía una notificación de alerta roja al Litigante responsable; el envío se registra en la bitácora. \\
    \bottomrule
  \end{tabularx}
\end{table}

\begin{table}[H]
  \centering
  \caption{RF-10: Calendario jurídico integrado}
  \label{tab:rf-10}
  \begin{tabularx}{\linewidth}{@{}L{3.2cm}X@{}}
    \toprule
    \textbf{ID} & RF-10 \\
    \midrule
    \textbf{Nombre} & Calendario jurídico integrado \\
    \textbf{Descripción} & El sistema debe proveer una visualización unificada de la agenda del despacho, integrando audiencias y plazos límite en un calendario interactivo. \\
    \textbf{Entrada} & Audiencias programadas; plazos activos; configuración de la vista seleccionada (día, semana o mes). \\
    \textbf{Salida} & Calendario unificado con codificación de color por tipo de evento y acceso directo al expediente relacionado desde cada entrada. \\
    \textbf{Actores involucrados} & Litigante, Paralegal, Owner. \\
    \textbf{Criterio de aceptación} & Dado un despacho con 5 o más casos activos, cuando el Litigante accede al calendario, entonces visualiza todos los vencimientos y audiencias de sus casos en una sola vista sin necesidad de navegar por cada expediente individualmente. \\
    \bottomrule
  \end{tabularx}
\end{table}

\subsubsection{Gestión Documental y Evidencias}

\begin{table}[H]
  \centering
  \caption{RF-11: Cadena de custodia digital}
  \label{tab:rf-11}
  \begin{tabularx}{\linewidth}{@{}L{3.2cm}X@{}}
    \toprule
    \textbf{ID} & RF-11 \\
    \midrule
    \textbf{Nombre} & Cadena de custodia digital \\
    \textbf{Descripción} & El sistema debe proveer un módulo para el registro de evidencias físicas y digitales, garantizando la trazabilidad de sus movimientos y resguardos mediante un registro cronológico inmutable. \\
    \textbf{Entrada} & Evidencia física o digital ingresada al sistema; operación de movimiento o consulta de la evidencia. \\
    \textbf{Salida} & Registro cronológico append-only con: tipo de operación, identificador de la evidencia, usuario responsable, timestamp y estado de la evidencia. \\
    \textbf{Actores involucrados} & Litigante, Paralegal; sistema de almacenamiento NoSQL (registro de auditoría). \\
    \textbf{Criterio de aceptación} & Dado un registro de evidencia existente, cuando cualquier usuario intenta modificar o eliminar una entrada de la cadena de custodia, entonces el sistema rechaza la operación; cada nuevo movimiento genera únicamente un registro de inserción nuevo. \\
    \bottomrule
  \end{tabularx}
\end{table}

\begin{table}[H]
  \centering
  \caption{RF-12: Repositorio documental}
  \label{tab:rf-12}
  \begin{tabularx}{\linewidth}{@{}L{3.2cm}X@{}}
    \toprule
    \textbf{ID} & RF-12 \\
    \midrule
    \textbf{Nombre} & Repositorio documental \\
    \textbf{Descripción} & El sistema debe permitir cargar, clasificar, consultar y compartir documentos procesales asociados a expedientes, asegurando su organización, disponibilidad y control de versiones. \\
    \textbf{Entrada} & Archivo digital (PDF, DOCX, multimedia); clasificación y metadatos ingresados por el usuario; expediente al que se asocia. \\
    \textbf{Salida} & Documento almacenado cifrado (AES-256-GCM), indexado con su hash SHA-256 y disponible para consulta por los roles autorizados. \\
    \textbf{Actores involucrados} & Litigante, Paralegal; sistema de almacenamiento cifrado. \\
    \textbf{Criterio de aceptación} & Dado un documento cargado al repositorio, cuando un Paralegal intenta acceder a un documento de un expediente no asignado, entonces el sistema deniega el acceso; el documento debe ser recuperable en menos de 3 segundos para archivos de hasta 10 MB. \\
    \bottomrule
  \end{tabularx}
\end{table}

\subsubsection{Servicios Criptográficos}

\begin{table}[H]
  \centering
  \caption{RF-13: Verificación de integridad de documentos}
  \label{tab:rf-13}
  \begin{tabularx}{\linewidth}{@{}L{3.2cm}X@{}}
    \toprule
    \textbf{ID} & RF-13 \\
    \midrule
    \textbf{Nombre} & Verificación de integridad de documentos \\
    \textbf{Descripción} & El sistema debe calcular y almacenar el hash SHA-256 (véase §2.3.1) de cada documento al momento de su carga. En cada consulta posterior, debe recalcular el hash y compararlo con el valor de referencia, notificando al usuario si se detecta una discrepancia. \\
    \textbf{Entrada} & Archivo digital cargado al sistema (PDF, DOCX u otro formato admitido). \\
    \textbf{Salida} & Hash SHA-256 almacenado en base de datos; alerta visible al usuario en caso de discrepancia detectada en consulta posterior. \\
    \textbf{Actores involucrados} & Litigante, Paralegal; sistema de almacenamiento (backend). \\
    \textbf{Criterio de aceptación} & Dado un documento cargado, cuando se modifica 1 byte del archivo almacenado, entonces el sistema detecta la discrepancia y muestra una alerta en la siguiente consulta sin requerir intervención manual. \\
    \bottomrule
  \end{tabularx}
\end{table}

\begin{table}[H]
  \centering
  \caption{RF-14: Cifrado autenticado de expedientes y documentos}
  \label{tab:rf-14}
  \begin{tabularx}{\linewidth}{@{}L{3.2cm}X@{}}
    \toprule
    \textbf{ID} & RF-14 \\
    \midrule
    \textbf{Nombre} & Cifrado autenticado de expedientes y documentos \\
    \textbf{Descripción} & El sistema debe cifrar los documentos y expedientes almacenados mediante AES-256-GCM (véase §2.3.2), proveyendo simultáneamente confidencialidad e integridad. Cada operación de cifrado debe emplear un vector de inicialización (IV) único generado aleatoriamente. \\
    \textbf{Entrada} & Documento o expediente a cifrar; clave maestra del despacho; IV generado aleatoriamente por operación. \\
    \textbf{Salida} & Contenido cifrado con AES-256-GCM almacenado en el repositorio; IV almacenado de forma separada de los datos cifrados. \\
    \textbf{Actores involucrados} & Sistema de cifrado (backend); módulo de gestión de claves. \\
    \textbf{Criterio de aceptación} & Dado un documento cifrado almacenado, cuando se modifica 1 byte del texto cifrado directamente en el almacenamiento, entonces la verificación de la etiqueta de autenticación GCM falla y el sistema rechaza el descifrado; cada operación debe utilizar un IV único, verificable mediante pruebas TDD. \\
    \bottomrule
  \end{tabularx}
\end{table}

\begin{table}[H]
  \centering
  \caption{RF-15: Firma digital de contratos}
  \label{tab:rf-15}
  \begin{tabularx}{\linewidth}{@{}L{3.2cm}X@{}}
    \toprule
    \textbf{ID} & RF-15 \\
    \midrule
    \textbf{Nombre} & Firma digital de contratos \\
    \textbf{Descripción} & El sistema debe permitir la firma digital de contratos mediante certificados X.509 (véase §2.3.3), vinculando criptográficamente la identidad del firmante con el contenido del documento. \\
    \textbf{Entrada} & Contrato en formato PDF listo para firmar; certificado X.509 del firmante; confirmación de la acción de firma por el usuario. \\
    \textbf{Salida} & Documento firmado con hash SHA-256 cifrado con la clave privada del firmante; metadatos de firma (timestamp local, identificador del certificado) almacenados junto al documento. \\
    \textbf{Actores involucrados} & Litigante (firmante); sistema de firma digital; módulo de certificados X.509. \\
    \textbf{Criterio de aceptación} & Dado un contrato firmado digitalmente, cuando se modifica cualquier byte del documento posterior a la firma, entonces la verificación retorna inválida; la firma solo puede verificarse exitosamente con la clave pública del certificado X.509 del firmante. \\
    \bottomrule
  \end{tabularx}
\end{table}

\begin{table}[H]
  \centering
  \caption{RF-16: Sellado de tiempo (NOM-151)}
  \label{tab:rf-16}
  \begin{tabularx}{\linewidth}{@{}L{3.2cm}X@{}}
    \toprule
    \textbf{ID} & RF-16 \\
    \midrule
    \textbf{Nombre} & Sellado de tiempo (NOM-151) \\
    \textbf{Descripción} & El sistema debe obtener un token RFC 3161 para cada contrato firmado mediante la TSA configurada, acreditando criptográficamente la existencia del digest en un momento determinado. La constancia NOM-151 de un PSC autorizado no se afirma para la TSA local. \\
    \textbf{Entrada} & Hash SHA-256 del contrato firmado y solicitud al adaptador de TSA configurado. \\
    \textbf{Salida} & Token de sello de tiempo (TimeStampToken) en formato ASN.1/.tsr almacenado y vinculado al contrato; la naturaleza técnica o regulada del emisor queda identificada en el paquete. \\
    \textbf{Actores involucrados} & Sistema backend (cliente TSA); TSA local de demostración o PSC autorizado futuro; módulo de firma digital. \\
    \textbf{Criterio de aceptación} & Dado un contrato firmado, cuando se solicita el sello a la TSA configurada, el sistema almacena un token RFC 3161 que cubre el digest y puede verificarse de forma independiente con OpenSSL sin requerir acceso al sistema. La campaña de estabilidad de un PSC externo queda fuera del alcance actual. \\
    \bottomrule
  \end{tabularx}
\end{table}

\begin{table}[H]
  \centering
  \caption{RF-17: Verificación de firmas y sellos de tiempo}
  \label{tab:rf-17}
  \begin{tabularx}{\linewidth}{@{}L{3.2cm}X@{}}
    \toprule
    \textbf{ID} & RF-17 \\
    \midrule
    \textbf{Nombre} & Verificación de firmas y sellos de tiempo \\
    \textbf{Descripción} & El sistema debe permitir verificar la validez de las firmas digitales y los sellos de tiempo asociados a un contrato, incluyendo la cadena de certificados X.509, vigencia del certificado e integridad del documento. \\
    \textbf{Entrada} & Contrato firmado con su token de sello de tiempo; cadena de certificados X.509 del firmante. \\
    \textbf{Salida} & Resultado de verificación con detalle de: validez de la firma, estado del certificado (vigente/revocado/expirado), validez del sello de tiempo e integridad del documento. \\
    \textbf{Actores involucrados} & Cualquier usuario autenticado con acceso al contrato; sistema de verificación criptográfica. \\
    \textbf{Criterio de aceptación} & Dado un contrato con firma válida e integridad verificada, cuando se ejecuta la verificación, entonces todos los componentes se reportan como válidos; dado un contrato modificado, cuando se verifica, entonces el sistema reporta la firma como inválida con indicación explícita de la causa. \\
    \bottomrule
  \end{tabularx}
\end{table}

\subsubsection{Operación y Monitoreo}

\begin{table}[H]
  \centering
  \caption{RF-18: Dashboard ejecutivo}
  \label{tab:rf-18}
  \begin{tabularx}{\linewidth}{@{}L{3.2cm}X@{}}
    \toprule
    \textbf{ID} & RF-18 \\
    \midrule
    \textbf{Nombre} & Dashboard ejecutivo \\
    \textbf{Descripción} & El sistema debe proveer un tablero de control que visualice indicadores operativos clave: casos activos, plazos próximos y carga de trabajo por abogado. \\
    \textbf{Entrada} & Datos de casos activos, plazos próximos y asignación de carga del despacho. \\
    \textbf{Salida} & Tablero visual con indicadores actualizados: número de casos activos, plazos por vencer en 48 h/7 días, contratos pendientes de firma. \\
    \textbf{Actores involucrados} & Owner, Litigante; sistema de consulta de datos agregados. \\
    \textbf{Criterio de aceptación} & Dado un Owner que accede al dashboard, cuando existen plazos fatales por vencer en menos de 48 horas, entonces el sistema los destaca visualmente con indicador de alerta roja diferenciado del resto de elementos del tablero. \\
    \bottomrule
  \end{tabularx}
\end{table}

\begin{table}[H]
  \centering
  \caption{RF-19: Generación de reportes}
  \label{tab:rf-19}
  \begin{tabularx}{\linewidth}{@{}L{3.2cm}X@{}}
    \toprule
    \textbf{ID} & RF-19 \\
    \midrule
    \textbf{Nombre} & Generación de reportes \\
    \textbf{Descripción} & El sistema debe permitir generar informes sobre el estado de los casos, desempeño de abogados y métricas operativas del despacho. \\
    \textbf{Entrada} & Parámetros de filtrado seleccionados por el usuario (rango de fechas, abogado, estado del caso, tipo de reporte). \\
    \textbf{Salida} & Reporte generado en formato PDF o exportable a CSV con los datos solicitados. \\
    \textbf{Actores involucrados} & Owner, Litigante; sistema de generación de reportes asíncrono (cola de trabajo). \\
    \textbf{Criterio de aceptación} & Dado un Owner que solicita un reporte mensual, cuando el procesamiento supera 3 segundos, entonces la interfaz permanece responsiva, muestra indicador de progreso y entrega el reporte mediante notificación sin bloquear la sesión. \\
    \bottomrule
  \end{tabularx}
\end{table}

\begin{table}[H]
  \centering
  \caption{RF-20: Bitácora de actividad (Audit Trail)}
  \label{tab:rf-20}
  \begin{tabularx}{\linewidth}{@{}L{3.2cm}X@{}}
    \toprule
    \textbf{ID} & RF-20 \\
    \midrule
    \textbf{Nombre} & Bitácora de actividad (Audit Trail) \\
    \textbf{Descripción} & El sistema debe mantener un registro cronológico inmutable de todas las acciones realizadas, incluyendo el identificador del usuario, la operación realizada, la marca de tiempo y el recurso afectado. \\
    \textbf{Entrada} & Cualquier operación realizada por un usuario autenticado (consulta, carga, firma, modificación de rol, etc.). \\
    \textbf{Salida} & Registro append-only en base de datos NoSQL con los campos: \texttt{user\_id}, \texttt{action\_type}, \texttt{resource\_id}, \texttt{timestamp}, \texttt{ip\_address}. \\
    \textbf{Actores involucrados} & Todos los roles del sistema; motor NoSQL de auditoría. \\
    \textbf{Criterio de aceptación} & Dado cualquier usuario autenticado, cuando ejecuta cualquier operación en el sistema, entonces se genera un registro en la bitácora en menos de 500 ms; el registro no puede eliminarse ni modificarse por ningún rol, incluido Owner. \\
    \bottomrule
  \end{tabularx}
\end{table}

\subsection{Requerimientos No Funcionales}\label{sec:requerimientos-no-funcionales}

Los requerimientos no funcionales describen las restricciones y atributos de calidad que el sistema debe satisfacer de forma transversal a todas sus funcionalidades. Se formulan en términos del atributo de calidad, no de la tecnología específica que lo implementa, y se agrupan en cuatro categorías: seguridad y protección de datos, rendimiento y operación, usabilidad y experiencia de usuario, y cumplimiento normativo y mantenibilidad.

\subsubsection{Seguridad y Protección de Datos}

\begin{table}[H]
  \centering
  \caption{RNF-01: Aislamiento lógico de la información del despacho}
  \label{tab:rnf-01}
  \begin{tabularx}{\linewidth}{@{}L{3.2cm}X@{}}
    \toprule
    \textbf{ID} & RNF-01 \\
    \midrule
    \textbf{Nombre} & Aislamiento lógico de la información del despacho \\
    \textbf{Descripción} & El sistema debe garantizar la segregación estricta de la información del despacho: cada instancia (single-tenant) atiende a un solo despacho con su propia base de datos y material criptográfico, y dentro de la instancia ninguna consulta, operación o error de lógica debe exponer información a usuarios sin autorización sobre el recurso. \\
    \textbf{Entrada} & Consulta o operación de cualquier usuario autenticado sobre cualquier recurso del sistema. \\
    \textbf{Salida} & Respuesta que contiene únicamente datos que el rol y las asignaciones del usuario autenticado le permiten consultar. \\
    \textbf{Actores involucrados} & Todos los usuarios; capa de middleware de autorización; base de datos relacional. \\
    \textbf{Criterio de aceptación} & Dado un usuario autenticado, cuando realiza cualquier consulta (incluso manipulando parámetros e identificadores de recursos), entonces el sistema nunca devuelve datos fuera de su ámbito de autorización; verificable mediante pruebas automatizadas que simulen manipulación de identificadores de recursos ajenos. \\
    \bottomrule
  \end{tabularx}
\end{table}

\begin{table}[H]
  \centering
  \caption{RNF-02: Gestión segura de credenciales}
  \label{tab:rnf-02}
  \begin{tabularx}{\linewidth}{@{}L{3.2cm}X@{}}
    \toprule
    \textbf{ID} & RNF-02 \\
    \midrule
    \textbf{Nombre} & Gestión segura de credenciales \\
    \textbf{Descripción} & Las contraseñas de los usuarios deben almacenarse mediante funciones de hashing resistentes (Argon2 o bcrypt) con salt (sal) aleatorio por contraseña. El sistema no debe almacenar contraseñas en texto plano ni usar SHA-256/MD5 para este fin. \\
    \textbf{Entrada} & Contraseña en texto plano proporcionada durante el registro o cambio de contraseña. \\
    \textbf{Salida} & Hash de la contraseña almacenado con Argon2 o bcrypt con salt aleatorio; contraseña original descartada inmediatamente. \\
    \textbf{Actores involucrados} & Sistema de autenticación; base de datos de credenciales. \\
    \textbf{Criterio de aceptación} & Dado el acceso directo a la base de datos, cuando se inspeccionan las credenciales almacenadas, entonces ninguna contraseña aparece en texto plano ni como hash SHA-256 o MD5; el tiempo de verificación debe superar los 100 ms para dificultar ataques de fuerza bruta. \\
    \bottomrule
  \end{tabularx}
\end{table}

\begin{table}[H]
  \centering
  \caption{RNF-03: Gestión de material criptográfico}
  \label{tab:rnf-03}
  \begin{tabularx}{\linewidth}{@{}L{3.2cm}X@{}}
    \toprule
    \textbf{ID} & RNF-03 \\
    \midrule
    \textbf{Nombre} & Gestión de material criptográfico \\
    \textbf{Descripción} & El sistema debe garantizar que las claves de cifrado se almacenen de forma segura y separada de los datos cifrados, que los IV para AES-GCM sean únicos por operación (véase §3.11.4) y que las claves privadas nunca se transmitan por la red ni se almacenen en texto plano. \\
    \textbf{Entrada} & Operación de cifrado o descifrado; solicitud de uso de clave privada por el firmante. \\
    \textbf{Salida} & Operación ejecutada sin exposición de material criptográfico sensible en logs, respuestas API o almacenamiento temporal. \\
    \textbf{Actores involucrados} & Sistema de cifrado (backend); módulo de gestión de claves; desarrolladores (auditores). \\
    \textbf{Criterio de aceptación} & Dado el análisis de los logs del sistema en cualquier nivel de verbosidad, cuando se buscan patrones de claves o IV en texto plano, entonces no debe encontrarse ninguna clave maestra, clave privada ni IV almacenado junto al contenido cifrado en el mismo registro. \\
    \bottomrule
  \end{tabularx}
\end{table}

\begin{table}[H]
  \centering
  \caption{RNF-04: Inmutabilidad de registros críticos}
  \label{tab:rnf-04}
  \begin{tabularx}{\linewidth}{@{}L{3.2cm}X@{}}
    \toprule
    \textbf{ID} & RNF-04 \\
    \midrule
    \textbf{Nombre} & Inmutabilidad de registros críticos \\
    \textbf{Descripción} & Los registros de historial procesal, cadena de custodia y bitácora de actividad deben seguir un modelo de solo inserción (append-only), impidiendo la modificación o eliminación de trazas históricas. \\
    \textbf{Entrada} & Intento de operación UPDATE o DELETE sobre tablas de historial procesal, cadena de custodia o bitácora de actividad. \\
    \textbf{Salida} & Operación rechazada con código de error; registro del intento en la bitácora de sistema. \\
    \textbf{Actores involucrados} & Cualquier usuario o proceso del sistema; motor de base de datos NoSQL. \\
    \textbf{Criterio de aceptación} & Dado un registro existente en la bitácora de auditoría, cuando cualquier proceso —incluidos scripts administrativos— intenta modificarlo o eliminarlo, entonces el motor rechaza la operación; verificable mediante pruebas de integración automatizadas. \\
    \bottomrule
  \end{tabularx}
\end{table}

\begin{table}[H]
  \centering
  \caption{RNF-05: Privacidad de la información del despacho}
  \label{tab:rnf-05}
  \begin{tabularx}{\linewidth}{@{}L{3.2cm}X@{}}
    \toprule
    \textbf{ID} & RNF-05 \\
    \midrule
    \textbf{Nombre} & Privacidad de la información del despacho \\
    \textbf{Descripción} & El sistema debe prohibir cualquier fuga de información hacia usuarios sin autorización o hacia el exterior de la instancia del despacho, verificable mediante pruebas automatizadas de aislamiento. \\
    \textbf{Entrada} & Consultas con parámetros manipulados que incluyan identificadores de recursos fuera del ámbito del usuario. \\
    \textbf{Salida} & Respuesta vacía o error de autorización; nunca datos fuera del ámbito autorizado. \\
    \textbf{Actores involucrados} & Middleware de autorización; capa de acceso a datos; equipo de QA. \\
    \textbf{Criterio de aceptación} & Dado un conjunto de pruebas automatizadas de aislamiento ejecutadas contra el entorno de staging, entonces el 100 \% de los intentos de acceso a datos fuera del ámbito de autorización deben resultar en error de autorización. \\
    \bottomrule
  \end{tabularx}
\end{table}

\subsubsection{Rendimiento y Operación}

\begin{table}[H]
  \centering
  \caption{RNF-06: Procesamiento asíncrono}
  \label{tab:rnf-06}
  \begin{tabularx}{\linewidth}{@{}L{3.2cm}X@{}}
    \toprule
    \textbf{ID} & RNF-06 \\
    \midrule
    \textbf{Nombre} & Procesamiento asíncrono \\
    \textbf{Descripción} & Las tareas intensivas (envío de notificaciones, generación de reportes, operaciones criptográficas de larga duración) deben procesarse en segundo plano mediante colas de trabajo, sin bloquear la interfaz de usuario. \\
    \textbf{Entrada} & Solicitud de tarea intensiva por parte del usuario. \\
    \textbf{Salida} & Confirmación inmediata de recepción; entrega del resultado mediante notificación cuando el procesamiento concluye. \\
    \textbf{Actores involucrados} & Usuario solicitante; sistema de colas de trabajo (backend); módulo de notificaciones. \\
    \textbf{Criterio de aceptación} & Dado un usuario que solicita la generación de un reporte, cuando el procesamiento supera 2 segundos, entonces la interfaz permanece responsiva con indicador de progreso; el usuario no necesita mantener la pantalla activa para recibir el resultado. \\
    \bottomrule
  \end{tabularx}
\end{table}

\begin{table}[H]
  \centering
  \caption{RNF-07: Tiempo de respuesta de interfaz}
  \label{tab:rnf-07}
  \begin{tabularx}{\linewidth}{@{}L{3.2cm}X@{}}
    \toprule
    \textbf{ID} & RNF-07 \\
    \midrule
    \textbf{Nombre} & Tiempo de respuesta de interfaz \\
    \textbf{Descripción} & Las operaciones interactivas del sistema deben ofrecer tiempos de respuesta perceptiblemente inmediatos, con actualizaciones reactivas que no requieran recargas completas de página. \\
    \textbf{Entrada} & Interacción del usuario con cualquier elemento de la interfaz (clic, búsqueda, navegación). \\
    \textbf{Salida} & Actualización reactiva de la interfaz sin recarga completa de página. \\
    \textbf{Actores involucrados} & Usuario final; frontend del sistema; API REST del backend. \\
    \textbf{Criterio de aceptación} & Dado un usuario en condiciones de red estándar (>10 Mbps), cuando realiza cualquier operación interactiva que no involucre procesamiento asíncrono, entonces la interfaz responde con retroalimentación visual en menos de 300 ms. \\
    \bottomrule
  \end{tabularx}
\end{table}

\subsubsection{Usabilidad y Experiencia de Usuario}

\begin{table}[H]
  \centering
  \caption{RNF-08: Diseño responsivo}
  \label{tab:rnf-08}
  \begin{tabularx}{\linewidth}{@{}L{3.2cm}X@{}}
    \toprule
    \textbf{ID} & RNF-08 \\
    \midrule
    \textbf{Nombre} & Diseño responsivo \\
    \textbf{Descripción} & La interfaz debe adaptarse fluidamente a dispositivos de escritorio y tabletas, respetando los lineamientos de legibilidad y jerarquía visual del sistema de diseño. \\
    \textbf{Entrada} & Acceso al sistema desde dispositivos con resoluciones entre 768 px (tablet) y 1920 px (escritorio). \\
    \textbf{Salida} & Interfaz adaptada correctamente al ancho de pantalla del dispositivo sin desbordamientos ni pérdida de funcionalidad. \\
    \textbf{Actores involucrados} & Litigante, Paralegal, Owner; sistema de diseño (componentes UI). \\
    \textbf{Criterio de aceptación} & Dado un dispositivo tablet con resolución de 768 px de ancho, cuando el usuario accede a cualquier módulo, entonces todos los elementos interactivos son accesibles y legibles sin requerir desplazamiento horizontal. \\
    \bottomrule
  \end{tabularx}
\end{table}

\begin{table}[H]
  \centering
  \caption{RNF-09: Soporte de localización}
  \label{tab:rnf-09}
  \begin{tabularx}{\linewidth}{@{}L{3.2cm}X@{}}
    \toprule
    \textbf{ID} & RNF-09 \\
    \midrule
    \textbf{Nombre} & Soporte de localización \\
    \textbf{Descripción} & La plataforma debe manejar formatos de fecha, moneda y textos en español de México (es-MX) como configuración predeterminada. \\
    \textbf{Entrada} & Configuración regional del sistema establecida como es-MX. \\
    \textbf{Salida} & Fechas en formato DD/MM/AAAA, montos en MXN con separador de miles y dos decimales, textos de interfaz en español de México. \\
    \textbf{Actores involucrados} & Sistema de internacionalización (i18n); todos los usuarios. \\
    \textbf{Criterio de aceptación} & Dado cualquier módulo del sistema, cuando se muestran fechas o montos, entonces el formato corresponde a es-MX sin que el usuario deba configurar manualmente su preferencia regional. \\
    \bottomrule
  \end{tabularx}
\end{table}

\subsubsection{Cumplimiento Normativo y Mantenibilidad}

\begin{table}[H]
  \centering
  \caption{RNF-10: Conformidad con el CNPP}
  \label{tab:rnf-10}
  \begin{tabularx}{\linewidth}{@{}L{3.2cm}X@{}}
    \toprule
    \textbf{ID} & RNF-10 \\
    \midrule
    \textbf{Nombre} & Conformidad con el CNPP \\
    \textbf{Descripción} & El modelado de datos y los flujos de trabajo deben adherirse estrictamente a la terminología y plazos del Código Nacional de Procedimientos Penales. \\
    \textbf{Entrada} & Flujo de trabajo relacionado con etapas procesales, plazos o tipos de audiencia. \\
    \textbf{Salida} & Terminología, plazos y secuencias de estados alineados al CNPP vigente. \\
    \textbf{Actores involucrados} & Litigante; equipo legal del despacho (validador); desarrolladores. \\
    \textbf{Criterio de aceptación} & Dado el catálogo de tipos de audiencia implementado, cuando se compara contra los artículos 307-403 del CNPP, entonces cada tipo debe tener correspondencia exacta con el artículo aplicable y activar los plazos correctos. \\
    \bottomrule
  \end{tabularx}
\end{table}

\begin{table}[H]
  \centering
  \caption{RNF-11: Conformidad con la NOM-151}
  \label{tab:rnf-11}
  \begin{tabularx}{\linewidth}{@{}L{3.2cm}X@{}}
    \toprule
    \textbf{ID} & RNF-11 \\
    \midrule
    \textbf{Nombre} & Conformidad con la NOM-151 \\
    \textbf{Descripción} & Los mecanismos de sellado de tiempo y conservación de mensajes de datos deben cumplir con los requisitos de la NOM-151-SCFI-2016 y el protocolo RFC 3161. \\
    \textbf{Entrada} & Solicitud de sellado de tiempo de un contrato firmado. \\
    \textbf{Salida} & Token RFC 3161 válido y verificable de forma independiente; la constancia de un PSC autorizado es una condición de despliegue regulado, no de la TSA local de demostración. \\
    \textbf{Actores involucrados} & Sistema backend; TSA configurada; equipo legal. \\
    \textbf{Criterio de aceptación} & Dado un token emitido por el sistema, cuando se verifica con OpenSSL usando el ancla de confianza correspondiente, la verificación resulta exitosa. Si el emisor es local, el resultado se reporta como evidencia técnica sin atribuir una fuente externa sincronizada con CENAM. \\
    \bottomrule
  \end{tabularx}
\end{table}

\begin{table}[H]
  \centering
  \caption{RNF-12: Cumplimiento de la LFPDPPP}
  \label{tab:rnf-12}
  \begin{tabularx}{\linewidth}{@{}L{3.2cm}X@{}}
    \toprule
    \textbf{ID} & RNF-12 \\
    \midrule
    \textbf{Nombre} & Cumplimiento de la LFPDPPP \\
    \textbf{Descripción} & El tratamiento de datos personales debe cumplir con las obligaciones de la LFPDPPP, particularmente las medidas de seguridad exigidas por su artículo 19. \\
    \textbf{Entrada} & Dato personal de un cliente del despacho ingresado al sistema. \\
    \textbf{Salida} & Dato almacenado cifrado, accesible únicamente por roles autorizados, con registro de acceso en la bitácora. \\
    \textbf{Actores involucrados} & Todos los usuarios; sistema de cifrado; Oficial de Privacidad del despacho. \\
    \textbf{Criterio de aceptación} & Dado el conjunto de medidas técnicas implementadas, cuando se evalúan contra el artículo 19 de la LFPDPPP, entonces el sistema debe evidenciar: cifrado AES-256-GCM (confidencialidad), hash SHA-256 (integridad), RBAC (control de acceso) y bitácora append-only (trazabilidad). \\
    \bottomrule
  \end{tabularx}
\end{table}

\begin{table}[H]
  \centering
  \caption{RNF-13: Arquitectura mantenible}
  \label{tab:rnf-13}
  \begin{tabularx}{\linewidth}{@{}L{3.2cm}X@{}}
    \toprule
    \textbf{ID} & RNF-13 \\
    \midrule
    \textbf{Nombre} & Arquitectura mantenible \\
    \textbf{Descripción} & El código debe seguir principios de diseño que faciliten la escalabilidad y el mantenimiento futuro, incluyendo separación de responsabilidades, bajo acoplamiento y alta cohesión. \\
    \textbf{Entrada} & Solicitud de cambio o nueva funcionalidad que afecte un módulo existente. \\
    \textbf{Salida} & Módulo modificado sin introducir regresiones; pruebas existentes continúan pasando. \\
    \textbf{Actores involucrados} & Equipo de desarrollo; sistema de integración continua. \\
    \textbf{Criterio de aceptación} & Dado el código fuente del sistema, cuando se analiza con métricas de código, entonces la complejidad ciclomática promedio por función no debe superar 10 y la cobertura de pruebas de los módulos criptográficos debe ser igual o mayor al 90 \%. \\
    \bottomrule
  \end{tabularx}
\end{table}

\begin{table}[H]
  \centering
  \caption{RNF-14: Aseguramiento de calidad}
  \label{tab:rnf-14}
  \begin{tabularx}{\linewidth}{@{}L{3.2cm}X@{}}
    \toprule
    \textbf{ID} & RNF-14 \\
    \midrule
    \textbf{Nombre} & Aseguramiento de calidad \\
    \textbf{Descripción} & El sistema debe someterse a pruebas de aceptación de usuario (UAT) al finalizar cada ciclo de desarrollo, complementadas con pruebas funcionales automatizadas y pruebas específicas de los servicios criptográficos. \\
    \textbf{Entrada} & Entregable de cada iteración XP; matriz de pruebas definida en la Fase I. \\
    \textbf{Salida} & Reporte de resultados de pruebas funcionales, de integración y criptográficas; lista de defectos identificados y su estado de resolución. \\
    \textbf{Actores involucrados} & Equipo de desarrollo; usuarios representativos del despacho (UAT); directora del proyecto. \\
    \textbf{Criterio de aceptación} & Dado el cierre de cada iteración, cuando se ejecuta la suite de pruebas automatizadas, entonces el 100 \% de las pruebas de módulos criptográficos (hash, cifrado, firma, sello de tiempo) deben pasar sin errores antes de considerar la iteración completada. \\
    \bottomrule
  \end{tabularx}
\end{table}


\section{Análisis de Costos}\label{sec:costos}

El presente análisis estima el costo total del proyecto en su fase de prototipo, considerando todos los recursos necesarios para el desarrollo, la infraestructura y los servicios externos requeridos. La estimación se basa en el costo de oportunidad del equipo de desarrollo y en la identificación exhaustiva de cada rubro de gasto, independientemente de que su costo sea cero en la fase actual.

\subsection{Metodología de Estimación}\label{sec:costos-metodologia}

Se utiliza un enfoque de costo de oportunidad para valorar el recurso humano: el salario de mercado que cada miembro del equipo percibiría como Desarrollador Full-Stack Junior/Mid en la Zona Metropolitana del Valle de México, convertido a tarifa horaria. Los rangos salariales se triangularon a partir de seis fuentes de datos actualizadas al periodo 2025--2026.

\begin{table}[H]
  \centering
  \caption{Fuentes salariales consultadas para la estimación del costo de recurso humano.}
  \label{tab:fuentes-salariales}
  \begin{tabularx}{\linewidth}{@{}L{0.42\linewidth}L{0.36\linewidth}R{0.16\linewidth}@{}}
    \toprule
    \textbf{Fuente} & \textbf{Rango Mensual (MXN)} & \textbf{Fecha de Datos} \\
    \midrule
    Glassdoor MX (Jr Full-Stack) & \$11,854 -- \$20,114 & Marzo 2026 \\
    Jobted MX (Jr Full-Stack) & \$11,650 promedio Jr & 2026 \\
    Indeed MX (Programador Jr) & \$12,784 promedio & Febrero 2026 \\
    Coderhouse (Jr Full-Stack) & \$22,000 -- \$28,000 & 2025 \\
    Talently (Jr Full-Stack MX) & ${\approx}$\$29,000 (${\approx}$\$1,436 USD) & Mayo 2024 \\
    Data México (ENOE Q1-2025) & \$8,410 promedio nacional* & Q1 2025 \\
    \bottomrule
  \end{tabularx}
  \par\smallskip
  {\footnotesize * Data México reporta el promedio nacional incluyendo todas las seniorities y regiones, lo cual sesga el dato hacia abajo. Se utiliza como límite inferior de referencia.}
\end{table}

A partir de estas fuentes, se definen tres escenarios salariales. El cálculo asume una jornada mensual estándar de 160 horas (40 hrs/semana $\times$ 4 semanas) bajo esquema de honorarios/freelance, lo cual incorpora el sobrecosto por ausencia de prestaciones de ley ($\approx$30--35\,\% sobre nómina tradicional).

\begin{table}[H]
  \centering
  \caption{Escenarios salariales y tarifa horaria derivada.}
  \label{tab:escenarios-salariales}
  \begin{tabularx}{\linewidth}{@{}L{0.13\linewidth}L{0.20\linewidth}L{0.13\linewidth}X@{}}
    \toprule
    \textbf{Escenario} & \textbf{Salario Mensual Equivalente} & \textbf{Tarifa / Hora} & \textbf{Justificación} \\
    \midrule
    Optimista & \$18,000 MXN & \$112.50 MXN & Piso del rango Jr según Indeed/Glassdoor; escenario donde el riesgo tiene bajo costo. \\
    Más Probable & \$25,000 MXN & \$156.25 MXN & Mediana del rango Jr-Mid; consistente con Coderhouse y Glassdoor p50. \\
    Pesimista & \$32,000 MXN & \$200.00 MXN & Techo del rango Mid; incluye prima por conocimiento criptográfico especializado. \\
    \bottomrule
  \end{tabularx}
\end{table}

Valor adoptado para los cálculos del proyecto: \textbf{\$156.25 MXN/hora} (escenario Más Probable). Los escenarios Optimista y Pesimista se reservan para el análisis de sensibilidad en la sección de riesgos.

\subsection{Capacidad del Equipo}\label{sec:costos-capacidad}

La capacidad efectiva de programación se reduce respecto a una jornada laboral típica por factores inherentes al contexto del proyecto. Los integrantes del equipo son estudiantes de último semestre con carga académica concurrente, y la metodología XP impone overhead por prácticas como Pair Programming, TDD y ceremonias iterativas.

\begin{table}[H]
  \centering
  \caption{Factores de ajuste sobre la capacidad de desarrollo.}
  \label{tab:factores-capacidad}
  \begin{tabularx}{\linewidth}{@{}L{0.28\linewidth}L{0.18\linewidth}X@{}}
    \toprule
    \textbf{Factor} & \textbf{Reducción Estimada} & \textbf{Razón} \\
    \midrule
    Carga académica concurrente & $-$30\,\% & Otras materias, exámenes y prácticas reducen la disponibilidad a 30 hrs/semana por desarrollador. \\
    Pair Programming (XP) & Factor $\times$2 en costo & Ambos desarrolladores trabajan simultáneamente en la misma tarea; se contabiliza el esfuerzo de ambos. \\
    TDD (Test-Driven Development) & $-$15\,\% & Overhead por escritura de pruebas antes de código funcional; compensado por menor deuda técnica. \\
    Ceremonias XP (standups, planning, retros) & $-$5\,\% & Iteraciones cortas (1--2 semanas) requieren reuniones frecuentes de coordinación. \\
    \bottomrule
  \end{tabularx}
\end{table}

Partiendo de una disponibilidad base de 30 horas semanales por desarrollador (jornada reducida por carga académica), y aplicando el overhead combinado de TDD y ceremonias XP ($-$20\,\%), se obtienen 24 horas efectivas de programación por desarrollador por semana. Con un equipo de 2 desarrolladores y 5 semanas de codificación intensiva (Fase II del plan de iteraciones), la capacidad total es de 240 horas-persona.

\begin{table}[H]
  \centering
  \caption{Resumen de capacidad del equipo y presupuesto base.}
  \label{tab:capacidad-resumen}
  \begin{tabularx}{\linewidth}{@{}X r@{}}
    \toprule
    \textbf{Métrica} & \textbf{Valor} \\
    \midrule
    Horas efectivas por desarrollador / semana & 24 horas \\
    Horas efectivas del equipo / semana & 48 horas-persona \\
    Semanas de codificación (Fase II) & 5 semanas \\
    Capacidad total & 240 horas-persona \\
    Presupuesto base (recurso humano) & \$37,500.00 MXN \\
    \bottomrule
  \end{tabularx}
\end{table}

\textbf{Nota:} En sesiones de Pair Programming, ambos desarrolladores trabajan en la misma tarea. Esto se refleja contabilizando las horas de ambos (48 hrs-equipo/semana), pero la producción de funcionalidades corresponde a una línea de trabajo, no dos paralelas.

\subsection{Desglose de Costos por Rubro}\label{sec:costos-desglose}

La siguiente tabla presenta el desglose completo de costos del proyecto en su fase de prototipo. Los rubros con costo cero se documentan explícitamente para transparentar que fueron evaluados y que su gratuidad corresponde a decisiones técnicas deliberadas, no a omisiones en la estimación.

\begin{table}[H]
  \centering
  \caption{Desglose de costos del proyecto en fase de prototipo.}
  \label{tab:desglose-costos}
  \begin{tabularx}{\linewidth}{@{}L{0.22\linewidth}L{0.20\linewidth}L{0.14\linewidth}X@{}}
    \toprule
    \textbf{Rubro} & \textbf{Proveedor / Herramienta} & \textbf{Costo en Fase Prototipo} & \textbf{Justificación} \\
    \midrule
    Recurso humano (2 desarrolladores) & Costo de oportunidad & \$37,500.00 MXN & 240 hrs-persona $\times$ \$156.25/hr. Escenario Más Probable. \\
    Sellado de tiempo externo (PSC) & Cincel PSC -- Sandbox & No contratado & La API no ofreció garantía operativa suficiente y el sandbox requiere pago con precios no transparentes. La evidencia de esta entrega utiliza la TSA local; no se presupone un costo externo. \\
    Infraestructura de servidor & VPS preexistente (aportación del equipo) & \$0.00 MXN & Infraestructura personal aportada por el equipo. No genera costo incremental para el proyecto. \\
    Correo transaccional & Resend -- Plan gratuito & \$0.00 MXN & Plan gratuito suficiente para el volumen de notificaciones del prototipo. \\
    Framework y herramientas de desarrollo & Rust, Axum, Astro, Svelte, Git (open source) & \$0.00 MXN & Stack completamente open source. Sin licencias comerciales requeridas. \\
    Base de datos & MySQL + almacenamiento NoSQL (open source) & \$0.00 MXN & Motores de base de datos open source desplegados en el VPS del equipo. \\
    Certificados SSL/TLS & Let's Encrypt & \$0.00 MXN & Autoridad de certificación gratuita y automatizada. \\
    Reserva de contingencia & Derivada del análisis de riesgos (sección \ref{sec:riesgos}) & \$45,578.13 MXN & VME total de línea base. Representa el costo esperado ponderado antes de actualizar el estado de materialización. \\
    \bottomrule
  \end{tabularx}
\end{table}

\subsection{Resumen Financiero}\label{sec:costos-resumen}

\begin{table}[H]
  \centering
  \caption{Resumen financiero del proyecto.}
  \label{tab:resumen-financiero}
  \begin{tabularx}{\linewidth}{@{}X r@{}}
    \toprule
    \textbf{Concepto} & \textbf{Monto} \\
    \midrule
    Costo base del proyecto (recurso humano) & \$37,500.00 MXN \\
    Costos conocidos de infraestructura y servicios & \$0.00 MXN \\
    Sandbox del PSC & No incorporado \\
    Costo total conocido sin contingencia & \$37,500.00 MXN \\
    Reserva de contingencia (VME total de línea base) & \$45,578.13 MXN \\
    Presupuesto conocido, sin sandbox del PSC & \$83,078.13 MXN \\
    \bottomrule
  \end{tabularx}
\end{table}

La reserva de contingencia de línea base representa el 121.5\,\% del costo base, lo cual refleja el alto riesgo técnico inherente a la integración de servicios criptográficos con tolerancia cero a fallos en un cronograma de 5 semanas de codificación. Este porcentaje es consistente con las recomendaciones del PMI para proyectos con incertidumbre técnica elevada. El presupuesto es una cota inferior: el costo del sandbox del PSC permanece separado hasta definir el tamaño de la campaña y contar con una tarifa verificable.

Es importante señalar que los costos conocidos de infraestructura y servicios son cero exclusivamente en la fase de prototipo y no incluyen el sandbox de pago del PSC. Un despliegue en producción requeriría contratar un servicio de sellado con volumen y estabilidad verificables, infraestructura cloud dedicada con respaldos, y licencias de dominio y certificados de producción, rubros que quedan fuera del alcance de la presente estimación.

\section{Análisis de Riesgos}\label{sec:riesgos}

El presente análisis identifica, evalúa y prioriza los riesgos que pueden afectar el desarrollo del prototipo, considerando tanto su probabilidad de ocurrencia como su impacto en tiempo y costo. La evaluación combina un enfoque cualitativo (matriz de severidad) con uno cuantitativo (Valor Monetario Esperado) para proporcionar una base objetiva en la toma de decisiones sobre mitigación y asignación de recursos.

\subsection{Modelo de Evaluación}\label{sec:riesgos-modelo}

El impacto financiero de cada riesgo materializado se calcula mediante la fórmula: IF = H(retraso) $\times$ N(devs) $\times$ T(hora), donde H(retraso) es el número de horas de retraso estimadas, N(devs) es el número de desarrolladores afectados (1 o 2) y T(hora) es la tarifa horaria del escenario Más Probable (\$156.25 MXN). El Valor Monetario Esperado (VME) pondera el impacto financiero por la probabilidad de ocurrencia: VME = P $\times$ IF.

\begin{table}[H]
  \centering
  \caption{Escala de probabilidad de ocurrencia.}
  \label{tab:escala-probabilidad}
  \begin{tabularx}{\linewidth}{@{}L{0.10\linewidth}L{0.08\linewidth}L{0.14\linewidth}X@{}}
    \toprule
    \textbf{Nivel} & \textbf{Valor} & \textbf{Rango} & \textbf{Criterio Orientativo} \\
    \midrule
    Bajo & 1 & $<$ 20\,\% & El riesgo es poco probable; el equipo tiene experiencia previa en el área o existen mitigaciones naturales. \\
    Medio & 2 & 20 -- 60\,\% & El riesgo tiene posibilidades razonables de materializarse; depende de factores parcialmente controlables. \\
    Alto & 3 & $>$ 60\,\% & El riesgo es altamente probable; involucra tecnología no dominada, dependencias externas críticas o restricciones de tiempo severas. \\
    \bottomrule
  \end{tabularx}
\end{table}

\begin{table}[H]
  \centering
  \caption{Escala de impacto en tiempo.}
  \label{tab:escala-impacto}
  \begin{tabularx}{\linewidth}{@{}L{0.10\linewidth}L{0.08\linewidth}L{0.20\linewidth}X@{}}
    \toprule
    \textbf{Nivel} & \textbf{Valor} & \textbf{Retraso} & \textbf{Criterio Orientativo} \\
    \midrule
    Bajo & 1 & 1 -- 2 días ($\leq$ 16 hrs) & Impacto menor; se puede absorber dentro de la iteración XP sin reprogramar. \\
    Medio & 2 & 3 -- 5 días (17--40 hrs) & Impacto significativo; requiere repriorizar historias de usuario y posiblemente reducir alcance de la iteración. \\
    Alto & 3 & $>$ 5 días ($>$ 40 hrs) & Impacto crítico; compromete el hito de entrega de la iteración o del proyecto completo; requiere escalamiento. \\
    \bottomrule
  \end{tabularx}
\end{table}

La multiplicación de ambas escalas produce un índice de severidad que permite priorizar la atención a riesgos:

\begin{table}[H]
  \centering
  \caption{Matriz de severidad (Probabilidad $\times$ Impacto).}
  \label{tab:matriz-severidad}
  \begin{tabular}{@{}cccc@{}}
    \toprule
    \textbf{Prob. \textbackslash{} Impacto} & \textbf{Bajo (1)} & \textbf{Medio (2)} & \textbf{Alto (3)} \\
    \midrule
    Bajo (1) & 1 -- Trivial & 2 -- Menor & 3 -- Moderado \\
    Medio (2) & 2 -- Menor & 4 -- Significativo & 6 -- Grave \\
    Alto (3) & 3 -- Moderado & 6 -- Grave & 9 -- Crítico \\
    \bottomrule
  \end{tabular}
\end{table}

Clasificación por rango de severidad: 1--2 (Verde): Aceptar y monitorear sin acción inmediata. 3--4 (Amarillo): Mitigar, definir plan de contingencia y asignar responsable. 6 (Naranja): Priorizar, implementar mitigación activa y considerar reducción de alcance. 9 (Rojo): Escalar, riesgo crítico que amenaza la viabilidad del proyecto y requiere acción inmediata.

\subsection{Riesgos Técnicos y Criptográficos}\label{sec:riesgos-tecnicos}

\begin{table}[H]
  \centering
  \caption{Ficha del riesgo RTC-01: Inestabilidad en la integración con el PSC (RFC 3161).}
  \label{tab:rtc-01}
  \begin{tabularx}{\linewidth}{@{}>{\bfseries}l X@{}}
    \toprule
    ID & RTC-01 \\
    Nombre & Inestabilidad en la integración con el PSC (RFC 3161) \\
    Cat. & Técnico \\
    Sev. & 6 -- Grave \\
    \midrule
    Declaración & \emph{Existe el riesgo de \textbf{que la integración con la API del PSC (Cincel) no sea suficientemente estable durante el flujo de sellado de tiempo}, debido a \textbf{respuestas externas no deterministas, disponibilidad fuera del control del equipo y acceso de pago al sandbox}, lo que podría resultar en \textbf{pruebas no reproducibles, costo variable de reintentos y bloqueo temporal de la emisión de constancias externas}.} \\
    Mitigación & Validar el contrato remoto mediante un Spike, aislarlo detrás del puerto \texttt{TimestampService}, cubrir sus estados con un stub HTTP y mantener una TSA RFC 3161 local para pruebas y demostraciones deterministas. \\
    Impacto & 2 -- Medio (30 h) \\
    VME & \$6,562.50 MXN \\
    Estado al corte & Materializado. El registro se completó, pero la API observada no ofreció estabilidad suficiente para una campaña repetible y el sandbox requiere pago. \\
    Contingencia & Emitir tokens RFC 3161 con una TSA local cuyo certificado deriva de la CA interna. Esta vía conserva el flujo técnico, pero se etiqueta como evidencia interna y no sustituye la confianza independiente de un PSC autorizado. \\
    \bottomrule
  \end{tabularx}
\end{table}

\begin{table}[H]
  \centering
  \caption{Ficha del riesgo RTC-02: Implementación insegura de AES-256-GCM.}
  \label{tab:rtc-02}
  \begin{tabularx}{\linewidth}{@{}>{\bfseries}l X@{}}
    \toprule
    ID & RTC-02 \\
    Nombre & Implementación insegura de AES-256-GCM \\
    Cat. & Técnico \\
    Sev. & 6 -- Grave \\
    \midrule
    Declaración & \emph{Existe el riesgo de \textbf{que el esquema de cifrado autenticado quede criptográficamente comprometido}, debido a \textbf{una generación predecible o reutilización del Vector de Inicialización (IV) en AES-256-GCM}, lo que podría resultar en \textbf{la pérdida simultánea de confidencialidad e integridad de todos los expedientes cifrados almacenados en el sistema}.} \\
    Mitigación & Aplicar Pair Programming obligatorio en este módulo. Desarrollar pruebas unitarias TDD que validen la aleatoriedad y no reutilización del IV mediante \texttt{random\_bytes}. \\
    Impacto & 3 -- Alto (45 h) \\
    VME & \$5,625.00 MXN \\
    Contingencia & Aislar y revertir el código vulnerado, purgar la base de datos de pruebas, rotar las llaves maestras de cifrado y refactorizar la lógica criptográfica. \\
    \bottomrule
  \end{tabularx}
\end{table}

\begin{table}[H]
  \centering
  \caption{Ficha del riesgo RTC-03: Mala configuración de infraestructura.}
  \label{tab:rtc-03}
  \begin{tabularx}{\linewidth}{@{}>{\bfseries}l X@{}}
    \toprule
    ID & RTC-03 \\
    Nombre & Mala configuración de infraestructura \\
    Cat. & Técnico \\
    Sev. & 2 -- Menor \\
    \midrule
    Declaración & \emph{Existe el riesgo de \textbf{que el módulo de almacenamiento de objetos quede inaccesible o con permisos incorrectos}, debido a \textbf{una configuración errónea de políticas de acceso o del servicio de almacenamiento en el servidor de desarrollo}, lo que podría resultar en \textbf{el bloqueo de las operaciones de carga y descarga de documentos, impidiendo avanzar con las iteraciones de gestión documental}.} \\
    Mitigación & Diseñar la configuración de infraestructura con políticas de mínimo privilegio desde la Iteración 4 y documentarla en el repositorio compartido. \\
    Impacto & 1 -- Bajo (15 h) \\
    VME & \$703.13 MXN \\
    Contingencia & Un desarrollador asume la resolución de infraestructura mientras el otro migra temporalmente el almacenamiento al disco local del servidor para no detener el avance. \\
    \bottomrule
  \end{tabularx}
\end{table}

\begin{table}[H]
  \centering
  \caption{Ficha del riesgo RTC-04: Pérdida de clave maestra de cifrado.}
  \label{tab:rtc-04}
  \begin{tabularx}{\linewidth}{@{}>{\bfseries}l X@{}}
    \toprule
    ID & RTC-04 \\
    Nombre & Pérdida de clave maestra de cifrado \\
    Cat. & Técnico \\
    Sev. & 3 -- Moderado \\
    \midrule
    Declaración & \emph{Existe el riesgo de \textbf{que los expedientes cifrados se vuelvan irrecuperables de forma permanente}, debido a \textbf{la ausencia de un procedimiento formal de resguardo, rotación y recuperación de la clave maestra AES del despacho}, lo que podría resultar en \textbf{la pérdida total e irreversible de la información cifrada de todos los clientes del despacho, sin posibilidad de restauración}.} \\
    Mitigación & Definir desde la Iteración 4 el procedimiento de resguardo de claves con almacenamiento separado de los datos cifrados (ej. variable de entorno o servicio de gestión de secretos). \\
    Impacto & 3 -- Alto ($>$40 h) \\
    VME & \$7,031.25 MXN \\
    Contingencia & Restaurar desde la copia de seguridad del almacén de claves; si no existe respaldo, declarar incidente crítico de pérdida de datos y notificar a los interesados. \\
    \bottomrule
  \end{tabularx}
\end{table}

\begin{table}[H]
  \centering
  \caption{Ficha del riesgo RTC-05: Expiración de certificado X.509 en producción.}
  \label{tab:rtc-05}
  \begin{tabularx}{\linewidth}{@{}>{\bfseries}l X@{}}
    \toprule
    ID & RTC-05 \\
    Nombre & Expiración de certificado X.509 en producción \\
    Cat. & Técnico \\
    Sev. & 4 -- Significativo \\
    \midrule
    Declaración & \emph{Existe el riesgo de \textbf{que el flujo de firma digital quede inoperante de forma inesperada en producción}, debido a \textbf{la falta de monitoreo automatizado de la fecha de vencimiento del certificado X.509 utilizado por el despacho}, lo que podría resultar en \textbf{la interrupción del servicio de no repudio y el potencial incumplimiento de compromisos contractuales con clientes activos del despacho}.} \\
    Mitigación & Implementar en el módulo de administración una alerta preventiva configurable (ej. 30 y 7 días antes del vencimiento) vinculada al atributo de expiración del certificado X.509. \\
    Impacto & 2 -- Medio (20 h) \\
    VME & \$3,125.00 MXN \\
    Contingencia & Emitir un certificado de reemplazo de forma urgente ante el PSC correspondiente; notificar a los firmantes afectados y reprogramar las operaciones de firma pendientes. \\
    \bottomrule
  \end{tabularx}
\end{table}

\subsection{Riesgos Legales y de Cumplimiento}\label{sec:riesgos-legales}

\begin{table}[H]
  \centering
  \caption{Ficha del riesgo RLC-01: Discrepancia de hashes SHA-256 (NOM-151).}
  \label{tab:rlc-01}
  \begin{tabularx}{\linewidth}{@{}>{\bfseries}l X@{}}
    \toprule
    ID & RLC-01 \\
    Nombre & Discrepancia de hashes SHA-256 (NOM-151) \\
    Cat. & Legal \\
    Sev. & 6 -- Grave \\
    \midrule
    Declaración & \emph{Existe el riesgo de \textbf{que las constancias de conservación emitidas por el PSC sean declaradas inválidas ante una instancia judicial}, debido a \textbf{una discrepancia entre el hash SHA-256 calculado por el sistema y el hash recibido por la TSA, originada por inconsistencias en la serialización del documento previo al sellado}, lo que podría resultar en \textbf{la pérdida del valor probatorio de todos los contratos firmados y el incumplimiento formal de la NOM-151-SCFI-2016, exponiendo al despacho a impugnaciones}.} \\
    Mitigación & Estudiar exhaustivamente el anexo técnico de la NOM-151. Auditar los digests generados en Rust cruzando resultados con herramientas certificadas (OpenSSL CLI) en cada Spike. \\
    Impacto & 3 -- Alto (48 h) \\
    VME & \$7,500.00 MXN \\
    Contingencia & Detener el flujo de firmas, validar técnicamente la estructura del hash y refactorizar la lógica de serialización antes de cualquier nuevo envío al PSC. \\
    \bottomrule
  \end{tabularx}
\end{table}

\begin{table}[H]
  \centering
  \caption{Ficha del riesgo RLC-02: Mutabilidad en registros de auditoría (CNPP).}
  \label{tab:rlc-02}
  \begin{tabularx}{\linewidth}{@{}>{\bfseries}l X@{}}
    \toprule
    ID & RLC-02 \\
    Nombre & Mutabilidad en registros de auditoría (CNPP) \\
    Cat. & Legal \\
    Sev. & 2 -- Menor \\
    \midrule
    Declaración & \emph{Existe el riesgo de \textbf{que la cadena de custodia digital pierda su valor probatorio ante un tribunal revisor}, debido a \textbf{la ausencia de controles de inmutabilidad en la capa de persistencia que permitan operaciones UPDATE o DELETE sobre los registros de auditoría}, lo que podría resultar en \textbf{la ruptura de la cadena de custodia conforme a los artículos 227 y 228 del CNPP, invalidando la evidencia digital presentada en juicio}.} \\
    Mitigación & Diseñar la capa de persistencia y el modelo NoSQL para prohibir operaciones UPDATE o DELETE en las bitácoras, forzando un esquema estrictamente append-only desde el inicio. \\
    Impacto & 2 -- Medio (24 h) \\
    VME & \$1,125.00 MXN \\
    Contingencia & Restaurar la base de datos desde respaldo previo a la mutación, auditar los registros afectados y aplicar un parche urgente a nivel de base de datos bloqueando la mutabilidad. \\
    \bottomrule
  \end{tabularx}
\end{table}

\begin{table}[H]
  \centering
  \caption{Ficha del riesgo RLC-03: Fuga de datos personales por fallas de autorización (LFPDPPP).}
  \label{tab:rlc-03}
  \begin{tabularx}{\linewidth}{@{}>{\bfseries}l X@{}}
    \toprule
    ID & RLC-03 \\
    Nombre & Fuga de datos personales por fallas de autorización (LFPDPPP) \\
    Cat. & Legal \\
    Sev. & 3 -- Moderado \\
    \midrule
    Declaración & \emph{Existe el riesgo de \textbf{que datos confidenciales del despacho sean accesibles por usuarios sin autorización}, debido a \textbf{una falla en el middleware de autorización que permita consultas sin restricción de rol o de ámbito del usuario}, lo que podría resultar en \textbf{una brecha de datos personales que viola el artículo 19 de la LFPDPPP, con exposición a sanciones administrativas y daño reputacional irreversible para el despacho}.} \\
    Mitigación & Inyectar restricción obligatoria de autorización (rol y ámbito del usuario) en toda consulta relacional y verificar el aislamiento mediante pruebas automatizadas en cada iteración. \\
    Impacto & 3 -- Alto (50 h) \\
    VME & \$1,562.50 MXN \\
    Contingencia & Activar modo de mantenimiento, purgar registros expuestos, rotar credenciales comprometidas y refactorizar los middlewares de aislamiento antes de restaurar el servicio. \\
    \bottomrule
  \end{tabularx}
\end{table}

\begin{table}[H]
  \centering
  \caption{Ficha del riesgo RLC-04: Indisponibilidad o degradación del PSC (Cincel).}
  \label{tab:rlc-04}
  \begin{tabularx}{\linewidth}{@{}>{\bfseries}l X@{}}
    \toprule
    ID & RLC-04 \\
    Nombre & Indisponibilidad o degradación del PSC (Cincel) \\
    Cat. & Legal \\
    Sev. & 4 -- Significativo \\
    \midrule
    Declaración & \emph{Existe el riesgo de \textbf{que el módulo de sellado externo quede indisponible o opere de forma degradada durante ventanas críticas de firma}, debido a \textbf{interrupciones o inestabilidad en la infraestructura del Prestador de Servicios de Certificación seleccionado}, lo que podría resultar en \textbf{retrasos en la obtención de la constancia NOM-151 y afectación al flujo operativo}.} \\
    Mitigación & Mantener una interfaz abstracta que permita configurar otro PSC sin modificar la lógica central y separar las pruebas deterministas de la disponibilidad del proveedor comercial. \\
    Impacto & 2 -- Medio (20 h) \\
    VME & \$3,125.00 MXN \\
    Estado al corte & Materializado parcialmente como degradación e inestabilidad de la API, no como una interrupción total del registro o del servicio. \\
    Contingencia & Encolar las solicitudes externas y reintentar con límites explícitos; para continuidad técnica usar la TSA interna, sin conmutación silenciosa y sin presentar sus tokens como constancias emitidas por un tercero. \\
    \bottomrule
  \end{tabularx}
\end{table}

\subsection{Riesgos Operativos y de Metodología XP}\label{sec:riesgos-operativos}

\begin{table}[H]
  \centering
  \caption{Ficha del riesgo ROP-01: Curva de aprendizaje TDD en módulos criptográficos.}
  \label{tab:rop-01}
  \begin{tabularx}{\linewidth}{@{}>{\bfseries}l X@{}}
    \toprule
    ID & ROP-01 \\
    Nombre & Curva de aprendizaje TDD en módulos criptográficos \\
    Cat. & Operativo \\
    Sev. & 4 -- Significativo \\
    \midrule
    Declaración & \emph{Existe el riesgo de \textbf{que el desarrollo de los módulos criptográficos supere el tiempo planificado de forma significativa}, debido a \textbf{la dificultad adicional de formular pruebas unitarias para operaciones criptográficas cuyo comportamiento correcto no siempre es intuitivo (verificación de IV, validación de cadenas de certificados)}, lo que podría resultar en \textbf{el consumo del margen de contingencia del cronograma de la Fase II, poniendo en riesgo la entrega del módulo de no repudio en las Iteraciones 7 y 8}.} \\
    Mitigación & Aplicar TDD estricto solo en los flujos criptográficos core (hash, cifrado AES, validación X.509) y relajarlo para CRUDs básicos que no comprometan la seguridad. \\
    Impacto & 2 -- Medio (35 h) \\
    VME & \$5,468.75 MXN \\
    Contingencia & Transicionar a pruebas de integración (test-after) en componentes secundarios si un test unitario excede 4 horas de bloqueo productivo sin resolución. \\
    \bottomrule
  \end{tabularx}
\end{table}

\begin{table}[H]
  \centering
  \caption{Ficha del riesgo ROP-02: Scope creep en interfaces gráficas.}
  \label{tab:rop-02}
  \begin{tabularx}{\linewidth}{@{}>{\bfseries}l X@{}}
    \toprule
    ID & ROP-02 \\
    Nombre & Scope creep en interfaces gráficas \\
    Cat. & Operativo \\
    Sev. & 4 -- Significativo \\
    \midrule
    Declaración & \emph{Existe el riesgo de \textbf{que el esfuerzo dedicado a interfaces gráficas consuma tiempo crítico de las semanas 9 y 10}, debido a \textbf{la tendencia a invertir iteraciones en refinamiento visual más allá de lo planificado, en detrimento de las historias de gestión procesal y documental que comparten la Iteración 9}, lo que podría resultar en \textbf{el desplazamiento de dichas historias y de las pruebas finales de la Iteración 10 hacia fuera del cronograma, comprometiendo el objetivo de gestión procesal del prototipo}.} \\
    Mitigación & Congelar las historias de usuario visuales y utilizar bibliotecas de componentes preconstruidas para la interfaz, prohibiendo el diseño a medida en la Fase II. La mitigación se refuerza estructuralmente porque el desarrollo de interfaz se traslada a la Iteración 9, posterior al hito criptográfico de la Iteración 8; el riesgo residual se concentra en el objetivo de gestión procesal (OE-2) por la compresión de esas historias en una sola iteración. \\
    Impacto & 2 -- Medio (20 h) \\
    VME & \$3,125.00 MXN \\
    Contingencia & Recortar características no críticas (RF-19 Dashboard Ejecutivo o reportes visuales) para canalizar el esfuerzo íntegramente a las historias de gestión procesal y documental de la Iteración 9. \\
    \bottomrule
  \end{tabularx}
\end{table}

\begin{table}[H]
  \centering
  \caption{Ficha del riesgo ROP-03: Ruptura de cadencia de Pair Programming.}
  \label{tab:rop-03}
  \begin{tabularx}{\linewidth}{@{}>{\bfseries}l X@{}}
    \toprule
    ID & ROP-03 \\
    Nombre & Ruptura de cadencia de Pair Programming \\
    Cat. & Operativo \\
    Sev. & 2 -- Menor \\
    \midrule
    Declaración & \emph{Existe el riesgo de \textbf{que la práctica de Pair Programming quede interrumpida por ausencia temporal de un desarrollador}, debido a \textbf{indisposición o compromisos académicos imprevistos de uno de los integrantes del equipo durante la Fase II de codificación}, lo que podría resultar en \textbf{la generación de silos de conocimiento en módulos críticos de seguridad y la pérdida de velocidad de desarrollo durante el periodo de ausencia}.} \\
    Mitigación & Mantener política de propiedad colectiva del código e integración continua diaria para garantizar que ambos desarrolladores conozcan todos los módulos. \\
    Impacto & 1 -- Bajo (16 h) \\
    VME & \$625.00 MXN \\
    Contingencia & El desarrollador activo pivota a tareas de menor riesgo técnico (maquetación, diccionarios de datos) hasta que la célula de pares se restablezca. \\
    \bottomrule
  \end{tabularx}
\end{table}

\subsection{Resumen Cuantitativo de Riesgos}\label{sec:riesgos-resumen}

\begin{table}[H]
  \centering
  \caption{Resumen cuantitativo de los riesgos identificados con su severidad y VME.}
  \label{tab:riesgos-cuantitativo}
  \small
  \begin{tabularx}{\linewidth}{@{}l X L{0.13\linewidth}L{0.13\linewidth}L{0.16\linewidth}r@{}}
    \toprule
    \textbf{ID} & \textbf{Nombre} & \textbf{Prob.} & \textbf{Impacto} & \textbf{Severidad} & \textbf{VME (MXN)} \\
    \midrule
    \textbf{RTC-01} & Inestabilidad en integración con PSC (RFC 3161) & 3 Alto (70\,\%) & 2 Medio (30 h) & \textbf{6 -- Grave} & \$6,562.50 MXN \\
    \textbf{RTC-02} & Implementación insegura de AES-256-GCM & 2 Medio (40\,\%) & 3 Alto (45 h) & \textbf{6 -- Grave} & \$5,625.00 MXN \\
    \textbf{RTC-03} & Mala configuración de infraestructura & 2 Medio (30\,\%) & 1 Bajo (15 h) & \textbf{2 -- Menor} & \$703.13 MXN \\
    \textbf{RTC-04} & Pérdida de clave maestra de cifrado & 1 Bajo (15\,\%) & 3 Alto ($>$40 h) & \textbf{3 -- Moderado} & \$7,031.25 MXN \\
    \textbf{RTC-05} & Expiración de certificado X.509 en producción & 2 Medio (40\,\%) & 2 Medio (20 h) & \textbf{4 -- Significativo} & \$3,125.00 MXN \\
    \textbf{RLC-01} & Discrepancia de hashes SHA-256 (NOM-151) & 2 Medio (50\,\%) & 3 Alto (48 h) & \textbf{6 -- Grave} & \$7,500.00 MXN \\
    \textbf{RLC-02} & Mutabilidad en registros de auditoría (CNPP) & 1 Bajo (15\,\%) & 2 Medio (24 h) & \textbf{2 -- Menor} & \$1,125.00 MXN \\
    \textbf{RLC-03} & Fuga de datos personales por fallas de autorización (LFPDPPP) & 1 Bajo (10\,\%) & 3 Alto (50 h) & \textbf{3 -- Moderado} & \$1,562.50 MXN \\
    \textbf{RLC-04} & Indisponibilidad o degradación del PSC (Cincel) & 2 Medio (30\,\%) & 2 Medio (20 h) & \textbf{4 -- Significativo} & \$3,125.00 MXN \\
    \textbf{ROP-01} & Curva de aprendizaje TDD en módulos criptográficos & 2 Medio (50\,\%) & 2 Medio (35 h) & \textbf{4 -- Significativo} & \$5,468.75 MXN \\
    \textbf{ROP-02} & Scope creep en interfaces gráficas & 2 Medio (50\,\%) & 2 Medio (20 h) & \textbf{4 -- Significativo} & \$3,125.00 MXN \\
    \textbf{ROP-03} & Ruptura de cadencia de Pair Programming & 2 Medio (25\,\%) & 1 Bajo (16 h) & \textbf{2 -- Menor} & \$625.00 MXN \\
    \bottomrule
  \end{tabularx}
\end{table}

\textbf{VME total del proyecto: \$45,578.13 MXN.} Si llegaran a materializarse simultáneamente los riesgos de mayor severidad (RTC-01, RLC-01 y ROP-01), el retraso combinado sería de 113 horas de reloj, equivalente a 2.35 semanas (113 $\div$ 48 horas-equipo/semana), lo que consumiría el 47\,\% del cronograma de codificación de 5 semanas. En términos de presupuesto, el impacto conjunto asciende a 226 horas-persona (113 $\times$ 2 desarrolladores), equivalente al 94\,\% de la capacidad total disponible.

\subsection{Actualización del riesgo del proveedor al cierre del núcleo criptográfico}\label{sec:riesgo-psc-materializado}

El registro en el entorno de Cincel pudo completarse; por tanto, el problema observado no fue una imposibilidad de alta. El riesgo se materializó durante el consumo del servicio: la API no ofreció respuesta y estabilidad suficientes para convertirla en una dependencia determinista de la suite, de la integración continua o de la demostración, y el sandbox requiere pago con precios no transparentes. La combinación introduce fallos ajenos al código propio, costo variable al repetir solicitudes y una dependencia operacional que el equipo no controla. En consecuencia, la prueba de humo se conserva como referencia técnica ignorada y la evidencia externa de Cincel queda fuera del alcance de esta entrega.

En términos de gestión, RTC-01 cambia de amenaza a incidencia materializada. Su probabilidad de 70\,\% y su VME de \$6,562.50 MXN se conservan como valores de la línea base previa para mantener trazabilidad, no como medición del costo ya incurrido. El impacto efectivo permanece sin cuantificar porque no se autoriza una campaña de pago; el riesgo residual se traslada a cualquier despliegue futuro que requiera confianza externa y deberá resolverse mediante un PSC con contrato y disponibilidad verificables.

La respuesta aplicada supera la contingencia de simulación prevista originalmente. La CA interna emite un certificado exclusivo de TSA y \texttt{LocalOpensslTsa} produce tokens RFC 3161 reales, verificables con \texttt{openssl ts -verify}, detrás del mismo puerto que utiliza el adaptador de Cincel. Esto permite reproducir el flujo criptográfico completo sin acoplar las pruebas a la disponibilidad o al costo del proveedor. La procedencia del token debe conservarse explícita: la TSA interna aporta continuidad y evidencia técnica, pero no independencia de tercero ni una constancia NOM-151 emitida por un PSC autorizado. En producción no debe existir una conmutación silenciosa entre ambos emisores; una falla externa debe encolarse o exponerse al usuario, y la selección de otra TSA regulada debe realizarse como sustitución consciente del adaptador y de su ancla de confianza.

\textbf{Recomendación:} El éxito del prototipo radica en su validez como evidencia técnica, no en su amplitud estética ni de reportes. El flujo RFC 3161 local ya fue comprobado de extremo a extremo y permite avanzar hacia la API, la persistencia y la interfaz web sin depender de Cincel. La integración con un PSC externo queda como sustitución futura y consciente del adaptador, no como condición para continuar el desarrollo.

\subsection{Trazabilidad Riesgo--Objetivo}\label{sec:riesgos-trazabilidad}

La siguiente tabla relaciona cada riesgo identificado en las secciones \ref{sec:riesgos-tecnicos} a \ref{sec:riesgos-operativos} con los objetivos específicos que se verían comprometidos si el riesgo se materializa. Este mapeo permite priorizar las mitigaciones cuya ausencia afectaría de forma más directa la entrega del prototipo.

\begin{table}[H]
  \centering
  \caption{Trazabilidad entre riesgos identificados y objetivos específicos.}
  \label{tab:riesgos-trazabilidad}
  \begin{tabularx}{\linewidth}{@{}lX L{0.30\linewidth}@{}}
    \toprule
    \textbf{ID} & \textbf{Descripción} & \textbf{Objetivo(s) afectado(s)} \\
    \midrule
    RTC-01 & Inestabilidad en integración con PSC (RFC 3161) & OE-3 \\
    RTC-02 & Implementación insegura de AES-256-GCM & OE-4 \\
    RTC-03 & Mala configuración de infraestructura & OE-1, OE-6 \\
    RTC-04 & Pérdida de clave maestra de cifrado & OE-4 \\
    RTC-05 & Expiración de certificado X.509 en producción & OE-3, OE-5 \\
    RLC-01 & Discrepancia de hashes SHA-256 (NOM-151) & OE-3, OE-4 \\
    RLC-02 & Mutabilidad en registros de auditoría (CNPP) & OE-2, OE-6 \\
    RLC-03 & Fuga de datos personales por fallas de autorización (LFPDPPP) & OE-4, OE-5 \\
    RLC-04 & Indisponibilidad o degradación del PSC (Cincel) & OE-3 \\
    ROP-01 & Curva de aprendizaje TDD en módulos criptográficos & OE-3, OE-4, OE-5 \\
    ROP-02 & Scope creep en interfaces gráficas & OE-1, OE-2 \\
    ROP-03 & Ruptura de cadencia de Pair Programming & Transversal (OE-1 a OE-6) \\
    \bottomrule
  \end{tabularx}
\end{table}


\section{Diseño de Arquitectura del Sistema}\label{sec:arquitectura}

Esta sección documenta las decisiones arquitectónicas que estructuran el prototipo y articula su correspondencia con los requerimientos funcionales y no funcionales especificados en §3.3 y §3.4. La arquitectura se construye alrededor de los atributos de calidad priorizados en §3.7.1, se materializa en el estilo hexagonal seleccionado en §3.7.3, y se enforce mediante el sistema de tipos de Rust y la organización del workspace de Cargo descrita en §3.7.6. Las vistas adicionales —de contexto, componentes, despliegue y seguridad— se derivan del modelo C4 propuesto por Brown \cite{brown2024} y del modelo 4+1 de Kruchten \cite{kruchten1995}, adaptados al alcance del prototipo. Las decisiones arquitectónicas significativas se registran en formato ADR (Architecture Decision Record) propuesto por Nygard \cite{nygard2011} en §3.7.10, y la coherencia interna del capítulo se verifica mediante la matriz de trazabilidad de §3.7.11.

\subsection{Drivers Arquitectónicos: Atributos de Calidad}\label{sec:drivers-calidad}

La arquitectura del sistema no se selecciona por adherencia a una moda tecnológica sino por su capacidad para satisfacer un conjunto explícito de atributos de calidad derivados de los requerimientos no funcionales del §3.4.2. La norma ISO/IEC 25010:2011 \cite{iso25010} clasifica los atributos de calidad de software en ocho características principales (adecuación funcional, eficiencia de desempeño, compatibilidad, usabilidad, fiabilidad, seguridad, mantenibilidad y portabilidad), de las cuales se priorizan aquellas que tienen mayor impacto sobre el dominio jurídico-criptográfico que aborda el prototipo.

El primer atributo prioritario es la seguridad, condicionada por los requerimientos RNF-01 (aislamiento lógico), RNF-02 (gestión segura de credenciales) y RNF-03 (gestión de material criptográfico). La seguridad no se concibe como una capa añadida al final del diseño sino como una propiedad transversal que la arquitectura debe garantizar desde sus fundamentos. La adopción de Rust como lenguaje del backend, la separación estricta de material criptográfico mediante el patrón envelope encryption (véase §3.11.4), y el aislamiento de los puertos criptográficos mediante traits intercambiables son decisiones arquitectónicas que materializan este atributo.

El segundo atributo prioritario es la integridad y trazabilidad, condicionada por RNF-04 (inmutabilidad de registros críticos) y RNF-11 (conformidad con la NOM-151). Los documentos gestionados por el sistema poseen valor probatorio bajo los Artículos 33 del Código de Comercio y 89-bis del mismo ordenamiento, así como bajo la NOM-151-SCFI-2016. Esto exige que la arquitectura garantice la imposibilidad de modificación retroactiva de la bitácora de auditoría —implementada mediante el encadenamiento SHA-256 descrito en §3.8.1—, así como la conservación criptográficamente verificable de los hashes y sellos de tiempo de cada documento.

El tercer atributo es la evolvabilidad, definida por Ford, Parsons y Kua como la capacidad de un sistema para cambiar guiado \cite{ford2023}. El prototipo se diseña para que la sustitución de cualquier sistema externo —el Prestador de Servicios de Certificación (Cincel por SeguriData o CECOBAN), el motor de base de datos, el proveedor de correo— requiera únicamente la implementación de un nuevo adaptador, sin modificación del dominio. Esta propiedad se garantiza mediante la regla de dependencias de la arquitectura hexagonal, formalizada en §3.7.4.

El cuarto atributo es la testabilidad, condicionada por la metodología Extreme Programming adoptada en §3.1, que prescribe Test-Driven Development como práctica central. Una arquitectura es testeable cuando permite ejercitar la lógica de dominio sin requerir infraestructura real (bases de datos, servicios de red, sistemas externos). La separación explícita del dominio respecto a sus dependencias mediante puertos permite inyectar adaptadores de prueba (mocks o stubs) durante la ejecución de la suite de pruebas, manteniendo el dominio agnóstico respecto a la infraestructura.

El quinto atributo es la mantenibilidad —entendida como bajo acoplamiento y alta cohesión—, y el sexto es la observabilidad, esencial para la operación del sistema en producción. Atributos de menor prioridad, pero considerados son la performance (en tanto el sistema deba responder a consultas interactivas en menos de 200 ms), y la portabilidad (limitada al despliegue en sistemas Linux POSIX-compliant).

La \cref{tab:atributos-calidad} formaliza la priorización adoptada y su mapeo a los requerimientos no funcionales que la justifican. Esta priorización guía todas las decisiones arquitectónicas posteriores: ante un trade-off entre dos atributos en tensión, se favorece aquel de mayor prioridad.

\begin{table}[H]
  \centering
  \caption{Priorización de atributos de calidad y su mapeo a requerimientos no funcionales.}
  \label{tab:atributos-calidad}
  \begin{tabularx}{\linewidth}{@{}C{0.18\linewidth}L{0.34\linewidth}X@{}}
    \toprule
    \textbf{Prioridad} & \textbf{Atributo de calidad} & \textbf{RNFs que lo motivan} \\
    \midrule
    1 (más alta) & Seguridad & RNF-01, RNF-02, RNF-03 \\
    2 & Integridad y trazabilidad & RNF-04, RNF-11 \\
    3 & Evolvabilidad & Sustitución de PSC, BD, SMTP \\
    4 & Testabilidad & XP/TDD (§3.1) \\
    5 & Mantenibilidad & Bajo acoplamiento, alta cohesión \\
    6 (más baja) & Observabilidad y performance & Operación productiva \\
    \bottomrule
  \end{tabularx}
\end{table}

\subsection{Vista de Contexto del Sistema}\label{sec:vista-contexto}

La vista de contexto, correspondiente al primer nivel del modelo C4 de Brown \cite{brown2024}, representa al sistema como una caja negra que interactúa con actores humanos y sistemas externos. Su propósito es establecer las fronteras del alcance del prototipo y los protocolos de integración antes de descender a los detalles internos.

La Figura 4 presenta los tres tipos de actores humanos contemplados —socio (abogado titular), paralegal (rol operativo) y cliente (firmante de contratos de prestación de servicios)— y los sistemas externos con los que el sistema se integra: el Prestador de Servicios de Certificación Cincel para sellado de tiempo conforme al RFC 3161 y la NOM-151-SCFI-2016, el proveedor de correo electrónico SMTP para notificaciones, y los navegadores web utilizados como cliente de la aplicación. Las flechas etiquetan el protocolo o estándar que regula cada interacción.

\emph{Figura 4. Vista de contexto del sistema: actores humanos y sistemas externos con sus protocolos de integración.}

El sistema queda explícitamente fuera del alcance respecto a la firma criptográfica realizada por entidades distintas a las partes contractuales (por ejemplo, no se integra con el SAT para emisión de e.firma; las razones se discuten en §3.11.2). Las integraciones con sistemas judiciales públicos (Tribunales, Fiscalías) tampoco están en alcance del prototipo, aunque la arquitectura permite añadirlas en el futuro como adaptadores driving adicionales sin modificar el dominio.

\subsection{Selección del Estilo Arquitectónico}\label{sec:estilo-arquitectonico}

La selección del estilo arquitectónico es la decisión más estructurante del diseño porque condiciona todas las decisiones subsiguientes sobre organización del código, manejo de dependencias y estrategia de pruebas. Para el presente proyecto se evaluaron seis alternativas mayores, todas documentadas en la literatura canónica de arquitectura de software: arquitectura en capas tradicional (N-tier), Clean Architecture de Martin \cite{martin2017}, Onion Architecture de Palermo, arquitectura hexagonal o de puertos y adaptadores de Cockburn \cite{cockburn2005}, microservicios \cite{newman2021} y monolito modular sin patrón hexagonal.

La arquitectura en capas tradicional fue descartada por su acoplamiento estructural del dominio a la persistencia: en este estilo, el dominio depende típicamente del repositorio (que a su vez depende del motor de base de datos), lo que impide aislar la lógica de negocio para pruebas sin levantar infraestructura real. Este acoplamiento contradice directamente los atributos de testabilidad y evolvabilidad priorizados en §3.7.1, y resulta especialmente inadecuado para un proyecto que adopta TDD como práctica central.

Los microservicios fueron descartados por desproporción operativa respecto al alcance. Newman documenta el conjunto de prerrequisitos operacionales que un sistema basado en microservicios necesita —orquestación, descubrimiento de servicios, gestión de configuración distribuida, observabilidad cross-service, pipelines de despliegue independientes por servicio, infraestructura de comunicación inter-servicios— \cite{newman2021}. Para un prototipo desarrollado por un equipo de dos personas en un periodo académico, esta complejidad es claramente desproporcionada respecto al beneficio. Fowler ha argumentado explícitamente la conveniencia de iniciar con un monolito modular y extraer microservicios solo cuando la escala lo justifique \cite{fowler2015}.

Las tres arquitecturas restantes —Clean, Onion y Hexagonal— comparten el mismo principio fundamental: la regla de dependencias hacia adentro, que establece que las capas internas (dominio) no deben conocer las capas externas (infraestructura). Difieren en su formulación y en el grado de subdivisión que prescriben. Clean Architecture introduce cuatro capas concéntricas (Entities, Use Cases, Interface Adapters, Frameworks \& Drivers); Onion Architecture introduce variantes de dos a cinco capas según el autor; Hexagonal —en su formulación original de Cockburn— establece únicamente la distinción entre el dominio (interior) y los adaptadores (exterior), conectados por puertos (interfaces de entrada y salida).

Se selecciona la arquitectura hexagonal por tres razones específicas. Primero, su formulación es la más simple: dos categorías (puertos driving y puertos driven) capturan toda la diversidad de adaptadores necesarios sin obligar a clasificaciones más finas que el alcance no amerita. Segundo, la metáfora del hexágono comunica explícitamente que los adaptadores son intercambiables: el dominio queda en el centro y alrededor pueden orbitar tantos adaptadores driving (REST, CLI, WebSocket) y driven (PostgreSQL, Redis, Cincel) como sea necesario. Tercero, la correspondencia con el sistema de tipos de Rust es directa: los puertos se materializan como traits y los adaptadores como structs con bloques impl, y el compilador garantiza mecánicamente la regla de dependencias mediante el grafo de dependencias del workspace de Cargo (véase §3.7.6).

La \cref{tab:comparativa-estilos} sintetiza la comparativa entre las seis alternativas consideradas, evaluándolas contra los atributos de calidad priorizados en §3.7.1. La decisión final se enuncia formalmente como ADR-01 en §3.7.10.

\begin{table}[H]
  \centering
  \caption{Comparativa de estilos arquitectónicos contra atributos de calidad priorizados.}
  \label{tab:comparativa-estilos}
  \begin{tabularx}{\linewidth}{@{}L{0.30\linewidth}C{0.13\linewidth}C{0.13\linewidth}X@{}}
    \toprule
    \textbf{Estilo arquitectónico} & \textbf{Testabilidad} & \textbf{Evolvabilidad} & \textbf{Adecuación al alcance} \\
    \midrule
    N-tier (capas tradicional) & Baja & Baja & Inadecuada (acoplamiento BD) \\
    Clean Architecture & Alta & Alta & Adecuada pero más prescriptiva \\
    Onion Architecture & Alta & Alta & Equivalente a Hexagonal pero menos clara \\
    Hexagonal (Cockburn) & Alta & Alta & Óptima para el alcance \\
    Microservicios & Media & Alta & Desproporcionada operativamente \\
    Monolito modular sin hexagonal & Media & Media & Insuficiente para enforce dependencias \\
    \bottomrule
  \end{tabularx}
\end{table}

\subsection{Arquitectura Hexagonal: Definición y Aplicación}\label{sec:hexagonal-definicion}

La arquitectura hexagonal, propuesta originalmente por Alistair Cockburn en 2005 \cite{cockburn2005} y refinada por Vernon \cite{vernon2013} en el contexto del Domain-Driven Design \cite{evans2003}, organiza una aplicación de software en tres capas concéntricas: el dominio en el centro, la aplicación en la capa intermedia, y los adaptadores en la capa exterior.

El dominio contiene los conceptos centrales del problema modelado: entidades (Case, Document, User, Signature, Hearing, AuditEvent, Certificate), value objects (Nuc, ProcessStage, DocumentHash) y servicios de dominio (lógica de negocio que no pertenece naturalmente a una entidad específica). La regla fundamental que gobierna esta capa es que su código no importa absolutamente nada de las capas externas: depende exclusivamente de la biblioteca estándar de Rust y de un conjunto mínimo de crates de bajo nivel (serde para serialización, thiserror para errores tipados, time para manejo de fechas, uuid para identificadores únicos, y zeroize para el borrado garantizado del material secreto en memoria al liberarlo).

La capa de aplicación contiene los casos de uso del sistema: SignContract, RegisterCase, ScheduleHearing, VerifySignature, TransitionCaseStage, QueryAuditLog, entre otros. Cada caso de uso orquesta las entidades del dominio para satisfacer una intención específica del usuario, y declara las dependencias que requiere del exterior mediante traits (los puertos). La aplicación depende exclusivamente del dominio.

Los adaptadores son las implementaciones concretas que conectan el sistema con el mundo exterior. Se clasifican en dos categorías: los puertos driving (también llamados primary o inbound) son aquellos por los que el sistema recibe peticiones —el adaptador REST construido con Axum, los handlers de WebSocket para notificaciones en tiempo real, y opcionalmente un adaptador de línea de comandos para tareas administrativas y procesos batch—; los puertos driven (también llamados secondary u outbound) son aquellos a través de los cuales el sistema invoca servicios externos —repositorios sobre PostgreSQL vía sqlx, almacén de sesiones sobre Redis vía fred, servicios criptográficos sobre la biblioteca ring, cliente RFC 3161 hacia Cincel vía reqwest, y dispatcher SMTP vía lettre.

La regla de dependencias se enuncia formalmente así: ningún módulo de una capa interior puede importar símbolos de una capa exterior. El dominio no importa de aplicación ni de adaptadores; la aplicación no importa de adaptadores. Esta regla se garantiza mecánicamente mediante el grafo de dependencias del workspace de Cargo (véase §3.7.6): si un desarrollador intentara importar \texttt{infrastructure::PostgresCaseRepository} desde un módulo del crate \texttt{domain}, el compilador rechazaría la compilación por dependencia inexistente, dado que el \texttt{Cargo.toml} del crate \texttt{domain} no declara a \texttt{infrastructure} entre sus dependencias.

La \cref{fig:hexagonal} presenta la arquitectura hexagonal aplicada al prototipo en su forma conceptual. La versión completa, con todos los puertos nombrados como traits específicos y todos los adaptadores nombrados como structs específicas, se presenta en la Figura A.2 del Anexo A.

\begin{figure}[H]
  \centering
  \includegraphics[width=0.82\linewidth]{image24.png}
  \caption{Arquitectura Hexagonal: capas conceptuales y regla de dependencias. Las dependencias siempre apuntan hacia el dominio.}
  \label{fig:hexagonal}
\end{figure}

El compilador de Rust verifica en tiempo de compilación que PostgresCaseRepository cumple con el contrato declarado por el trait CaseRepository. Cualquier divergencia (firma incorrecta, tipo de retorno incompatible, ausencia de un método requerido) impide la compilación del programa. Esta verificación mecánica elimina una categoría completa de errores que en lenguajes con tipado dinámico solo se manifestarían en tiempo de ejecución.

\subsection{Vista de Componentes y Bounded Contexts}\label{sec:vista-componentes}

La vista de componentes desciende un nivel respecto a la vista de contexto y muestra la organización interna del sistema en términos de bounded contexts —contextos limitados— derivados directamente de las siete agrupaciones funcionales declaradas en §3.2. El concepto de bounded context, introducido por Eric Evans en el contexto de Domain-Driven Design \cite{evans2003}, denota una zona del modelo dentro de la cual los conceptos del lenguaje ubicuo tienen un significado específico y consistente.

El sistema se descompone en siete bounded contexts: Identidad y Acceso (autenticación, RBAC, MFA), Gestión de Casos (expedientes y transiciones procesales), Calendario y Plazos (audiencias y plazos del CNPP), Documentos y Custodia (carga, clasificación y cadena de custodia), Servicios Criptográficos (firma, sellado y verificación), Firma de Contratos (firma con sellado de tiempo) y Operaciones y Auditoría (tableros de control, informes y bitácora). Servicios Criptográficos es un servicio compartido del que dependen tanto Documentos y Custodia como Firma de Contratos.

\begin{figure}[H]
  \centering
  \includegraphics[width=0.9\linewidth]{image25.png}
  \caption{Vista de componentes: siete contextos limitados del sistema y sus dependencias principales.}
  \label{fig:componentes}
\end{figure}

El mapeo entre cada bounded context, los casos de uso que implementa y el crate del workspace de Cargo donde reside su código se presenta en la \cref{tab:bounded-contexts}. Esta correspondencia permite que las modificaciones a un bounded context queden contenidas dentro de un crate específico, facilitando la atribución de cambios y la asignación de responsabilidades.

\begin{table}[H]
  \centering
  \caption{Mapeo entre contextos delimitados, casos de uso y módulos del workspace de Cargo.}
  \label{tab:bounded-contexts}
  \begin{tabularx}{\linewidth}{@{}L{0.30\linewidth}C{0.18\linewidth}X@{}}
    \toprule
    \textbf{Contexto delimitado} & \textbf{Casos de uso} & \textbf{Crate / módulo} \\
    \midrule
    Identidad y Acceso & CU-01..04 & \texttt{domain::identity}, \texttt{application::identity} \\
    Gestión de Casos & CU-05..07 & \texttt{domain::case}, \texttt{application::case} \\
    Calendario y Plazos & CU-08..10 & \texttt{domain::calendar}, \texttt{application::calendar} \\
    Documentos y Custodia & CU-11..12 & \texttt{domain::document}, \texttt{application::document} \\
    Servicios Criptográficos & (transversal) & \texttt{domain::crypto}, \texttt{infrastructure::crypto} \\
    Firma de Contratos & CU-13..14 & \texttt{domain::contract}, \texttt{application::contract} \\
    Operaciones y Auditoría & CU-15..18 & \texttt{domain::audit}, \texttt{application::audit} \\
    \bottomrule
  \end{tabularx}
\end{table}

\subsection{Estructura del Workspace de Cargo}\label{sec:workspace-cargo}

El proyecto se organiza como un workspace de Cargo —el sistema de compilación y gestión de paquetes de Rust— compuesto por cinco crates miembros con responsabilidades estrictamente delimitadas y dependencias permitidas que enforce mecánicamente la regla de la arquitectura hexagonal. Esta estructura se documenta en el \texttt{Cargo.toml} raíz del proyecto y se valida en cada compilación.

Los cinco crates son: domain (entidades, value objects, traits/puertos y errores de dominio; sus únicas dependencias externas son serde, thiserror, time, uuid y zeroize), application (casos de uso, comandos y consultas; depende exclusivamente de domain y de async-trait), infrastructure (adaptadores driven que implementan los traits del dominio; depende de domain y application), web (adaptadores driving HTTP construidos con Axum, middleware Tower, DTOs y serialización HTTP; depende de application y domain), y bin (composition root que instancia los adaptadores concretos y los inyecta en los casos de uso; depende de todos los anteriores). La \cref{fig:workspace} representa esta estructura y las dependencias que el grafo de Cargo permite o prohíbe.

\begin{figure}[H]
  \centering
  \includegraphics[width=0.9\linewidth]{image26.png}
  \caption{Estructura del workspace de Cargo: cinco crates miembros con dependencias permitidas (flechas negras) y prohibidas (flechas rojas punteadas).}
  \label{fig:workspace}
\end{figure}

La consecuencia técnica de esta organización es que las dependencias prohibidas no son una convención que el equipo deba recordar respetar: son una imposibilidad mecánica. Si un módulo del crate domain intenta importar un símbolo del crate infrastructure mediante una declaración \texttt{use infrastructure::PostgresCaseRepository}, el compilador emite el error «unresolved import» porque el \texttt{Cargo.toml} de domain no incluye infrastructure entre sus dependencias. La regla de dependencias se verifica mecánicamente por el compilador en cada build.

El sistema de Cargo features se utiliza para alternar adaptadores intercambiables en tiempo de compilación. Por ejemplo, la feature cincel-tsa activa el adaptador para el PSC Cincel, mientras que una futura feature seguridata-tsa activaría un adaptador alternativo para el PSC SeguriData. La selección se realiza al compilar el binario, manteniendo un único path de código para cada despliegue específico y evitando la complejidad de un sistema de inyección de dependencias en runtime.

\subsection{Stack Tecnológico}\label{sec:stack-tecnologico}

Esta subsección documenta y justifica las tecnologías concretas seleccionadas para implementar el sistema. Cada elección se sustenta en los atributos de calidad priorizados en §3.7.1 y se compara explícitamente con sus alternativas serias del ecosistema actual. Las decisiones más estructurantes se registran adicionalmente como ADRs formales en §3.7.10.

\subsubsection{Justificación de Rust como lenguaje del backend}\label{sec:rust-backend}

El backend del prototipo se implementa en Rust \cite{klabnik2023}. La selección responde principalmente a los atributos de seguridad e integridad priorizados en §3.7.1, en cuanto Rust elimina por construcción —mediante su sistema de ownership y verificación en tiempo de compilación— las clases de errores de memoria que han sido históricamente la fuente principal de vulnerabilidades explotables en software de sistemas. Microsoft ha documentado que aproximadamente el 70\% de las vulnerabilidades reportadas en sus productos a lo largo de una década correspondieron a errores de manejo de memoria \cite{msrc2019}; un porcentaje similar ha sido reportado por el equipo de seguridad de Chromium. Rust elimina categóricamente estos errores sin recurrir a un recolector de basura, lo que lo distingue de Go, Java o cualquier lenguaje con runtime managed.

Frente a Go, Rust ofrece tipos algebraicos (Result, Option, enums con datos asociados) que permiten modelar errores y casos especiales como parte del sistema de tipos, en lugar de mediante valores centinela o múltiples valores de retorno como hace Go. Esta expresividad es particularmente valiosa para un dominio en el que los errores deben ser explícitos y exhaustivos —una operación de firma puede fallar por certificado revocado, por reloj fuera de sincronía con la TSA, por documento que excede el tamaño máximo, entre otros casos—, y donde olvidar manejar uno de ellos no debe ser silencioso. Rust también ofrece zero-cost abstractions, permitiendo construir abstracciones de alto nivel (traits genéricos, iteradores) sin penalización en tiempo de ejecución, lo que Go renuncia a soportar por su modelo de simplicidad explícita.

Frente a Java o Kotlin con Spring, Rust evita el costo del runtime JVM —tiempo de arranque del orden de segundos, huella de memoria elevada, dependencia de un intérprete de bytecode— que es desproporcionado para un prototipo desarrollado en un equipo de desarrollo individual. Frente a Node.js con TypeScript, Rust proporciona tipado estático verificado en compilación (TypeScript es type erasure: los tipos desaparecen en runtime y no protegen contra errores reales del mundo) y rendimiento criptográfico sustancialmente superior, esencial para un sistema cuyo eje son las operaciones de firma, hash y cifrado intensivo. Frente a Python con FastAPI, Rust elimina el Global Interpreter Lock que limita la concurrencia real de Python, y ofrece un ecosistema criptográfico que ha sido sometido a auditorías formales independientes (particularmente la biblioteca ring, que utilizamos para las primitivas).

Una propiedad operacional adicional de Rust pertinente al alcance del prototipo es que el binario resultante es ejecutable único, sin dependencia de runtime, intérprete ni gestor de paquetes en producción. El despliegue se reduce a copiar un archivo de aproximadamente 15 MB al servidor objetivo, lo que simplifica la operación y reduce significativamente la superficie de ataque del entorno desplegado.

\subsubsection{Backend: Axum sobre el ecosistema tokio, tower y hyper}\label{sec:axum}

Como framework web se selecciona Axum, desarrollado por el equipo de tokio. Axum se construye sobre tres capas estándar de facto en el ecosistema Rust: tokio como runtime asíncrono multi-thread con scheduler work-stealing, tower como abstracción de middleware componible basada en el trait \texttt{Service<Request, Response>}, y hyper como implementación de los protocolos HTTP/1.1 y HTTP/2. Esta estratificación permite que cada nivel de la pila se evalúe y evolucione independientemente, y que los componentes desarrollados para tower puedan reutilizarse fuera del contexto HTTP.

Axum se distingue de alternativas como actix-web o rocket por no ser opinado respecto a la organización del código de la aplicación, lo que facilita la implementación del estilo arquitectónico hexagonal seleccionado en §3.7.3 sin tener que combatir las convenciones del framework. Los handlers de Axum reciben sus dependencias mediante extractores tipados (extractors), permitiendo inyectar los casos de uso de la capa de aplicación de manera explícita y verificable en compilación.

El pipeline de middleware Tower se configura para ejecutar, en orden, los siguientes interceptores transversales: CORS para el control de origen cruzado de peticiones del navegador; Rate Limiting con ventana deslizante para protección contra abuso; Tracing con propagación de identificador de petición para observabilidad; autenticación JWT para verificación de identidad del solicitante; logging estructurado para auditoría operativa. Cada middleware es composable —su ejecución se especifica declarativamente al construir el router— y verificable en compilación por el sistema de tipos.

\subsubsection{Frontend: Astro con Svelte Islands}\label{sec:frontend-astro}

El frontend del sistema se implementa con Astro como meta-framework y Svelte como librería de componentes interactivos, siguiendo el patrón Islands Architecture propuesto por Sylor-Miller y Miller \cite{sylormiller2019}. Bajo este patrón, las páginas se renderizan como HTML estático por defecto, y solo los componentes que requieren interactividad (calendario jurídico, dashboard, interfaz de verificación de firma) se hidratan selectivamente como islas reactivas. Esta estrategia minimiza el JavaScript enviado al navegador a aquello estrictamente necesario, mejorando los tiempos de carga y reduciendo la superficie de ataque del cliente.

Frente a Next.js con React, Astro y Svelte tienen tres ventajas específicas para el alcance del prototipo. Primero, Svelte compila a JavaScript imperativo sin Virtual DOM ni runtime de framework, produciendo bundles considerablemente más ligeros que React. Segundo, la reactividad de Svelte es sintáctica —se declara con \texttt{\$:} y asignaciones a variables— lo que reduce el boilerplate frente al modelo de hooks de React. Tercero, Astro está diseñado desde su concepción para emitir el mínimo JavaScript posible, mientras que Next.js parte del modelo de SPA hidratada completa y solo recientemente ha introducido los Server Components para mitigar ese problema.

Frente a SvelteKit puro, Astro se elige porque las páginas mayoritariamente informativas del sistema (landing, documentación, vistas de solo lectura) no se benefician de hidratación completa; ahí Astro emite HTML estático directamente. SvelteKit, en cambio, hidrata todas las páginas como SPA, perdiendo la oportunidad de servir contenido estático puro. La combinación Astro + Svelte aprovecha lo mejor de ambos: contenido estático por defecto y reactividad rica solo donde se necesita.

\subsubsection{Persistencia híbrida}\label{sec:persistencia-hibrida}

El sistema emplea un modelo de persistencia híbrido que combina PostgreSQL como base de datos relacional para los datos con valor probatorio, con Redis como almacén clave-valor en memoria para datos volátiles y comunicación en tiempo real. La justificación completa de esta separación, los esquemas de cada motor, y el mecanismo de encadenamiento criptográfico SHA-256 que garantiza la inmutabilidad de la bitácora de auditoría se documentan en §3.8 (Modelado de Base de Datos Híbrida). Los crates utilizados para acceder a cada motor son sqlx para PostgreSQL —que ofrece verificación de queries SQL en tiempo de compilación contra el esquema real— y fred para Redis —cliente asíncrono con soporte nativo para pub/sub.

\subsubsection{Inventario completo de crates}\label{sec:inventario-crates}

El proyecto utiliza un conjunto curado de crates externos del ecosistema Rust, todos publicados en crates.io con versiones pinneadas en los archivos \texttt{Cargo.toml} de cada crate del workspace. Los crates se agrupan por capa arquitectónica para facilitar la auditoría y para enforce las reglas de dependencias entre capas discutidas en §3.7.6. La \cref{tab:inventario-crates} enumera los crates principales del proyecto, agrupados por capa y con sus respectivas funciones. La versión mínima soportada del compilador (Minimum Supported Rust Version, MSRV) es Rust 1.78 estable.

\begin{longtable}{@{}L{0.22\linewidth}L{0.22\linewidth}L{0.46\linewidth}@{}}
  \caption{Inventario de crates externos utilizados por el proyecto, agrupados por capa arquitectónica.}
  \label{tab:inventario-crates} \\
  \toprule
  \textbf{Capa} & \textbf{Crate} & \textbf{Función} \\
  \midrule
  \endfirsthead
  \toprule
  \textbf{Capa} & \textbf{Crate} & \textbf{Función} \\
  \midrule
  \endhead
  \midrule
  \multicolumn{3}{r@{}}{\emph{Continúa en la página siguiente}} \\
  \endfoot
  \bottomrule
  \endlastfoot
  Dominio & serde & Serialización/deserialización tipada \\
  Dominio & thiserror & Errores tipados por capa \\
  Dominio & time & Manejo de fechas con zona horaria \\
  Dominio & uuid & Identificadores únicos universales \\
  Dominio & zeroize & Borrado de material secreto en memoria al liberar (higiene de memoria) \\
  Aplicación & async-trait & Traits asíncronos \\
  Infraestructura cripto & ring & Hash, AES-GCM, HMAC (auditado) \\
  Infraestructura cripto & x509-cert + der & Manejo de certificados X.509 RFC 5280 \\
  Infraestructura cripto & rsa & Firma digital RSA-3072 PKCS\#1 v1.5 \\
  Infraestructura cripto & argon2 & Hashing de contraseñas (RFC 9106) \\
  Infraestructura cripto & totp-rs & Autenticación multifactor TOTP (RFC 6238) \\
  Infraestructura cripto & cms + x509-tsp & Parseo del token de sellado de tiempo RFC 3161 (CMS SignedData y TSTInfo) \\
  Infraestructura datos & sqlx & Cliente PostgreSQL con verificación en compile-time \\
  Infraestructura datos & fred & Cliente Redis asíncrono con pub/sub \\
  Infraestructura datos & fs2 & Bloqueo advisory de archivos compartidos entre procesos \\
  Infraestructura red & reqwest & Cliente HTTP asíncrono (Cincel TSA) \\
  Infraestructura red & lettre & Cliente SMTP con TLS y SASL \\
  Web/HTTP & axum & Framework web sobre tokio/tower/hyper \\
  Web/HTTP & tower / tower-http & Middleware composable \\
  Web/HTTP & tokio & Runtime asíncrono multi-thread \\
  Interfaz de línea de comandos & clap & Análisis declarativo de argumentos y subcomandos \\
  Observabilidad & tracing + tracing-subscriber & Logs estructurados con spans \\
  Validación & validator & Validación de DTOs HTTP \\
  Configuración & config + dotenvy & Carga de configuración desde entorno \\
  Errores (boundaries) & anyhow & Errores polimórficos en \texttt{main.rs} y tests \\
  Testing & mockall, proptest, rstest, wiremock & Mocks, pruebas basadas en propiedades, fixtures, mock HTTP \\
\end{longtable}

\subsection{Patrones Arquitectónicos Transversales}\label{sec:patrones-transversales}

Esta subsección documenta cuatro patrones que atraviesan transversalmente todas las capas de la arquitectura. Su tratamiento explícito y consistente es lo que permite que el sistema mantenga coherencia operativa en escala: errores tipados, concurrencia estructurada, inyección de dependencias por composición, y observabilidad estructurada.

\subsubsection{Manejo de errores}\label{sec:manejo-errores}

El sistema adopta una estrategia de errores tipados por capa, implementada con la macro derive de thiserror. Cada capa define su propio tipo de error con variantes específicas: DomainError captura las violaciones de invariantes de negocio (caso inválido, transición prohibida, firma rechazada); ApplicationError envuelve errores de orquestación de casos de uso e incluye DomainError como variante; los crates de infraestructura definen errores específicos de su backend (PostgresError, CincelError, SmtpError); finalmente, ApiError representa errores HTTP y mapea desde las capas anteriores mediante el trait IntoResponse de Axum.

Esta tipificación contrasta con el uso indiscriminado de anyhow, que se reserva exclusivamente para el binario principal (\texttt{bin/main.rs}) y para las pruebas, donde la conveniencia del polimorfismo de errores supera la necesidad de discriminar tipos. En cambio, en los crates de dominio, aplicación e infraestructura, los errores deben ser exhaustivos y discriminables porque cada caller necesita decidir qué hacer ante cada variante específica.

\subsubsection{Modelo de concurrencia asíncrona}\label{sec:concurrencia-asincrona}

El sistema utiliza tokio como runtime asíncrono multi-thread con scheduler work-stealing, configurado por defecto con tantos worker threads como núcleos lógicos disponibles. Las operaciones I/O bound (consultas a PostgreSQL, llamadas a Cincel, envíos SMTP) se ejecutan asíncronamente sin bloquear los workers, permitiendo concurrencia masiva con un consumo modesto de threads del sistema operativo. La concurrencia entre tareas se estructura mediante \texttt{tokio::spawn} con joins explícitos —siguiendo el principio de structured concurrency— para garantizar que ninguna tarea escape del scope de su parent y produzca leaks o panics no manejados.

Las operaciones que pueden quedar bloqueadas indefinidamente (invocaciones HTTP a sistemas externos, esperas a locks distribuidos) se envuelven con \texttt{tokio::time::timeout} para garantizar terminación acotada. Las operaciones cancelables utilizan CancellationToken propagado desde el handler raíz, permitiendo que un cliente que cierra la conexión HTTP libere inmediatamente todos los recursos asociados a su petición. Los traits Send + Sync requeridos para los puertos garantizan en tiempo de compilación que cualquier implementación es segura para usar entre threads, eliminando una categoría completa de errores de concurrencia.

\subsubsection{Inyección de dependencias por composición}\label{sec:inyeccion-dependencias}

La inyección de dependencias se realiza por composición explícita en el composition root —el archivo \texttt{bin/main.rs}—, sin recurrir a un contenedor de DI en runtime. Los casos de uso reciben sus dependencias como parámetros genéricos restringidos por traits (\texttt{impl Trait}) o como trait objects (\texttt{Arc<dyn Trait>}). Por ejemplo, el caso de uso SignContract se define como \texttt{SignContract<R, S, T>} donde R: CaseRepository, S: DocumentSigner, T: TimestampService; \texttt{main.rs} construye PostgresCaseRepository, RsaPkcs1Signer y CincelTsaAdapter, los inyecta en SignContract, y registra el resultado en el router de Axum. Esta estrategia es completamente verificable en tiempo de compilación —si un adaptador no satisface el contrato del puerto, el programa no compila— y elimina la categoría de errores característicos de los frameworks DI (dependencias no resueltas, ciclos detectados solo en runtime, configuraciones mal escritas).

\subsubsection{Observabilidad estructurada}\label{sec:observabilidad}

El sistema utiliza el ecosistema tracing para producir logs estructurados con spans tipados. Cada petición HTTP genera un span raíz con un \texttt{request\_id} propagado a través de todas las invocaciones asíncronas mediante el contexto de tokio. Los logs se emiten en formato JSON line-delimited (JSONL), facilitando su ingesta por agregadores como Vector, Loki o Elasticsearch. Los niveles de verbosidad son configurables por módulo a través de la variable de entorno \texttt{RUST\_LOG}, permitiendo aumentar temporalmente el detalle de una capa específica sin reiniciar el servicio. Para trazabilidad distribuida, el crate tracing-opentelemetry permite exportar spans a backends compatibles (Jaeger, Tempo) sin modificar el código de la aplicación. La política de redacción de información sensible (referencia cruzada a RNF-03) está implementada como un layer personalizado del subscriber, que aplica patrones regex sobre los campos antes de su emisión para sanitizar secretos y material criptográfico que pudieran filtrarse.

\subsection{Vista de Seguridad: Fronteras de Confianza}\label{sec:vista-seguridad}

Esta subsección desarrolla la vista de seguridad arquitectónica, complementaria al diseño criptográfico detallado en §3.11. Mientras §3.11 documenta los algoritmos criptográficos específicos (RSA-3072, AES-256-GCM, SHA-256, RFC 3161, PKI con CA interna) y sus parámetros, esta subsección documenta cómo esos mecanismos se distribuyen entre las distintas zonas de confianza del sistema y qué controles se aplican en cada cruce de frontera.

El modelo adoptado sigue el principio de defensa en profundidad formalizado por Saltzer y Schroeder \cite{saltzer1975} y desarrollado operativamente en el OWASP Application Security Verification Standard \cite{owasp2021}. La Figura A.1 en el Anexo A presenta el modelo completo con cuatro zonas de confianza y los controles de seguridad aplicados en cada cruce.

El sistema se modela como cuatro zonas de confianza separadas por tres fronteras explícitas. La primera zona, la red pública, comprende todo aquello que no está bajo control directo del sistema: el navegador del usuario final, los actores adversarios potenciales, y los servicios externos como Cincel PSC y los proveedores SMTP que aunque sean legítimos no operan dentro del perímetro de seguridad del prototipo. La segunda zona, el perímetro de aplicación, contiene el reverse proxy con terminación TLS y la API pública expuesta. La tercera zona, el núcleo aplicativo, contiene el binario Rust con sus tres módulos internos (web, application, domain). La cuarta y más confiable zona, datos persistentes y secretos, contiene PostgreSQL, Redis y el subsistema de gestión de claves criptográficas.

El cruce de la primera frontera —de la red pública al perímetro de aplicación— se controla mediante terminación TLS 1.3 con suites de cifrado modernas; rate limiting con ventana deslizante para protección contra abuso de la API; un Web Application Firewall que filtra patrones de ataque conocidos; y cabeceras CORS estrictas que restringen el origen de las peticiones legítimas. La superficie expuesta a esta frontera es deliberadamente mínima: solo los endpoints HTTPS estrictamente necesarios.

El cruce de la segunda frontera —del perímetro al núcleo aplicativo— se controla mediante autenticación basada en JWT (JSON Web Token) firmado por la propia aplicación; validación exhaustiva de toda entrada del usuario antes de cualquier procesamiento; y autorización basada en roles (RBAC) que restringe qué operaciones puede invocar cada rol (socio, paralegal, cliente). Cualquier petición que no satisfaga las tres verificaciones es rechazada con un código de error apropiado y registrada en la bitácora de auditoría.

El cruce de la tercera frontera —del núcleo aplicativo a datos persistentes y secretos— se controla mediante conexiones TLS internas entre el binario Rust y los motores de PostgreSQL y Redis; uso obligatorio de consultas parametrizadas mediante sqlx (que verifica en compile-time la conformidad de cada query con el esquema, eliminando la posibilidad de inyección SQL); y cifrado en reposo de todos los documentos mediante el patrón envelope encryption con KEK protegida por el módulo de gestión de claves (cuyo diseño detallado se documenta en §3.11.4).

Los servicios externos se modelan como entidades de la zona no confiable a las que el sistema invoca explícitamente con canales seguros: un PSC externo futuro para sellado de tiempo (RFC 3161) utilizará TLS y verificación de la firma del sello recibido; la TSA local de demostración opera sin salir del perímetro del prototipo; la integración con el proveedor SMTP utiliza TLS más autenticación SASL. En ningún caso se confía en estos servicios para operaciones de seguridad sin verificación local del resultado: un sello recibido se verifica criptográficamente contra la cadena de certificados del emisor antes de ser aceptado e incorporado al expediente.

El modelo de amenazas STRIDE de Shostack \cite{shostack2014} se aplicó como herramienta de revisión sistemática durante el diseño para verificar que cada una de las seis categorías de amenazas (Spoofing, Tampering, Repudiation, Information Disclosure, Denial of Service, Elevation of Privilege) tuviese al menos un control asignado. La matriz resultante demuestra que la combinación de los controles enumerados ofrece cobertura para todas las categorías, con énfasis particular en Tampering (mediante el encadenamiento SHA-256 de la bitácora) y Repudiation (mediante el sellado de tiempo independiente del PSC).

\subsection{Decisiones Arquitectónicas Significativas (ADRs)}\label{sec:adrs}

Las decisiones arquitectónicas más estructurantes del sistema se registran formalmente como Architecture Decision Records (ADRs) en el formato propuesto por Michael Nygard \cite{nygard2011}. Cada ADR documenta de forma autocontenida cuatro elementos: el contexto del problema que motivó la decisión, las alternativas consideradas, la decisión adoptada, y las consecuencias —tanto positivas como negativas— de la decisión. Los seis ADRs siguientes cubren las decisiones cuya reversión implicaría una refactorización significativa.

\subsubsection{ADR-01. Adopción de Arquitectura Hexagonal}\label{sec:adr-01}

Contexto: el sistema requiere alta testabilidad (XP/TDD) y evolvabilidad (sustitución de PSC, motor de BD). Alternativas: N-tier, Clean, Onion, Hexagonal, Microservicios, Monolito Modular sin patrón hexagonal. Decisión: Monolito Modular con Arquitectura Hexagonal de Cockburn. Consecuencias positivas: dominio aislado y testeable, adaptadores intercambiables, regla de dependencias verificable mecánicamente. Consecuencias negativas: ligero overhead de indirección por cada operación que cruza una capa, requiere disciplina del equipo para no introducir atajos.

\subsubsection{ADR-02. Rust como lenguaje del backend}\label{sec:adr-02}

Contexto: prototipo de seguridad con operaciones criptográficas intensivas. Alternativas: Go, Java/Kotlin con Spring, Node.js con TypeScript, Python con FastAPI. Decisión: Rust. Consecuencias positivas: seguridad de memoria sin GC, sistema de tipos algebraicos, ecosistema criptográfico auditado, binario único sin runtime. Consecuencias negativas: curva de aprendizaje del ownership y borrow checker, ecosistema más pequeño que Java o Node, tiempos de compilación más largos.

\subsubsection{ADR-03. Axum como framework web}\label{sec:adr-03}

Contexto: necesidad de un framework HTTP que no impusiera estructura incompatible con el estilo hexagonal. Alternativas: actix-web, rocket, warp. Decisión: Axum sobre tokio/tower/hyper. Consecuencias positivas: framework no opinado, ecosistema tower reutilizable para otros servicios, mantenido por el equipo de tokio. Consecuencias negativas: API ergonómica menos madura que actix-web, documentación todavía en evolución.

\subsubsection{ADR-04. Astro + Svelte Islands para el frontend}\label{sec:adr-04}

Contexto: páginas mayoritariamente informativas con islas de interactividad rica (calendario, dashboard, verificación de firma). Alternativas: Next.js + React, SvelteKit puro, Remix. Decisión: Astro como meta-framework con Svelte como librería de componentes interactivos. Consecuencias positivas: HTML estático por defecto, JavaScript mínimo enviado al cliente, Svelte sin Virtual DOM. Consecuencias negativas: dos sistemas de templating (Astro y Svelte) que el desarrollador debe conocer, ecosistema de componentes Svelte más pequeño que React.

\subsubsection{ADR-05. Monolito Modular sobre Microservicios}\label{sec:adr-05}

Contexto: equipo de desarrollo de dos personas en periodo académico, alcance del prototipo no requiere despliegue distribuido. Alternativas: descomposición en microservicios desde el inicio, monolito modular. Decisión: Monolito Modular siguiendo la recomendación MonolithFirst de Fowler \cite{fowler2015}. Consecuencias positivas: simplicidad operativa, despliegue como binario único, sin overhead de comunicación inter-servicios, facilidad de refactorización. Consecuencias negativas: escalado monolítico (ya aceptable para el alcance del prototipo), imposibilidad de desplegar componentes independientemente.

\subsubsection{ADR-06. Persistencia híbrida PostgreSQL + Redis}\label{sec:adr-06}

Contexto: el sistema maneja tanto datos con valor probatorio (que requieren persistencia transaccional duradera) como datos volátiles para sesiones, caché y notificaciones en tiempo real. Alternativas: PostgreSQL exclusivamente, Redis exclusivamente, PostgreSQL + Redis. Decisión: PostgreSQL para datos persistentes con responsabilidades específicas en §3.8.1, Redis para datos volátiles con TTL conforme a §3.8.2. Consecuencias positivas: cada motor utilizado para lo que mejor hace, baja latencia para sesiones y notificaciones, integridad transaccional para operaciones críticas. Consecuencias negativas: necesidad de operar dos motores de almacenamiento, mayor complejidad operativa, necesidad de coherencia eventual entre ambos.

\subsection{Trazabilidad Arquitectura ↔ Requerimientos}\label{sec:trazabilidad}

Esta subsección presenta la matriz de trazabilidad que cierra el Capítulo 3 demostrando que cada requerimiento funcional y no funcional cuenta con un componente arquitectónico responsable, y que ningún componente arquitectónico carece de justificación en un requerimiento. La \cref{tab:trazabilidad} sintetiza este mapeo en una vista de alto nivel; la matriz completa, con los diecisiete casos de uso (CU-01 a CU-17) y los catorce requerimientos no funcionales (RNF-01 a RNF-14 de §3.4), se entrega como anexo electrónico complementario al documento.

\begin{table}[H]
  \centering
  \caption{Mapeo de alto nivel entre componentes arquitectónicos y los requerimientos que satisfacen.}
  \label{tab:trazabilidad}
  \begin{tabularx}{\linewidth}{@{}L{0.42\linewidth}X@{}}
    \toprule
    \textbf{Componente arquitectónico} & \textbf{Requerimientos atendidos} \\
    \midrule
    Bounded Context «Identidad y Acceso» & CU-01..04, RNF-01, RNF-02 \\
    Bounded Context «Gestión de Casos» & CU-05..07, RNF-04 (parcial) \\
    Bounded Context «Calendario y Plazos» & CU-08..10 \\
    Bounded Context «Documentos y Custodia» & CU-11..12, RNF-03, RNF-04 \\
    Bounded Context «Servicios Criptográficos» & CU-11..14, RNF-03 \\
    Bounded Context «Firma de Contratos» & CU-13..14 \\
    Bounded Context «Operaciones y Auditoría» & CU-15..17, RNF-04, RF-20 \\
    Crate \texttt{domain} con traits (puertos) & Soporta evolvabilidad transversal \\
    Workspace de Cargo con regla de dependencias & Mantenibilidad transversal \\
    Pipeline Tower middleware & RF-01 (autenticación y sesiones), seguridad cross-cutting \\
    Modelo de fronteras de confianza (§3.7.9) & RNF-01, RNF-02, RNF-03 \\
    Subsistema observabilidad (tracing) & RF-20 (bitácora de actividad) \\
    \bottomrule
  \end{tabularx}
\end{table}

\subsubsection{Fitness functions}\label{sec:fitness-functions}

Siguiendo el enfoque de arquitecturas evolutivas formalizado por Ford, Parsons y Kua \cite{ford2023}, la coherencia entre la arquitectura prescrita y la arquitectura efectiva del código se verifica mediante un conjunto de fitness functions: pruebas automatizadas que comprueban propiedades arquitectónicas que no podrían verificarse mediante pruebas funcionales tradicionales. Estas fitness functions se ejecutan como parte del pipeline de Integración Continua y, en caso de fallo, bloquean la integración del cambio en la rama principal.

Las fitness functions implementadas para el prototipo son: (1) verificación de que ningún módulo del crate domain importa símbolos de los crates infrastructure ni web (verificada por el compilador en cada build); (2) verificación de que todo handler HTTP envuelve su lógica en un span de tracing (verificada con un lint personalizado de clippy); (3) verificación de que ningún log emitido contiene literales que matcheen patrones de claves criptográficas (verificada con un test que ejecuta toda la suite e inspecciona los logs producidos); (4) verificación de que cada trait que define un puerto driven tiene al menos una implementación concreta y al menos un mock para tests (verificada por inspección del workspace en cada build); (5) verificación de que el binario compilado no excede 25 MB (verificada por el pipeline de release).


\section{Modelado de Base de Datos Híbrida}\label{sec:bd-hibrida}

El sistema emplea un modelo de persistencia híbrido que combina \textbf{PostgreSQL} (base de datos relacional) con \textbf{Redis} (almacén clave-valor en memoria), con responsabilidades estrictamente delimitadas. Esta decisión responde a los requisitos funcionales y no funcionales identificados en el protocolo, particularmente RF-20 (bitácora de actividad), RNF-04 (integridad append-only) y la necesidad de notificaciones en tiempo real.

\subsection{PostgreSQL: Datos Persistentes con Valor Probatorio}\label{sec:postgresql}

PostgreSQL se selecciona como motor relacional principal por las siguientes razones: (1) soporte nativo de \textbf{JSONB} para metadatos variables de expedientes, lo que permite modelar la diversidad del sistema acusatorio oral mexicano (CNPP, Art. 211-215) sin sacrificar la integridad referencial; (2) \textbf{RULE} para enforce de append-only en la bitácora de auditoría, impidiendo operaciones UPDATE y DELETE a nivel de motor; (3) extensiones criptográficas nativas (pgcrypto); (4) tipo de dato INET para registro de direcciones IP en la bitácora.

La bitácora de auditoría (audit\_log) implementa un mecanismo de \textbf{encadenamiento SHA-256}: cada registro incluye un campo sha256\_chain calculado como el hash SHA-256 (véase §2.3.1) del valor sha256\_chain del registro anterior concatenado con los datos del registro actual, es decir, $\mathrm{sha256\_chain}(n) = \mathrm{SHA\text{-}256}(\mathrm{sha256\_chain}(n-1) \,\|\, \mathrm{registro}(n))$. Esto crea una cadena verificable de integridad que permite detectar cualquier manipulación retroactiva, cumpliendo con el Art. 33 del Código de Comercio (retención mínima de 10 años).

\subsubsection{Decisión de diseño: Single-Tenant}

El sistema se diseña para un solo despacho legal (single-tenant), eliminando la complejidad de aislamiento multi-tenant contemplada en el protocolo original (cuyos requerimientos se reformulan en RNF-01 y RNF-05 como aislamiento y privacidad de la información del despacho dentro de su instancia). Esta simplificación es coherente con el alcance del prototipo y reduce significativamente la superficie de ataque.

\subsection{Redis: Datos Volátiles y Comunicación en Tiempo Real}\label{sec:redis}

Redis se emplea exclusivamente para datos efímeros que \textbf{no requieren persistencia a largo plazo ni tienen valor probatorio}:

\begin{itemize}
  \item \textbf{Sesiones JWT}: almacenamiento de tokens de sesión con TTL configurable (24h por defecto), permitiendo invalidación inmediata por revocación.
  \item \textbf{Caché de permisos}: materialización temporal de las autorizaciones del usuario (TTL 5 minutos) para evitar consultas repetidas a PostgreSQL en cada petición.
  \item \textbf{Pub/Sub para notificaciones}: canal de comunicación entre el backend y los clientes WebSocket. Cuando ocurre un evento relevante (nuevo documento, firma completada, deadline próximo), se publica en un canal Redis que el adaptador WebSocket consume y retransmite al navegador del usuario.
  \item \textbf{Rate limiting}: contadores con ventana deslizante para protección contra abuso de la API.
  \item \textbf{Tokens temporales}: tokens de un solo uso para restablecimiento de contraseña y verificación MFA (TTL 15 minutos).
\end{itemize}

\subsection{Diagrama del Modelo de Datos}\label{sec:diagrama-datos}

La \cref{fig:modelo-datos} presenta el modelo de datos completo, mostrando las tablas de PostgreSQL (izquierda) con sus relaciones y las estructuras de Redis (derecha) con sus políticas de TTL.

\begin{figure}[H]
  \centering
  \includegraphics[width=0.9\linewidth]{image27.png}
  \caption{Modelo de datos híbrido: PostgreSQL para datos persistentes con valor probatorio, Redis para datos volátiles.}
  \label{fig:modelo-datos}
\end{figure}

Nótese que la tabla audit\_log está resaltada en rojo para indicar su naturaleza especial: es la única tabla del sistema con restricción RULE que impide UPDATE y DELETE. Su campo sha256\_chain implementa el encadenamiento criptográfico descrito en la \cref{sec:postgresql}. El campo metadata JSONB en la tabla cases permite almacenar atributos específicos de cada tipo de procedimiento sin necesidad de alterar el esquema relacional.

\section{Diagrama de Despliegue y Stack Tecnológico}\label{sec:despliegue}

La \cref{fig:despliegue} presenta la vista de despliegue del sistema, ilustrando la interacción entre las capas del frontend, el backend y los servicios de infraestructura.

\begin{figure}[H]
  \centering
  \includegraphics[width=0.9\linewidth]{image28.png}
  \caption{Diagrama de despliegue: Astro+Svelte (cliente) → Axum/Rust (backend) → PostgreSQL + Redis + servicios externos.}
  \label{fig:despliegue}
\end{figure}

El backend Rust compila a un \textbf{binario único} de aproximadamente 15 MB que incluye todas las dependencias. El binario Rust se ejecuta directamente sin runtime externo, eliminando la necesidad de intérpretes, servidores web adicionales o gestores de paquetes. Esto simplifica significativamente el despliegue y reduce la superficie de ataque del servidor.

La capa de \textbf{Tower Middleware} opera como un pipeline composable donde cada middleware es una capa que envuelve al siguiente. El flujo de una petición HTTP es: CORS → Rate Limiting → Autenticación JWT → Logging → Handler de Axum → Dominio. Esta composición es declarativa y verificada en tiempo de compilación por el sistema de tipos de Rust.

La comunicación con \textbf{Cincel PSC} para sellado de tiempo RFC 3161 se realiza a través de reqwest (cliente HTTP asíncrono). El diseño del puerto TimestampService como trait abstracto permite sustituir Cincel por cualquier otra TSA (SeguriData, CECOBAN; véase §2.3.3.4) implementando únicamente un nuevo adaptador, sin modificar el dominio ni los demás componentes del sistema.

Las conexiones a \textbf{PostgreSQL} se gestionan mediante sqlx, que ofrece verificación de queries SQL en tiempo de compilación contra el esquema real de la base de datos. Las conexiones a \textbf{Redis} se gestionan mediante fred, un cliente asíncrono compatible con cluster que soporta pub/sub nativo para las notificaciones en tiempo real.

\section{Modelo de Datos Jurídico-Técnico}\label{sec:modelo-juridico}

La fundamentación del proceso penal expuesta en el marco teórico tiene una consecuencia directa sobre el diseño del prototipo: la base de datos no puede ser un conjunto de tablas genéricas, sino que debe reflejar fielmente las dependencias y restricciones que el CNPP impone. Esta sección articula los vínculos explícitos entre la teoría jurídica y los requerimientos funcionales.

Un caso en el prototipo no es un registro estático, sino una entidad que evoluciona a través de los estados que el CNPP define como etapas \cite{cnpp2021}. La transición de un estado a otro solo es posible mediante la carga del documento procesal habilitante: de Investigación a Intermedia únicamente mediante el escrito de acusación; de Intermedia a Juicio Oral mediante el auto de apertura. El historial de transiciones debe mantenerse de forma inmutable como parte del expediente, dado que un tribunal revisor puede requerir verificar la secuencia procesal completa.

\section{Diseño Criptográfico del Prototipo}\label{sec:diseno-cripto}

Esta sección documenta las decisiones de diseño criptográfico adoptadas para el prototipo, fundamentadas en los conceptos teóricos expuestos en el capítulo 2 y en las restricciones impuestas por el marco legal mexicano. Cada decisión se justifica con referencia a los estándares y publicaciones del NIST, las RFC del IETF y las guías de OWASP que las sustentan.

\subsection{Nivel de Seguridad Objetivo}\label{sec:nivel-seguridad}

El prototipo adopta un nivel de fortaleza de seguridad de 128 bits conforme a NIST SP 800-57 Part 1 Rev 5 \cite{barker2020}. Este nivel unificado exige que cada primitivo criptográfico ofrezca al menos 128 bits de resistencia: SHA-256 proporciona 128 bits de resistencia a colisiones, AES-256 proporciona 256 bits de fortaleza, y RSA-3072 proporciona aproximadamente 128 bits. La elección de 128 bits como objetivo coherente garantiza que ningún eslabón de la cadena criptográfica sea más débil que los demás. NIST considera este nivel seguro al menos hasta 2030 para aplicaciones de uso general.

La Tabla 2 del NIST SP 800-57 \cite{barker2020} clasifica los niveles de fortaleza en función de los tamaños de clave recomendados para cada familia de algoritmos. El nivel de 112 bits, correspondiente a RSA-2048, está clasificado como legacy-use: aceptable para procesar información existente, pero no recomendado para proteger información nueva más allá de 2030. El nivel de 128 bits, correspondiente a RSA-3072, AES-128 o AES-256 y SHA-256, es el primer nivel que NIST clasifica como aceptable para uso general sin restricción temporal. Dado que el prototipo gestiona documentos con valor probatorio que pueden requerirse como prueba judicial años después de su firma, la elección del nivel de 128 bits proporciona un margen de seguridad adecuado frente a la evolución del poder computacional. Cabe señalar que la e.firma del SAT utiliza RSA-2048 (112 bits), un nivel inferior al adoptado por el prototipo; la libertad de elegir parámetros que otorga la CA interna permite superar ese piso normativo.

La coherencia entre primitivos se verifica de la siguiente manera: SHA-256 proporciona 128 bits de resistencia a colisiones y 256 bits de resistencia a preimagen; RSA-3072 proporciona aproximadamente 128 bits de fortaleza para firma digital; AES-256 proporciona 256 bits de fortaleza para cifrado simétrico, superando el nivel objetivo pero seleccionado por ser el tamaño estándar de clave para GCM en aplicaciones de nivel 128 bits o superior; y Argon2id con hash de salida de 32 bytes proporciona 256 bits de resistencia a preimagen para el almacenamiento de contraseñas. Ningún primitivo del sistema presenta una fortaleza inferior a 128 bits, lo que garantiza la ausencia de un eslabón débil en la cadena criptográfica. Esta propiedad de coherencia es especialmente relevante en un sistema que combina múltiples servicios criptográficos —integridad, confidencialidad, no repudio y autenticación—, ya que un atacante racional concentraría su esfuerzo en el primitivo más débil.

\subsection{Firma Digital}\label{sec:firma-digital}

El esquema de firma digital emplea RSA-3072 con SHA-256 y padding PKCS\#1 v1.5, con sha256\-With\-RSA\-Encryption como OID del algoritmo de firma. Se descartó ECDSA P-256 tras verificar que la Política de Certificados de la ACR2-SE \cite{economia_acr2se} establece que los certificados emitidos bajo la jerarquía PKI mexicana deben contener claves públicas RSA con sha256WithRSAEncryption como algoritmo de firma; ECDSA no está contemplado para certificados legalmente reconocidos bajo PSC mexicano. La elección de PKCS\#1 v1.5 sobre PSS se justifica por la alineación con el esquema que usa la PKI mexicana real (ACR2-SE y e.firma), lo que facilita la migración futura a certificados emitidos por un PSC sin rediseñar el verificador.

Los certificados X.509 v3 de los socios del despacho son emitidos por una Autoridad Certificadora (CA) interna del prototipo, implementada con OpenSSL scriptado en una topología de una sola capa: un certificado raíz autofirmado con vigencia de cinco años firma directamente los certificados de los socios con vigencia de un año. Esta elección expone al equipo al ciclo de vida completo de una PKI —generación de la clave raíz, emisión de CSRs (Certificate Signing Request, solicitudes de firma de certificado), firma de certificados, validación de cadena y revocación— que una integración con la e.firma del SAT o un PSC comercial ocultaría detrás de una API. El módulo de firma acepta cualquier certificado X.509 RSA + SHA-256, de modo que en un despliegue productivo la CA interna podría sustituirse por la e.firma o un PSC autorizado sin modificar la lógica de verificación.

La validación del estado de revocación se realiza mediante una Lista de Certificados Revocados (CRL) conforme a RFC 5280 \cite{cooper2008}, publicada como archivo PEM en el endpoint GET /pki/crl.pem con vigencia de siete días. La CRL se regenera bajo demanda tras cada revocación y periódicamente antes de su expiración. Se descartó OCSP por la escala reducida del prototipo (\textasciitilde{}2--5 certificados), donde la CRL es de tamaño despreciable y no requiere un servicio persistente adicional.

La elección de una CA interna sobre la e.firma del SAT o un PSC comercial responde a múltiples consideraciones. En primer lugar, las firmas generadas con certificados de la CA interna emplean los mismos mecanismos técnicos que las generadas con certificados de la e.firma o de un PSC: RSA-3072 con SHA-256, estructura X.509 v3 y verificación de cadena. Lo que cambia es el emisor y, con él, el nivel de confianza jurídica atribuible al certificado. Esta elección permite demostrar integridad y autoría dentro del alcance académico, pero no debe presentarse como equivalente a una firma avanzada. En segundo lugar, un sello RFC 3161 emitido por un PSC autorizado puede aportar la independencia temporal prevista por la NOM-151 \cite{adams2001,nom151_2016}; un sello emitido por la TSA interna conserva el mismo flujo criptográfico, pero permanece dentro de la frontera de confianza del prototipo.

En tercer lugar, usar la e.firma habría obligado a que el socio subiera su archivo .key y su contraseña al prototipo, o a implementar firma del lado del cliente parseando el formato propietario del SAT. Ambas rutas introducen un problema de custodia de claves que contradice el objetivo declarado de no repudio: si el servidor custodia la clave privada del firmante, la firma pierde su propiedad vinculante porque no puede atribuirse exclusivamente a su titular. Con la CA interna, el ciclo de vida de las claves queda bajo control del prototipo y es auditable. Finalmente, montar la CA internamente expone al equipo al ciclo de vida completo de una PKI —generación de la clave raíz, emisión de CSRs, firma de certificados con extensiones X.509 v3, construcción de la cadena de confianza, validación y revocación— que una integración con la e.firma o un PSC comercial ocultaría detrás de una API externa. Este aprendizaje es parte del valor formativo declarado del Trabajo Terminal.

La implementación de la CA se realiza con OpenSSL scriptado en lugar de herramientas como step-ca o EJBCA. OpenSSL expone cada paso del proceso de forma transparente y auditable: la generación de la clave raíz (openssl genrsa), la creación del certificado autofirmado (openssl req -x509), la firma de solicitudes de certificado (openssl x509 -req -CA) y la generación de CRLs (openssl ca -gencrl) se encapsulan en aproximadamente cuatro invocaciones cuyo resultado es verificable con openssl x509 -text. step-ca automatiza y oculta esa mecánica detrás de una API REST, lo cual es deseable en producción pero contraproducente para un ejercicio académico. Además, OpenSSL no introduce dependencias adicionales al stack del prototipo. Respecto a las vigencias, el certificado raíz tiene una validez de cinco años para sobrevivir ampliamente al ciclo de vida del proyecto y al período de evaluación, mientras que los certificados de los socios tienen una validez de un año, alineada con la práctica estándar de la industria y suficiente para demostrar que el prototipo contempla la renovación de certificados.

\subsection{Sellado de Tiempo}\label{sec:sellado-tiempo}

El sellado de tiempo se diseña detrás del puerto \texttt{TimestampService}, con \texttt{LocalOpensslTsa} como adaptador actual de demostración y un adaptador de PSC externo como sustitución futura. En ambos casos el flujo opera conforme al protocolo RFC 3161 \cite{adams2001}: se calcula el hash SHA-256 del documento firmado, se solicita el sello al emisor configurado y se almacena la respuesta \texttt{TimeStampResp} en formato ASN.1 DER como artefacto independiente (.tsr). Esta abstracción permite sustituir la TSA local por Cincel, SeguriData, CECOBAN u otra TSA compatible sin modificar la lógica central. La equivalencia es de protocolo, no de confianza: únicamente el camino hacia un PSC autorizado y contratado puede producir la constancia externa prevista por la NOM-151-SCFI-2016.

Para facilitar la verificación independiente por un tercero (tribunal, contraparte, perito), el prototipo genera un paquete ZIP descargable que contiene el documento original, la firma RSA del socio, el sello de tiempo .tsr, los certificados del firmante y de la CA raíz, e instrucciones con comandos OpenSSL para la verificación completa. Este diseño garantiza que la verificación no dependa de la disponibilidad de la plataforma.

La selección inicial de Cincel como proveedor de sellado de tiempo obedeció a tres factores. Primero, Cincel es un PSC acreditado por la Secretaría de Economía para la emisión de Sellos Digitales de Tiempo (SDT) y constancias de conservación de mensajes de datos (CCMD), con acreditación publicada en el Diario Oficial de la Federación el 24 de noviembre de 2023 \cite{comparasoftware2026}. Segundo, ofrece un entorno de sandbox y recursos de integración. Tercero, la interfaz documentada permite enviar el hash SHA-256 y recuperar el artefacto ASN.1 DER conforme a RFC 3161 \cite{cincel2024,adams2001}, incluidos estados de procesamiento diferido que el prototipo maneja mediante reintentos configurables. La experiencia de integración corrigió la segunda premisa: aunque el registro se completó, el sandbox requiere pago con precios no transparentes y la API no ofreció respuesta y estabilidad suficientes para una campaña repetible. En consecuencia, Cincel no se utiliza como evidencia de esta entrega; permanece únicamente como adaptador externo opcional, mientras la TSA interna verifica el flujo completo sin convertir la disponibilidad comercial en requisito de desarrollo.

El sello de tiempo se almacena como la respuesta cruda \texttt{TimeStampResp} en formato ASN.1 DER (extensión .tsr) en lugar de empaquetarlo en una estructura CAdES-T o embebido en un PDF (PAdES-T). Esta decisión se justifica por tres razones. Primera, con artefactos independientes cada componente del flujo criptográfico —documento, firma, sello, certificado— puede inspeccionarse y verificarse por separado con herramientas estándar; el comando \texttt{openssl ts -verify} opera directamente sobre archivos .tsr sin necesidad de parsear contenedores CMS ni estructuras PDF. Segunda, CAdES-T requeriría construir estructuras CMS con atributos no firmados y PAdES-T requeriría manipular los mecanismos internos del formato PDF (\texttt{ByteRange}, diccionarios de firma y guardados incrementales), tareas de ingeniería de formatos que consumirían tiempo del cronograma sin aportar al objetivo criptográfico del Trabajo Terminal. Tercera, el contenedor no altera el token RFC 3161 subyacente ni la identidad de su emisor: el alcance jurídico depende de la autoridad y su cadena de confianza, no de almacenar el token como .tsr, CAdES-T o PAdES-T \cite{nom151_2016}.

El flujo de verificación completo del sello de tiempo procede en cuatro pasos: (1) se parsea la \texttt{TimeStampResp} para extraer el \texttt{TimeStampToken}; (2) se verifica que \texttt{messageImprint.\allowbreak{}hashedMessage} coincida con el hash SHA-256 recalculado del documento original; (3) se verifica la firma de la TSA sobre el token; y (4) se valida su certificado contra el ancla de confianza correspondiente. Para la TSA interna, esa ancla es la CA del prototipo; para una constancia externa, debe ser la cadena del PSC emisor hasta la ACR2-SE \cite{economia_acr2se}. Con OpenSSL, el proceso se ejecuta con \texttt{openssl ts -verify -data documento.pdf -in sello.tsr -CAfile ancla.pem}. Un despliegue productivo podría migrar a CAdES-T o PAdES-T sin cambiar el token subyacente ni el flujo de obtención.

\subsection{Cifrado en Reposo}\label{sec:cifrado-reposo}

Los documentos almacenados se cifran con AES-256-GCM conforme a NIST SP 800-38D \cite{dworkin2007}. Cada operación utiliza un nonce (número de uso único) aleatorio de 96 bits generado con un CSPRNG (Generador de Números Pseudoaleatorios Criptográficamente Seguro), cuya probabilidad de colisión es despreciable cuando cada documento emplea su propia clave (NIST SP 800-38D §8.3 establece un límite de 2³² invocaciones por clave). La gestión de claves sigue un esquema de envelope encryption: una Clave de Cifrado de Clave (KEK), almacenada en una variable de entorno del servidor, envuelve una Clave de Cifrado de Datos (DEK) generada aleatoriamente para cada documento conforme a NIST SP 800-38F \cite{dworkin2012}. Esta separación permite rotar la KEK re-cifrando únicamente las DEKs envueltas, sin necesidad de re-cifrar los documentos del repositorio. Los Datos Adicionales Autenticados (AAD) de GCM incluyen el identificador del documento y su versión, lo que previene ataques de intercambio (swap) y de reversión (rollback) entre documentos.

La elección de un nonce aleatorio sobre un nonce de contador responde a la prioridad de evitar el modo de fallo más catastrófico de GCM: la reutilización de un par (clave, nonce). Un nonce de contador requiere mantener un estado persistente que se incremente atómicamente en cada operación de cifrado; si ese estado se reinicia por un error de software, una restauración de respaldo o una condición de carrera, el contador repite un valor y el esquema GCM se rompe irrecuperablemente, exponiendo tanto el texto claro como la clave de autenticación. Un nonce aleatorio de 96 bits generado con un CSPRNG elimina esta dependencia de estado por diseño: cada invocación es independiente y no requiere coordinación con invocaciones anteriores. NIST SP 800-38D \cite{dworkin2007} §8.3 establece que con nonces aleatorios de 96 bits el número total de invocaciones por clave no debe exceder 2³² (aproximadamente 4,300 millones) para mantener la probabilidad de colisión por debajo de \textasciitilde{}2⁻³³. Con envelope encryption, donde cada documento emplea su propia DEK, el número de invocaciones por clave es efectivamente uno, lo que hace la probabilidad de colisión exactamente cero en la práctica.

La arquitectura de envelope encryption se seleccionó sobre dos alternativas: una clave única compartida para todo el repositorio y una derivación de claves mediante HKDF. Una sola clave compartida concentra el riesgo —su compromiso expone la totalidad de la información— y hace la rotación prohibitivamente costosa porque cada documento debe re-cifrarse. Un esquema basado en HKDF derivaría una clave por documento a partir de una clave maestra y un contexto único (por ejemplo, el identificador del documento), lo que evita almacenar claves envueltas pero presenta dos desventajas: si la clave maestra se compromete, todas las claves derivadas caen sin posibilidad de revocar una sola; y la unicidad absoluta del contexto de derivación reintroduce un problema de estado similar al del nonce de contador. Envelope encryption aísla el daño por diseño: el compromiso de una DEK solo expone un documento, y el compromiso de la KEK requiere además acceso al almacenamiento cifrado para descifrar las DEKs. Este es el patrón estándar de la industria, implementado por AWS KMS, Google Cloud KMS y Azure Key Vault \cite{dworkin2012}.

El flujo de cifrado procede de la siguiente manera al almacenar un documento: (1) se genera una DEK aleatoria de 32 bytes con un CSPRNG; (2) se genera un nonce aleatorio de 12 bytes; (3) se cifra el documento con AES-256-GCM utilizando la DEK, el nonce y los AAD (document\_id + version), obteniendo el ciphertext y un authentication tag de 16 bytes; (4) se genera un segundo nonce de 12 bytes para la envoltura; y (5) se cifra la DEK con la KEK usando AES-256-GCM. El paquete cifrado del documento se almacena en formato nonce (12 bytes) ‖ ciphertext ‖ tag (16 bytes), y la DEK envuelta sigue el mismo formato. El flujo de descifrado invierte el proceso: primero se descifra la DEK con la KEK verificando la integridad del tag, y después se descifra el documento con la DEK recuperada. Si cualquiera de los dos tags falla la verificación, la operación se rechaza íntegramente, lo que previene la entrega de documentos corrompidos o manipulados.

La rotación de la KEK es una operación ligera que no requiere re-cifrar ningún documento. Al generar una nueva KEK con un CSPRNG, se itera sobre las DEKs del repositorio: cada una se descifra con la KEK vigente y se re-cifra con la nueva. Dado que las DEKs son bloques de 32 bytes, esta operación es de tamaño fijo y se ejecuta en tiempo constante por documento, independientemente del tamaño del documento protegido. Los documentos en sí no se tocan: sus DEKs no cambian, solo cambia su envoltura. Esta separación es la ventaja principal de envelope encryption sobre un esquema de clave única: la rotación opera sobre claves (rápida, de tamaño fijo) y no sobre documentos (potencialmente de varios megabytes). Para el prototipo, basta con implementar y documentar el mecanismo de rotación; ejecutar una rotación periódica queda fuera del alcance de la demostración.

\subsection{Autenticación y Control de Acceso}\label{sec:autenticacion}

Las contraseñas de los usuarios se almacenan hasheadas con Argon2id (RFC 9106) \cite{biryukov2015,biryukov2021}, la variante híbrida recomendada por OWASP \cite{owasp_pwstorage2024} y NIST SP 800-63B Rev 4 \cite{nist2024} para nuevas implementaciones. Los parámetros adoptados son: memoria m = 256 MiB (262\,144 KiB), iteraciones t = 4, paralelismo p = 1, sal de 16 bytes y hash de salida de 32 bytes. En el hardware de referencia (AMD Ryzen 7 7730U, 16 hilos lógicos) cinco corridas promediaron 607.8 ms, dentro de la banda objetivo de 0.5--1.0 segundo por intento legítimo. Si el hardware de despliegue difiere, la medición debe repetirse antes de la puesta en producción.

Los clientes del despacho cuentan adicionalmente con autenticación de doble factor mediante TOTP (RFC 6238) \cite{mraihi2011,owasp_mfa2024} con HMAC-SHA1, códigos de 6 dígitos, período de 30 segundos y secreto de 20 bytes, alcanzando el nivel AAL2 de NIST SP 800-63B \cite{nist2024}. Se descartó WebAuthn/FIDO2 (AAL3) por la barrera de adopción que representa para clientes de un despacho legal que requiere solo un smartphone con una app autenticadora estándar. Los socios, por su parte, se autentican mediante certificados digitales X.509 emitidos por la CA interna (\cref{sec:firma-digital}), lo que provee un nivel de aseguramiento equivalente a AAL3 sin hardware dedicado adicional.

La selección de Argon2id sobre bcrypt se fundamenta en la propiedad de dureza de memoria (memory-hardness). bcrypt, diseñado en 1999, es resistente a fuerza bruta en CPU gracias a su costo iterativo, pero utiliza únicamente aproximadamente 4 KiB de memoria por invocación, lo que permite a un atacante con GPUs o ASICs ejecutar miles de instancias en paralelo a bajo costo por unidad de intento. Argon2id obliga a reservar 256 MiB por cada intento de hash en la configuración calibrada, lo que hace que escalar el ataque con hardware masivo sea prohibitivamente caro porque la memoria DRAM es el recurso que peor escala en paralelo. Argon2 fue el ganador de la Password Hashing Competition (PHC) en julio de 2015 \cite{biryukov2015}, una competencia abierta de dos años con 24 candidatos evaluados por un panel de criptógrafos. Esta validación pública es análoga a la que AES recibió del concurso NIST en 2001 y aporta confianza en que el diseño ha sido ampliamente escrutado.

Argon2 define tres variantes que difieren en su resistencia a distintos vectores de ataque. Argon2d maximiza la resistencia a ataques con GPU y ASIC, pero es vulnerable a ataques de canal lateral si el atacante puede observar los patrones de acceso a memoria del proceso. Argon2i es resistente a ataques de canal lateral al usar un patrón de acceso independiente de los datos, pero ofrece ligeramente menos resistencia a hardware especializado. Argon2id es la variante híbrida: ejecuta Argon2i en el primer pase para proteger contra canal lateral, y Argon2d en los pases subsiguientes para maximizar la resistencia a GPU. RFC 9106 \cite{biryukov2021} §4 recomienda explícitamente Argon2id para hashing de contraseñas por esta combinación de protecciones. Los parámetros adoptados (m = 256 MiB, t = 4, p = 1, sal de 16 bytes, hash de 32 bytes) fueron seleccionados por calibración del hardware de referencia; la sal de 128 bits excede el mínimo de 64 bits de RFC 9106 y es el valor recomendado por OWASP.

La selección de HMAC-SHA1 como función base del TOTP merece una justificación explícita dado que SHA-1 está desaconsejado para firmas digitales desde 2017 por la demostración práctica de colisiones. Sin embargo, la seguridad de HMAC no depende de la resistencia a colisiones de la función hash subyacente, sino de la pseudoaleatoriedad de su función de compresión interna. NIST SP 800-107 Rev 1 \cite{dang2012} confirma que HMAC-SHA-1 sigue siendo aceptable para autenticación de mensajes y generación de códigos. Adicionalmente, HMAC-SHA1 es el único algoritmo soportado universalmente por todas las aplicaciones autenticadoras del mercado; Google Authenticator históricamente solo soporta SHA-1, lo que haría que la adopción de HMAC-SHA256 o HMAC-SHA512 excluyera a una parte significativa de los dispositivos de los clientes del despacho.

El flujo de registro (enrollment) del segundo factor procede de la siguiente manera: (1) el cliente inicia sesión con usuario y contraseña; (2) el sistema genera un secreto de 20 bytes con un CSPRNG; (3) se codifica como URI otpauth (otpauth://totp/Despacho:cliente@email?secret=BASE32(secret)) y se presenta como código QR; (4) el cliente escanea con su aplicación autenticadora; (5) el cliente ingresa el código actual para confirmar la configuración; y (6) el sistema almacena el secreto cifrado con AES-256-GCM y genera ocho códigos de recuperación de un solo uso, hasheados con Argon2id antes de almacenarse. Durante la autenticación, el sistema calcula el TOTP esperado para tres ventanas temporales (t₀₋₁, t₀, t₀₊₁), lo que tolera un desfase de reloj de hasta 30 segundos entre servidor y dispositivo. Se aplica rate limiting de cinco intentos fallidos consecutivos con bloqueo temporal de 15 minutos para prevenir fuerza bruta en línea sobre el espacio de 10⁶ combinaciones (6 dígitos).

\subsection{Postura Legal del Prototipo}\label{sec:postura-legal}

La firma generada por el prototipo constituye una firma electrónica simple conforme al artículo 89 del Código de Comercio \cite{ccom2018}, dado que los certificados no provienen de un PSC autorizado (artículo 97, requisito 5). Su argumentación probatoria se apoya en: (1) la admisibilidad de la firma simple como prueba, ya que el artículo 89 no condiciona la admisibilidad a la intervención de un PSC; y (2) la cláusula de equivalencia funcional del artículo 100, que permite argumentar la fiabilidad del método con base en estándares reconocidos internacionalmente (RSA-3072 conforme a NIST SP 800-57 \cite{barker2020} y SHA-256 conforme a FIPS 180-4 \cite{nist2015fips1804}). Un sello emitido por un PSC autorizado añadiría fecha cierta e independencia de tercero conforme al marco de la NOM-151; el token de la TSA interna demuestra el mecanismo RFC 3161, pero no añade esa garantía externa. La distinción práctica respecto de una firma avanzada radica en la carga de la prueba: quien invoca la firma simple debe acreditar la fiabilidad del método, mientras que la firma avanzada goza de presunción legal. En un despliegue productivo, la sustitución de la CA interna por un PSC autorizado convertiría la firma en avanzada sin cambios arquitectónicos, dado que el módulo de firma acepta cualquier certificado X.509 RSA + SHA-256.

La arquitectura del prototipo está diseñada para que la transición de firma simple a firma avanzada sea una operación de configuración, no de rediseño. El módulo de firma define una interfaz abstracta que acepta cualquier certificado X.509 con clave RSA y algoritmo sha256WithRSAEncryption. En un despliegue productivo, el despacho podría sustituir la CA interna por certificados emitidos por un PSC autorizado como SeguriData \cite{seguridata} o CECOBAN \cite{cecoban}, convirtiendo la firma en avanzada conforme al artículo 97 del Código de Comercio \cite{ccom2018}; o bien integrar la e.firma del SAT si los socios ya disponen de ella, sin modificar la lógica de verificación. De forma análoga, el módulo de sellado de tiempo define una interfaz abstracta que permite sustituir Cincel por cualquier TSA que implemente RFC 3161 \cite{adams2001}. Esta portabilidad del emisor es un acierto arquitectónico deliberado: el prototipo demuestra la corrección de los mecanismos criptográficos con una infraestructura soberana del despacho, mientras que la producción puede adoptar la infraestructura regulada sin tocar el código de verificación.
